% Leçon 37 — Expression orale : internet — Chinois 1ère A — Cours signé OnBuch+
% Programme officiel MINESEC. Compiler avec : tectonic ZHA37-oral-internet.tex
\newcommand{\DOCMATIERE}{Chinois}
\newcommand{\DOCNIVEAU}{1ère A}
\newcommand{\DOCMODULE}{Module 5 : Mass médias et télécommunication}
\newcommand{\DOCLECON}{ZHA37}
\newcommand{\DOCTITRE}{Leçon 37 — Expression orale : internet}
% OnBuch+ — préambule « cours signés » ENRICHI (compilé avec tectonic / XeLaTeX)
% Système visuel « Pop » d'OnBuch+ : fond crème (jamais blanc), bordures encre
% épaisses, ombres « dures » décalées, titres Archivo Black en capitales.
% Version MATHÉMATIQUES : graphes (pgfplots/tikz), boîtes pédagogiques
% (définition, propriété, méthode, attention, le savais-tu, exercice résolu,
% exercice, TP, bilan), page de garde et signature OnBuch+ ancrée en bas.
%
% Macros attendues dans main.tex AVANT \input{preamble} :
%   \DOCMATIERE  \DOCNIVEAU  \DOCMODULE  \DOCLECON (numéro)  \DOCTITRE
\documentclass[11pt]{article}

\usepackage{fontspec}
\usepackage[french]{babel} % français = langue principale ; le chinois via \zh{..} / textebox (xeCJK)
\usepackage{xeCJK}
\usepackage{pifont} % \ding{51} \ding{55} utilisés par les modèles
\usepackage[a4paper,margin=2cm,top=3.4cm,bottom=2.8cm,headheight=1.4cm,footskip=1.3cm]{geometry}
\usepackage{amsmath,amssymb,amsfonts}
\usepackage{mathtools} % \xmapsto, \coloneqq... souvent utilisés par les modèles
\usepackage{cancel} % simplifications barrées (bilans de forces)
% \abs{z} (module, valeur absolue) et \norm{u} (norme), utilisés par les modèles
\providecommand{\abs}{}\let\abs\relax\DeclarePairedDelimiter{\abs}{\lvert}{\rvert}
\providecommand{\norm}{}\let\norm\relax\DeclarePairedDelimiter{\norm}{\lVert}{\rVert}
\usepackage[table]{xcolor} % [table] : fournit aussi \rowcolors (alternance de lignes)
\usepackage{pagecolor}
\usepackage{tikz}
\usetikzlibrary{arrows.meta,positioning,calc,shapes.geometric,shapes.misc,shapes.arrows,decorations.pathmorphing,decorations.pathreplacing,decorations.markings,patterns,shadows,mindmap,trees,fit,backgrounds,matrix,intersections,angles,quotes}
\usepackage{pgfplots}
\usepackage[version=4]{mhchem} % équations nucléaires \ce{^{235}_{92}U}
\usepackage[european,siunitx]{circuitikz} % schémas électriques
\pgfplotsset{compat=1.18}
\usepgfplotslibrary{fillbetween,groupplots} % aires entre courbes (calcul intégral)
\usepackage{enumitem}
\usepackage{fancyhdr}
% [most] charge toutes les bibliothèques tcolorbox (listings, minted, raster,
% documentation...) et gonfle inutilement la mémoire TeX sur un document long
% (~300 boîtes) : on ne charge que ce qui est réellement utilisé.
\usepackage[skins,breakable]{tcolorbox}
\usepackage{graphicx}
\usepackage{colortbl}
\usepackage{booktabs}
\usepackage{tabularx}
\usepackage{diagbox} % case d'en-tête diagonale des tableaux à double entrée
% Colonnes extensibles centrée / gauche / droite (C, L, R), souvent utilisées par les modèles
\newcolumntype{C}{>{\centering\arraybackslash}X}
\newcolumntype{L}{>{\raggedright\arraybackslash}X}
\newcolumntype{R}{>{\raggedleft\arraybackslash}X}
\usepackage{array}
\usepackage{multicol}
\usepackage{multirow}
\usepackage{siunitx}
% parse-numbers=false : accepte aussi \SI{2,5 \times 10^{-3}}{...}
\sisetup{locale=FR,output-decimal-marker={,},group-minimum-digits=4,parse-numbers=false}
\DeclareSIUnit\ohm{\ensuremath{\symup{\Omega}}} % U+2126 absent de la police
\DeclareSIUnit\division{div} % réglages d'oscilloscope (ms/div, V/div)
\DeclareSIUnit\tr{tr} % tours (vitesse de rotation, ex. \SI{1500}{\tr\per\minute})
\DeclareSIUnit\molar{M} % concentration molaire (ex. \SI{3}{\molar}, \SI{1}{\micro\molar})
\DeclareSIUnit\an{an} % année (le modèle confond parfois avec \year, un compteur TeX) — ex. \SI{5}{\kilo\meter\per\an}
\DeclareSIUnit\FCFA{FCFA} % franc CFA (ex. \SI{500}{\FCFA\per\kilo\gram})
\DeclareSIUnit\degreeBrix{\textdegree Brix} % teneur en sucre d'un sirop/jus (réfractométrie)
\DeclareSIUnit\month{mois} % le modèle utilise parfois \month, un compteur TeX, comme unité de durée
\usepackage{qrcode}
% \degree : commande fréquemment utilisée par les modèles pour un angle ou une
% température en dehors de \SI{}{} (non fournie par les paquets chargés).
\providecommand{\degree}{^{\circ}}
% \qed (fin de démonstration), utilisé par les modèles sans amsthm
\providecommand{\qed}{\hfill$\square$}
% \barre : double barre des valeurs interdites dans un tableau de variations (tkz-tab)
\providecommand{\barre}{\ensuremath{\|}}
% environnement proof (amsthm non chargé) utilisé par les modèles
\ifdefined\proof\else\newenvironment{proof}[1][Démonstration]{\par\noindent\textit{#1.}\ }{\qed\par}\fi
\usepackage{lastpage}
\usepackage{hyperref}
\hypersetup{hidelinks}

% ============================================================
% Palette « Pop » OnBuch+ — mêmes hex que la classe P (plus_ui.dart)
% ============================================================
\definecolor{poporange}{HTML}{F59321}
\definecolor{popdark}{HTML}{E07A0C}
\definecolor{popcream}{HTML}{FFF6E8}
\definecolor{popink}{HTML}{1C1714}
\definecolor{poppurple}{HTML}{7A5AE0}
\definecolor{popgreen}{HTML}{2E9E6A}
\definecolor{poppink}{HTML}{E0457B}
\definecolor{popblue}{HTML}{2F7DE1}
\definecolor{popgold}{HTML}{C98A12}
\definecolor{popmuted}{HTML}{8A7F76}
\definecolor{popline}{HTML}{EADFD2}
\definecolor{popgray}{HTML}{8A7F76}
\colorlet{popgrey}{popgray}
% Alias tolérants (le modèle nomme parfois les couleurs différemment)
\colorlet{popwhite}{white}
\colorlet{popblack}{popink}
\colorlet{popred}{poppink}
\colorlet{popyellow}{popgold}
% Teintes claires (fonds de tableaux/encadrés) — le modèle nomme parfois
% les couleurs avec un suffixe L (ex. popblueL) sans les avoir définies.
\colorlet{poporangeL}{poporange!20}
\colorlet{popdarkL}{popdark!20}
\colorlet{poppurpleL}{poppurple!20}
\colorlet{popgreenL}{popgreen!20}
\colorlet{poppinkL}{poppink!20}
\colorlet{popblueL}{popblue!20}
\colorlet{popgoldL}{popgold!20}
\colorlet{popredL}{popred!20}
\colorlet{popgrayL}{popgray!20}
% Teintes claires pour fonds de tableaux / zones
\colorlet{popblueL}{popblue!10!white}
\colorlet{poporangeL}{poporange!14!white}
\colorlet{popgreenL}{popgreen!12!white}
\colorlet{poppurpleL}{poppurple!12!white}
\colorlet{poppinkL}{poppink!12!white}
\colorlet{popgoldL}{popgold!14!white}

\pagecolor{popcream}

% ============================================================
% Typographie
% ============================================================
\defaultfontfeatures{Ligatures=TeX}
\setmainfont{PlusJakartaSans-Medium.ttf}[Path=../../../fonts/,BoldFont=PlusJakartaSans-Bold.ttf]
% Symboles, API et emojis absents de la police principale : repli sur DejaVu Sans (caractères actifs)
\newfontface\symfont{DejaVuSans.ttf}[Path=../../../fonts/]
\makeatletter
\newcommand\symactive[1]{\catcode#1=13 \begingroup\lccode`\~=#1 \lowercase{\endgroup\def~}{{\symfont\char#1}}}
\newcommand\symblank[1]{\catcode#1=13 \begingroup\lccode`\~=#1 \lowercase{\endgroup\def~}{}}
\newcommand\symhyph[1]{\catcode#1=13 \begingroup\lccode`\~=#1 \lowercase{\endgroup\def~}{\mbox{-}}}
\newcount\symc
\newcommand\symrange[2]{\symc=#1 \@whilenum\symc<#2 \do{\expandafter\symactive\expandafter{\the\symc}\advance\symc by 1 }}
\newcommand\symblankrange[2]{\symc=#1 \@whilenum\symc<#2 \do{\expandafter\symblank\expandafter{\the\symc}\advance\symc by 1 }}
\makeatother
\symrange{"0250}{"02FF}\symactive{"2713}\symactive{"2714}\symactive{"2717}\symactive{"2718}\symactive{"2610}\symactive{"2611}\symactive{"2612}
\symrange{"2600}{"27BF}
\catcode"2705=13 \begingroup\lccode`\~="2705 \lowercase{\endgroup\def~}{{\symfont\char"2713}}\symblank{"26BD}\symblank{"26A1}
\symrange{"2460}{"257F}\symactive{"223C}\symactive{"2248}\symactive{"2260}
\symactive{"01F8}\symactive{"01F9}\symactive{"1E3E}\symactive{"1E3F}
\symhyph{"2011}\symhyph{"2E1A}\symblankrange{"1F300}{"1FAFF}\symblank{"FE0F}
% Caractères chinois : WenQuanYi Zen Hei (sous-ensemble GB2312 + pinyin), gras/italique simulés
\setCJKmainfont{WenQuanYiZenHei-subset.ttf}[Path=../../../fonts/,AutoFakeBold=3,AutoFakeSlant=0.2]
\xeCJKsetup{CJKspace=false}
% Pinyin : voyelles à caron (ǎ ǐ ǒ ǔ ǖ ǘ ǚ ǜ) absentes de la police principale -> caractères actifs
\newfontface\pyfont{WenQuanYiZenHei-subset.ttf}[Path=../../../fonts/]
\newcommand\pyactive[1]{\catcode#1=13 \begingroup\lccode`\~=#1 \lowercase{\endgroup\def~}{{\pyfont\char#1}}}
\pyactive{"01CD}\pyactive{"01CE}\pyactive{"01CF}\pyactive{"01D0}\pyactive{"01D1}\pyactive{"01D2}\pyactive{"01D3}\pyactive{"01D4}\pyactive{"01D5}\pyactive{"01D6}\pyactive{"01D7}\pyactive{"01D8}\pyactive{"01D9}\pyactive{"01DA}\pyactive{"01DB}\pyactive{"01DC}
% Maths en sans-serif (Fira Math) avec chiffres et lettres droites de Plus
% Jakarta, pour que les formules se fondent dans le texte.
\usepackage[math-style=ISO,bold-style=ISO]{unicode-math}
\setmathfont{FiraMath-Regular.otf}[Path=../../../fonts/]
\setmathfont{PlusJakartaSans-Medium.ttf}[Path=../../../fonts/,range={up/num,up/{Latin,latin},\mathup}]
\usepackage{icomma} % virgule décimale sans espace en mode math
% Caractères mathématiques Unicode absents des polices de texte (Plus Jakarta,
% Archivo Black) : rendus par leur équivalent mathématique, en texte comme en
% formule — sinon ils s'affichent en carré vide (ex. « Arithmétique dans ℤ »).
\usepackage{newunicodechar}
\newunicodechar{ℝ}{\ensuremath{\mathbb{R}}}
\newunicodechar{ℤ}{\ensuremath{\mathbb{Z}}}
\newunicodechar{ℂ}{\ensuremath{\mathbb{C}}}
\newunicodechar{ℕ}{\ensuremath{\mathbb{N}}}
\newunicodechar{ℚ}{\ensuremath{\mathbb{Q}}}
\newunicodechar{𝔻}{\ensuremath{\mathbb{D}}}
\newunicodechar{α}{\ensuremath{\alpha}}
\newunicodechar{β}{\ensuremath{\beta}}
\newunicodechar{γ}{\ensuremath{\gamma}}
\newunicodechar{δ}{\ensuremath{\delta}}
\newunicodechar{ε}{\ensuremath{\varepsilon}}
\newunicodechar{θ}{\ensuremath{\theta}}
\newunicodechar{λ}{\ensuremath{\lambda}}
\newunicodechar{μ}{\ensuremath{\mu}}
\newunicodechar{σ}{\ensuremath{\sigma}}
\newunicodechar{τ}{\ensuremath{\tau}}
\newunicodechar{φ}{\ensuremath{\varphi}}
\newunicodechar{ψ}{\ensuremath{\psi}}
\newunicodechar{ω}{\ensuremath{\omega}}
\newunicodechar{Σ}{\ensuremath{\Sigma}}
% majuscules produites par \MakeUppercase dans les titres (α-aminés -> Α)
\newunicodechar{Α}{\ensuremath{\alpha}}
\newunicodechar{Β}{\ensuremath{\beta}}
\newunicodechar{Γ}{\ensuremath{\Gamma}}
\newunicodechar{Δ}{\ensuremath{\Delta}}
\newunicodechar{Ε}{\ensuremath{\varepsilon}}
\newunicodechar{Θ}{\ensuremath{\Theta}}
\newunicodechar{Λ}{\ensuremath{\Lambda}}
\newunicodechar{Μ}{\ensuremath{\mu}}
\newunicodechar{Τ}{\ensuremath{\tau}}
\newunicodechar{Φ}{\ensuremath{\Phi}}
\newunicodechar{Ψ}{\ensuremath{\Psi}}
\newunicodechar{Ω}{\ensuremath{\Omega}}
\newunicodechar{↦}{\ensuremath{\mapsto}}
\newunicodechar{∈}{\ensuremath{\in}}
\newunicodechar{∉}{\ensuremath{\notin}}
\newunicodechar{∧}{\ensuremath{\wedge}}
\newunicodechar{⇔}{\ensuremath{\Leftrightarrow}}
\newunicodechar{⟺}{\ensuremath{\Longleftrightarrow}}
\newunicodechar{⇒}{\ensuremath{\Rightarrow}}
\newunicodechar{⟹}{\ensuremath{\Longrightarrow}}
\newunicodechar{∘}{\ensuremath{\circ}}
\newunicodechar{∪}{\ensuremath{\cup}}
\newunicodechar{∩}{\ensuremath{\cap}}
\newunicodechar{≡}{\ensuremath{\equiv}}
\newunicodechar{⊂}{\ensuremath{\subset}}
\newunicodechar{⊥}{\ensuremath{\perp}}
\newunicodechar{∃}{\ensuremath{\exists}}
\newunicodechar{∀}{\ensuremath{\forall}}
\newunicodechar{∛}{\ensuremath{\sqrt[3]{\,}}}
\newunicodechar{∜}{\ensuremath{\sqrt[4]{\,}}}
\newunicodechar{⃗}{\ensuremath{{}^{\to}}}
\newunicodechar{ⁿ}{\ifmmode^{n}\else\textsuperscript{\ensuremath{n}}\fi}
\newunicodechar{ˣ}{\ifmmode^{x}\else\textsuperscript{\ensuremath{x}}\fi}
\newunicodechar{ᵇ}{\ifmmode^{b}\else\textsuperscript{\ensuremath{b}}\fi}
\newunicodechar{ʳ}{\ifmmode^{r}\else\textsuperscript{\ensuremath{r}}\fi}
\newunicodechar{ᵘ}{\ifmmode^{u}\else\textsuperscript{\ensuremath{u}}\fi}
\newunicodechar{ʸ}{\ifmmode^{y}\else\textsuperscript{\ensuremath{y}}\fi}
\newunicodechar{ᵃ}{\ifmmode^{a}\else\textsuperscript{\ensuremath{a}}\fi}
\newunicodechar{ᵉ}{\ifmmode^{e}\else\textsuperscript{\ensuremath{e}}\fi}
\newunicodechar{ᵏ}{\ifmmode^{k}\else\textsuperscript{\ensuremath{k}}\fi}
\newunicodechar{ᵖ}{\ifmmode^{p}\else\textsuperscript{\ensuremath{p}}\fi}
\newunicodechar{ᴺ}{\ifmmode^{N}\else\textsuperscript{\ensuremath{N}}\fi}
\newunicodechar{ᶜ}{\ifmmode^{c}\else\textsuperscript{\ensuremath{c}}\fi}
\newunicodechar{ʰ}{\ifmmode^{h}\else\textsuperscript{\ensuremath{h}}\fi}
\newunicodechar{ᵠ}{\ifmmode^{q}\else\textsuperscript{\ensuremath{q}}\fi}
\newunicodechar{ₙ}{\ifmmode_{n}\else\textsubscript{\ensuremath{n}}\fi}
\newunicodechar{ₐ}{\ifmmode_{a}\else\textsubscript{\ensuremath{a}}\fi}
\newunicodechar{ᵢ}{\ifmmode_{i}\else\textsubscript{\ensuremath{i}}\fi}
\newunicodechar{ᵣ}{\ifmmode_{r}\else\textsubscript{\ensuremath{r}}\fi}
\newunicodechar{ₖ}{\ifmmode_{k}\else\textsubscript{\ensuremath{k}}\fi}
\newunicodechar{ᵦ}{\ifmmode_{\beta}\else\textsubscript{\ensuremath{\beta}}\fi}
\newunicodechar{ₓ}{\ifmmode_{x}\else\textsubscript{\ensuremath{x}}\fi}
\newfontfamily\popdisplay{ArchivoBlack-Regular.ttf}[Path=../../../fonts/]
\newfontfamily\poplogo{SpaceGrotesk-Bold.ttf}[Path=../../../fonts/]
\newfontfamily\popsemi{PlusJakartaSans-SemiBold.ttf}[Path=../../../fonts/]
\newfontfamily\popxbold{PlusJakartaSans-ExtraBold.ttf}[Path=../../../fonts/]
\newfontfamily\popxboldls{PlusJakartaSans-ExtraBold.ttf}[Path=../../../fonts/,LetterSpace=3.0]

\setlist[enumerate]{leftmargin=1.7em,itemsep=5pt,topsep=4pt}
\setlist[itemize]{leftmargin=1.7em,itemsep=4pt,topsep=4pt,label={\color{poporange}\textbullet}}
\setlength{\parindent}{0pt}
\setlength{\parskip}{5pt}
\renewcommand{\arraystretch}{1.25}

% pgfplots : style Pop par défaut
\pgfplotsset{
  every axis/.append style={axis line style={popink,line width=1pt},tick style={popink},
    grid style={popline,line width=0.5pt},label style={font=\small},tick label style={font=\footnotesize}},
  popcurve/.style={poporange,line width=1.8pt,no markers,smooth},
}
% Même style disponible dans un \draw TikZ simple (hors pgfplots)
\tikzset{popcurve/.style={poporange,line width=1.8pt}}
\tikzset{popgrid/.style={popline,very thin}}
\colorlet{popcurve}{poporange} % le nom du style est parfois utilisé comme couleur

% ============================================================
% Pill / squiggle
% ============================================================
\newcommand{\poppill}[3]{%
  \tikz[baseline=-0.55ex]\node[draw=popink,line width=1.2pt,rounded corners=2.6pt,%
    fill=#2,inner xsep=6.5pt,inner ysep=3pt]{%
    \popxboldls\fontsize{7}{7}\selectfont\color{#3}\MakeUppercase{#1}};%
}
\newcommand{\popsquiggle}[1]{%
  \tikz[baseline=-0.4ex]{
    \draw[#1,line width=2.6pt,line cap=round]
      plot [smooth] coordinates {(0,0.09) (0.34,0.01) (0.68,0.21) (1.02,0.01) (1.36,0.10)};
  }%
}
% Mot-clé surligné (marqueur orange sous le texte)
\newcommand{\cle}[1]{{\popxbold\color{popdark}#1}}

% ============================================================
% En-tête / pied
% ============================================================
\pagestyle{fancy}
\fancyhf{}
\renewcommand{\headrulewidth}{0pt}
\renewcommand{\footrulewidth}{0pt}
\lhead{\raisebox{-0.15ex}{\poplogo\fontsize{13}{13}\selectfont\color{popink}OnBuch\color{poporange}+}}
\rhead{\poppill{\DOCMATIERE\ \textbullet\ \DOCNIVEAU\ \textbullet\ Leçon \DOCLECON}{white}{popink}}
\lfoot{\popxbold\fontsize{7.6}{7.6}\selectfont\color{popmuted}\MakeUppercase{Cours signé OnBuch+}\ \textbullet\ {\popsemi\DOCTITRE}}
\rfoot{\tikz[baseline=-0.6ex]\node[rounded corners=8.5pt,fill=popink,minimum height=17pt,minimum width=17pt,inner xsep=5pt,inner sep=0]{\popxbold\fontsize{8}{8}\selectfont\color{popcream}\thepage\,/\,\pageref*{LastPage}};}

% ============================================================
% Titres
% ============================================================
\newcommand{\chaptitre}[2]{%
  \begin{tcolorbox}[enhanced,colback=poporange,colframe=popink,boxrule=2.6pt,arc=11pt,
      drop shadow=popink,left=15pt,right=15pt,top=12pt,bottom=11pt,before skip=3pt,after skip=18pt]
    \poppill{#2}{popink}{popcream}\par\vspace{9pt}
    {\raggedright\hyphenpenalty=10000\exhyphenpenalty=10000\popdisplay\fontsize{22}{24}\selectfont\color{popcream}\MakeUppercase{#1}\par}
  \end{tcolorbox}
}
% Section numérotée : pastille orange avec le numéro + titre Archivo
\newcounter{popsec}
\newcommand{\coursec}[1]{%
  \stepcounter{popsec}\setcounter{popsub}{0}%
  \par\vspace{16pt}\noindent
  \begin{minipage}{\linewidth}
  \tikz[baseline=-0.7ex]\node[draw=popink,line width=1.6pt,fill=poporange,rounded corners=4pt,minimum size=22pt,inner sep=0]{\popdisplay\fontsize{12}{12}\selectfont\color{popcream}\thepopsec};%
  \hspace{8pt}{\popdisplay\fontsize{15}{17}\selectfont\color{popink}\MakeUppercase{#1}}\par
  \vspace{4pt}\noindent{\color{poporange}\rule{36pt}{3.6pt}}
  \end{minipage}\par\nopagebreak\vspace{8pt}
}
\newcounter{popsub}
\newcommand{\courssub}[1]{%
  \stepcounter{popsub}%
  \par\vspace{9pt}\noindent{\popxbold\fontsize{11.5}{13}\selectfont\color{popink}{\color{poporange}\thepopsec.\thepopsub}\hspace{6pt}#1}\par\nopagebreak\vspace{4pt}
}

% ============================================================
% Boîtes pédagogiques — même recette : fond blanc, bordure épaisse colorée,
% ombre dure encre, badge plein. Titre optionnel [#1].
% ============================================================
\tcbset{popbase/.style={breakable,enhanced,colback=white,boxrule=2.3pt,arc=9pt,
  shadow={3pt}{-3pt}{0pt}{popink},left=14pt,right=14pt,top=11pt,bottom=11pt,before skip=11pt,after skip=15pt,
  fontupper=\popsemi\fontsize{10.6}{14.5}\selectfont\color{popink},
  fontlower=\popsemi\fontsize{10.6}{14.5}\selectfont\color{popink}}}
% Coche dessinée (le glyphe ✓ n'existe pas dans les polices du document)
\newcommand{\popcheck}[1]{\tikz[baseline=-0.5ex]\draw[#1,line width=1.6pt,line cap=round,line join=round](-0.13,0.0)--(-0.04,-0.09)--(0.13,0.1);}
\providecommand{\coche}{\popcheck{popgreen}}
\providecommand{\resultat}[1]{\par\smallskip\noindent\fcolorbox{popgreen}{popgreenL}{\parbox{\dimexpr\linewidth-2\fboxsep-2\fboxrule\relax}{\textbf{Résultat.} #1}}\par\smallskip}
\newcommand{\popboxhead}[3]{\poppill{#1}{#2}{white}\ifx\relax#3\relax\else\hspace{6pt}{\popxbold\fontsize{10.6}{12}\selectfont\color{popink}#3}\fi\par\vspace{7pt}}

\newtcolorbox{aretenir}[1][]{popbase,colframe=popgreen,before upper={\popboxhead{À retenir}{popgreen}{#1}}}
% Texte en chinois (document de lecture, dialogue, extrait) : cadre
% d'étude. Un titre optionnel (source, auteur) s'affiche dans l'en-tête.
\newtcolorbox{textebox}[1][]{popbase,colframe=popblue,before upper={\popboxhead{文本 · Texte}{popblue}{#1}},after upper={}}
% Marques ✓ / ✗ (forme correcte / fautive) souvent utilisées par les modèles
\providecommand{\cross}{\textcolor{poppink}{\ensuremath{\times}}}
\providecommand{\xmark}{\cross}\providecommand{\crossmark}{\cross}\providecommand{\cmark}{\ensuremath{\checkmark}}
% Ligne de réponse / justification à compléter (inventée par les modèles)
\providecommand{\justification}[1]{\par\noindent\textit{Justification~:} \dotfill\par}
% \textipa (package tipa non chargé) : la police ne couvre pas l'API -> simple passage
\providecommand{\textipa}[1]{#1}
% Soulignement (ulem non chargé) : \uline, \ul
\providecommand{\uline}[1]{\underline{#1}}\providecommand{\ul}[1]{\underline{#1}}
% \glossaire{\item ...} inventée par les modèles : liste de glossaire
\providecommand{\glossaire}[1]{\begin{itemize}#1\end{itemize}}
% \alert (beamer) inventée par les modèles : texte mis en évidence
\providecommand{\alert}[1]{\textbf{\textcolor{poppink}{#1}}}
% variantes ASCII d'ordinaux écrites en macro
\providecommand{\iere}{\textsuperscript{re}}\providecommand{\ieme}{\textsuperscript{e}}\providecommand{\ere}{\textsuperscript{re}}\providecommand{\eme}{\textsuperscript{e}}\providecommand{\er}{\textsuperscript{er}}\providecommand{\ie}{\textsuperscript{e}}
% Ordinaux français écrits en macro par les modèles : 1\ière, 3\ième
\newcommand{\ière}{\textsuperscript{re}}\newcommand{\ième}{\textsuperscript{e}}
% \exemple (inventée par les modèles) : étiquette « Exemple : »
\providecommand{\exemple}{\textbf{Exemple~:}\ }
% \square utilisable aussi hors mode math (cases à cocher)
\AtBeginDocument{\let\squareMath\square\renewcommand{\square}{\ensuremath{\squareMath}}}
% Mot ou phrase en chinois à l'intérieur d'un paragraphe français
\newcommand{\zh}[1]{\ifmmode\text{#1}\else{#1}\fi}
\let\eng\zh \let\esp\zh \let\all\zh % alias tolérés
% flèches d'intonation inventées par les modèles
\providecommand{\textsearrow}{}\renewcommand{\textsearrow}{\ensuremath{\searrow}}\providecommand{\textnearrow}{}\renewcommand{\textnearrow}{\ensuremath{\nearrow}}
% \fr{..} : mot français (inventé par les modèles)
\providecommand{\fr}[1]{\textit{#1}}
% zhbox : encadré de texte chinois inventé par les modèles
\newenvironment{zhbox}{\begin{quote}}{\end{quote}}
% \courssubsub : titre de niveau 3 inventé par les modèles
\providecommand{\courssubsub}[1]{\par\medskip\noindent\textbf{#1}\par\smallskip}
% \zhbf : texte chinois en gras inventé par les modèles
\providecommand{\zhbf}[1]{\textbf{#1}}
% \vs : « versus » inventé par les modèles
\providecommand{\vs}{\textit{vs}\ }
% \blank : case à compléter inventée par les modèles
\providecommand{\blank}[1][]{\underline{\hspace{1.6em}}}
\newtcolorbox{exemplebox}[1][]{popbase,colframe=poporange,before upper={\popboxhead{Exemple}{poporange}{#1}}}
\newtcolorbox{definition}[1][]{popbase,colframe=popblue,before upper={\popboxhead{Définition}{popblue}{#1}}}
\newtcolorbox{propriete}[1][]{popbase,colframe=poppurple,before upper={\popboxhead{Propriété}{poppurple}{#1}}}
\newtcolorbox{methode}[1][]{popbase,colframe=popgold,before upper={\popboxhead{Méthode}{popgold}{#1}}}
\newtcolorbox{attention}[1][]{popbase,colframe=poppink,before upper={\popboxhead{Attention}{poppink}{#1}}}
\newtcolorbox{experience}[1][]{popbase,colframe=popblue,colback=popblueL,before upper={\popboxhead{Expérience}{popblue}{#1}}}
% « Le savais-tu ? » : carte violette pleine (contexte, culture, Cameroun)
\newtcolorbox{savaistu}[1][]{popbase,colback=poppurple,colframe=popink,
  fontupper=\popsemi\fontsize{10.4}{14.2}\selectfont\color{white},
  before upper={\poppill{Le savais-tu\,?}{white}{poppurple}\ifx\relax#1\relax\else\hspace{6pt}{\popxbold\color{white}#1}\fi\par\vspace{7pt}}}
% Exercice résolu : énoncé en haut, \tcblower puis la solution sur fond crème
\newcounter{popexo}\newcounter{popexr}
\newtcolorbox{exoresolu}[1][]{popbase,colframe=poporange,colbacklower=popcream,
  segmentation style={popink,line width=1.2pt,dashed},
  before upper={\stepcounter{popexr}\popboxhead{Exercice résolu \thepopexr}{poporange}{#1}},
  before lower={\poppill{Solution}{popink}{popcream}\par\vspace{6pt}}}
% Exercice (non résolu en place) — la correction est dans la partie Corrigés
\newtcolorbox{exercice}[1][]{popbase,colframe=popink,
  before upper={\stepcounter{popexo}\popboxhead{Exercice \thepopexo}{popink}{#1}}}
% Corrigé
\newtcolorbox{corrige}[1][]{popbase,colframe=popgreen,colback=popgreenL,
  before upper={\popboxhead{Corrigé}{popgreen}{#1}}}
% Objectifs / prérequis / bilan
\newtcolorbox{objectifs}{popbase,colframe=popink,colback=poporangeL,
  before upper={\popboxhead{Objectifs de la leçon}{popink}{}}}
\newtcolorbox{prerequis}{popbase,colframe=popink,colback=popblueL,
  before upper={\popboxhead{Prérequis}{popblue}{}}}
\newtcolorbox{bilan}{popbase,colframe=popink,colback=white,boxrule=2.8pt,
  before upper={\popboxhead{Fiche bilan}{poporange}{L'essentiel à connaître pour l'examen}}}
% Figure encadrée avec légende
\newtcolorbox{popfigure}[1][]{enhanced,breakable=false,colback=white,colframe=popline,boxrule=1.4pt,arc=8pt,
  left=10pt,right=10pt,top=10pt,bottom=8pt,before skip=10pt,after skip=12pt,halign=center,
  lower separated=false,halign lower=center,
  fontlower=\popsemi\fontsize{9}{11.5}\selectfont\color{popmuted},
  #1}
\newcommand{\legende}[1]{\par\vspace{6pt}{\popsemi\fontsize{9}{11.5}\selectfont\color{popmuted}#1}}

% ============================================================
% Signature OnBuch+
% ============================================================
\newcommand{\poplogoinline}[1]{{\poplogo\fontsize{#1}{#1}\selectfont\color{popink}OnBuch\color{poporange}+}}
\newcommand{\signatureonbuch}{%
  \begin{tcolorbox}[enhanced,colback=popink,colframe=popink,boxrule=2.2pt,arc=10pt,
      drop shadow=poporange,left=14pt,right=14pt,top=10pt,bottom=10pt,before skip=0pt,after skip=0pt]
    \tikz[baseline=-0.7ex]\node[circle,fill=popgreen,draw=popcream,line width=1.3pt,minimum size=22pt,inner sep=0]{\popcheck{white}};%
    \hspace{9pt}\begin{minipage}[c]{0.62\linewidth}
      {\popxbold\fontsize{12}{13}\selectfont\color{popcream}Signé\ \textbullet\ {\poplogo OnBuch\color{poporange}+}}\par\vspace{2pt}
      {\raggedright\popsemi\fontsize{8}{10.5}\selectfont\color{popline}\MakeUppercase{Équipe pédagogique OnBuch+}\\\MakeUppercase{Conforme au programme officiel MINESEC}\par}
    \end{minipage}\hfill
    {\poplogo\fontsize{22}{22}\selectfont\color{popcream}OnBuch\color{poporange}+}
  \end{tcolorbox}%
}

% ============================================================
% Page de garde
% #1 titre · #2 matière · #3 niveau · #4 module · #5 n° leçon · #6 durée ·
% #7 description / objectifs (texte ou itemize)
% ============================================================
\newcommand{\pagegarde}[7]{%
  \thispagestyle{empty}
  \newgeometry{a4paper,margin=1.7cm,top=1.6cm,bottom=1.4cm}%
  \begin{tikzpicture}[remember picture,overlay]
    \fill[poporange!18] ([xshift=-3.2cm,yshift=-3.6cm]current page.north east) circle (4.6cm);
    \fill[poppurple!14] ([xshift=2.4cm,yshift=7.6cm]current page.south west) circle (3.8cm);
  \end{tikzpicture}%
  {\poplogo\fontsize{30}{30}\selectfont\color{popink}OnBuch\color{poporange}+}\hfill
  \raisebox{0.6ex}{\poppill{Cours signé \textbullet\ Premium}{popink}{popcream}}\par
  \vspace{1pt}\popsquiggle{poppurple}\par
  \vspace{8pt}
  \begin{tcolorbox}[enhanced,colback=poporange,colframe=popink,boxrule=2.8pt,arc=13pt,
      drop shadow=popink,left=18pt,right=18pt,top=12pt,bottom=13pt,after skip=10pt]
    \poppill{#2 \textbullet\ #3}{popink}{popcream}\hspace{5pt}\poppill{#4}{white}{popink}\par\vspace{8pt}
    {\popdisplay\fontsize{14}{14}\selectfont\color{popink}LEÇON #5}\par\vspace{4pt}
    {\raggedright\hyphenpenalty=10000\exhyphenpenalty=10000\popdisplay\fontsize{25}{27}\selectfont\color{popcream}\MakeUppercase{#1}\par}
    \vspace{8pt}
    {\popxbold\fontsize{9.5}{9.5}\selectfont\color{popink}\MakeUppercase{Durée conseillée : #6}}
  \end{tcolorbox}
  \begin{tcolorbox}[enhanced,colback=white,colframe=popink,boxrule=2.2pt,arc=10pt,
      drop shadow=popink,left=16pt,right=16pt,top=10pt,bottom=10pt,after skip=8pt]
    \poppill{Dans cette leçon}{poppurple}{white}\par\vspace{5pt}
    \popsemi\fontsize{9.3}{12.6}\selectfont\color{popink}#7
  \end{tcolorbox}
  \noindent\begin{minipage}[t]{0.487\linewidth}
    \begin{tcolorbox}[enhanced,colback=popgreen,colframe=popink,boxrule=2.2pt,arc=10pt,
        drop shadow=popink,left=11pt,right=11pt,top=10pt,bottom=10pt,height=3.1cm,valign=center]
      \tcbox[colback=white,colframe=popink,boxrule=1.3pt,arc=5pt,boxsep=0pt,nobeforeafter,
        left=3pt,right=3pt,top=3pt,bottom=3pt,valign=center]{\qrcode[height=1.9cm]{https://play.google.com/store/apps/details?id=com.luvvix.onbuch&hl=fr}}%
      \hspace{8pt}\begin{minipage}[c]{0.5\linewidth}
        {\popdisplay\fontsize{11}{12.5}\selectfont\color{white}\MakeUppercase{Télécharge OnBuch}}\par\vspace{4pt}
        {\raggedright\popsemi\fontsize{8.4}{11}\selectfont\color{white}Gratuit \textbullet\ Google Play\\Tuteur IA, annales, concours, résultats\par}
      \end{minipage}
    \end{tcolorbox}
  \end{minipage}\hfill
  \begin{minipage}[t]{0.487\linewidth}
    \begin{tcolorbox}[enhanced,colback=poppurple,colframe=popink,boxrule=2.2pt,arc=10pt,
        drop shadow=popink,left=11pt,right=11pt,top=10pt,bottom=10pt,height=3.1cm,valign=center]
      \tikz[baseline=-0.7ex]\node[circle,fill=white,draw=popink,line width=1.3pt,minimum size=1.6cm,inner sep=0]{\popdisplay\fontsize{24}{24}\selectfont\color{poppurple}+};%
      \hspace{8pt}\begin{minipage}[c]{0.6\linewidth}
        {\popdisplay\fontsize{11}{12.5}\selectfont\color{white}\MakeUppercase{Passe à OnBuch+}}\par\vspace{4pt}
        {\raggedright\popsemi\fontsize{8.4}{11}\selectfont\color{white}Tous les cours signés, Tuteur IA illimité, fiches PDF — onglet OnBuch+ de l'app\par}
      \end{minipage}
    \end{tcolorbox}
  \end{minipage}
  \vspace{8pt}
  \popcontenu
  \vfill
  \signatureonbuch
  \restoregeometry
  \newpage
}

% Rangée « Ce cours contient » (page de garde)
\newcommand{\poptile}[3]{% #1 couleur #2 gros texte #3 légende
  \begin{tcolorbox}[enhanced,colback=#1,colframe=popink,boxrule=2pt,arc=9pt,shadow={2.5pt}{-2.5pt}{0pt}{popink},
      left=6pt,right=6pt,top=9pt,bottom=9pt,halign=center,valign=center,height=1.75cm,width=\linewidth,nobeforeafter]
    {\popdisplay\fontsize{12.5}{13}\selectfont\color{white}\MakeUppercase{#2}}\par\vspace{3pt}
    {\centering\popsemi\fontsize{7.8}{9.5}\selectfont\color{white}#3\par}
  \end{tcolorbox}}
\newcommand{\popcontenu}{%
  \par\noindent{\popxboldls\fontsize{7.5}{7.5}\selectfont\color{popmuted}CE COURS SIGNÉ CONTIENT}\par\vspace{7pt}
  \noindent\begin{minipage}[t]{0.235\linewidth}\poptile{popblue}{Cours}{complet, illustré, pas à pas}\end{minipage}\hfill
  \begin{minipage}[t]{0.235\linewidth}\poptile{poporange}{Méthodes}{et exercices résolus}\end{minipage}\hfill
  \begin{minipage}[t]{0.235\linewidth}\poptile{poppink}{Exercices}{type examen + corrigés}\end{minipage}\hfill
  \begin{minipage}[t]{0.235\linewidth}\poptile{popgreen}{Fiche bilan}{l'essentiel à retenir}\end{minipage}\par}

% Sommaire
\newtcolorbox{sommaire}{enhanced,colback=white,colframe=popink,boxrule=2.2pt,arc=10pt,
  drop shadow=popink,left=18pt,right=18pt,top=13pt,bottom=13pt,before skip=6pt,after skip=20pt,
  before upper={\poppill{Sommaire}{poppurple}{white}\par\vspace{9pt}\popsemi\fontsize{10.6}{14.5}\selectfont\color{popink}}}

% Signature de fin de cours — poussée tout en bas de la dernière page
\newcommand{\findecours}{%
  \par\vfill\nopagebreak
  \begin{center}\popsquiggle{poporange}\end{center}\vspace{4pt}
  \signatureonbuch
}
\AtBeginDocument{\def\llbracket{\mathopen{[\mkern-3.5mu[}}\def\rrbracket{\mathclose{]\mkern-3.5mu]}}} % intervalles d'entiers (glyphe ⟦ absent de la police)
% Glyphes absents de Fira Math : redéfinis à partir de glyphes disponibles
\AtBeginDocument{%
  \def\setminus{\mathbin{\backslash}}\let\smallsetminus\setminus
  \def\vdots{\mathord{\vbox{\baselineskip3.2pt\lineskiplimit0pt\kern4pt\hbox{.}\hbox{.}\hbox{.}}}}%
  \def\ddots{\mathinner{\mkern1mu\raise6.5pt\hbox{.}\mkern2mu\raise3.6pt\hbox{.}\mkern2mu\raise0.7pt\hbox{.}\mkern1mu}}%
  \def\triangle{\mathord{\scalebox{0.9}{$\symup{\Delta}$}}}%
  \def\nabla{\mathord{\raisebox{\depth}{\scalebox{1}[-1]{$\symup{\Delta}$}}}}%
  \def\checkmark{\popcheck{popdark}}%
  \def\star{\mathbin{\ast}}\def\bigstar{\mathord{\ast}}%
  \def\diamond{\mathbin{\scalebox{0.6}{$\Diamond$}}}%
  \def\bigcup{\mathop{\mathchoice{\raisebox{-0.3em}{\scalebox{1.7}{$\cup$}}}{\scalebox{1.25}{$\cup$}}{\cup}{\cup}}\displaylimits}%
  \def\bigcap{\mathop{\mathchoice{\raisebox{-0.3em}{\scalebox{1.7}{$\cap$}}}{\scalebox{1.25}{$\cap$}}{\cap}{\cap}}\displaylimits}%
  \def\lesseqqgtr{\mathrel{\lessgtr}}\def\bigoplus{\mathop{\oplus}\displaylimits}%
}
\providecommand{\ssi}{si et seulement si} % abréviation fréquente
\AtBeginDocument{\providecommand{\wideparen}{\overparen}} % arc de cercle

%% FIGFIX-BEGIN v3 — correctifs globaux des figures (docs/onbuch-generation/tools/figfix)
\usepackage{adjustbox}\usepackage{etoolbox}
% toute figure TikZ/pgfplots plus large que la ligne est réduite à la largeur disponible
\BeforeBeginEnvironment{tikzpicture}{\begin{adjustbox}{max width=\linewidth}}
\AfterEndEnvironment{tikzpicture}{\end{adjustbox}}
% légendes pgfplots lisibles par-dessus les courbes
\pgfplotsset{every axis/.append style={legend style={fill=white,fill opacity=0.92,text opacity=1,draw=black!15,inner sep=3pt}}}
% étiquettes d'arêtes (syntaxe edge["..."]) avec fond blanc
\tikzset{every edge quotes/.append style={fill=white,inner sep=1.2pt}}
% bulles de cartes mentales : texte à la ligne dans la bulle, bulle un peu plus grande
\tikzset{every concept/.append style={text width=2.0cm,align=center,minimum size=2.3cm,inner sep=2pt}}
%% FIGFIX-END
\begin{document}

\pagegarde{Leçon 37 — Expression orale : internet}{Chinois}{1ère A}{Module 5 : Mass médias et télécommunication}{ZHA37}{2 h}{Cette leçon d'expression orale apprend à présenter un exposé simple et à participer à un débat guidé sur le thème d'internet en chinois. Elle fournit les formules fonctionnelles pour présenter, donner son avis et relancer la parole, ainsi qu'un entraînement ciblé sur les tons.
\vspace{4pt}
\begin{multicols}{2}\small
\begin{itemize}
\item À la fin de cette leçon, je sais utiliser le vocabulaire de base lié à internet (上网, 网站, 手机, 网络...)
\item je sais présenter oralement un exposé simple et structuré en chinois sur internet
\item je sais donner mon avis avec les formules 我觉得..., 我认为..., 对我来说...
\item je sais exprimer l'accord et le désaccord poliment (我同意, 我不同意, 你说得对...)
\item je sais relancer la parole et poser des questions dans un débat (你呢? 你觉得怎么样?)
\item je sais prononcer correctement les tons des mots clés de la leçon
\end{itemize}\end{multicols}}

\begin{sommaire}
\begin{multicols}{2}
\begin{enumerate}
\item Vocabulaire d'internet et mise en route orale
\item Dialogue modèle 1 : parler de ses habitudes
\item Formules fonctionnelles : présenter, opiner, relancer
\item Dialogue modèle 2 : débattre des avantages et inconvénients
\item Travail phonétique : les tons dans le débat
\item Jeux de rôle et débat guidé de classe
\item Activité d'intégration
\item Méthodes et exercices résolus
\item Exercices
\item Corrigés des exercices
\item Fiche bilan
\end{enumerate}
\end{multicols}
\end{sommaire}

\begin{objectifs}
\begin{itemize}
\item À la fin de cette leçon, je sais utiliser le vocabulaire de base lié à internet (上网, 网站, 手机, 网络...)
\item je sais présenter oralement un exposé simple et structuré en chinois sur internet
\item je sais donner mon avis avec les formules 我觉得..., 我认为..., 对我来说...
\item je sais exprimer l'accord et le désaccord poliment (我同意, 我不同意, 你说得对...)
\item je sais relancer la parole et poser des questions dans un débat (你呢? 你觉得怎么样?)
\item je sais prononcer correctement les tons des mots clés de la leçon
\item je sais comparer les avantages et les inconvénients d'internet à l'oral
\item je sais jouer un dialogue avec un camarade sur l'usage d'internet au quotidien
\end{itemize}
\end{objectifs}

\begin{prerequis}
\begin{itemize}
\item Se présenter et parler de ses habitudes quotidiennes (leçons de Seconde)
\item Les pronoms personnels et la structure sujet-verbe-complément
\item Les nombres et les expressions de fréquence (每天, 常常, 有时候)
\item Les verbes modaux 会, 可以, 应该
\item Le vocabulaire des médias vu dans le Module 5 (电视, 电话, 手机)
\end{itemize}
\end{prerequis}

\newpage

\coursec{Vocabulaire d'internet et mise en route orale}

\courssub{Situation d'accroche et problématique}

À Yaoundé comme à Douala, dès la sortie des cours, les élèves sortent leur téléphone : WhatsApp pour discuter avec les amis, TikTok pour rire, Facebook pour suivre l'actualité, Google pour réviser le bac. En Chine, c'est exactement la même scène, mais avec d'autres applications : les jeunes Chinois passent leur temps sur \zh{微信} (Wēixìn, WeChat) pour discuter et payer, et sur \zh{抖音} (Dǒuyīn, Douyin) pour regarder de courtes vidéos.

Mais imagine la scène suivante : un correspondant chinois de Shanghai te pose une question toute simple — \zh{你用互联网做什么？} (Nǐ yòng hùliánwǎng zuò shénme ? — Que fais-tu sur internet ?). Sauras-tu lui répondre en chinois, décrire tes habitudes, puis débattre avec lui des avantages et des inconvénients d'internet ? C'est exactement l'objectif de cette leçon : à la fin, tu sauras \cle{présenter un petit exposé oral} sur internet et \cle{défendre ton avis} dans un mini-débat.

Pour cela, la première étape est de construire ensemble le vocabulaire indispensable et de s'échauffer la voix : en chinois, on apprend à parler en parlant.

\courssub{Observation : un premier texte pour découvrir les mots}

Avant d'apprendre la liste de vocabulaire, lis d'abord ce petit texte. Souligne mentalement les mots que tu devines.

\begin{textebox}[Les habitudes de Li Ming — 李明的上网习惯]
我每天用手机上网。我常常在网上查信息，有时候跟朋友聊天。\\
Wǒ měitiān yòng shǒujī shàngwǎng. Wǒ chángcháng zài wǎngshang chá xìnxī, yǒushíhou gēn péngyou liáotiān.\\
Je me connecte chaque jour avec mon téléphone. Je cherche souvent des informations en ligne, parfois je discute avec des amis.\\[4pt]
我很少玩游戏，从来不下载电影。我觉得互联网很有用，因为我可以在网上学习汉语。\\
Wǒ hěn shǎo wán yóuxì, cónglái bù xiàzài diànyǐng. Wǒ juéde hùliánwǎng hěn yǒuyòng, yīnwèi wǒ kěyǐ zài wǎngshang xuéxí Hànyǔ.\\
Je joue rarement aux jeux, je ne télécharge jamais de films. Je trouve qu'internet est très utile, parce que je peux apprendre le chinois en ligne.
\end{textebox}

Observons ce texte. Trois faits sautent aux yeux :

\begin{enumerate}
\item Le mot \zh{网} (wǎng, « filet, réseau ») revient partout : \zh{上网}, \zh{网上}, \zh{互联网}. C'est la clé de voûte de tout le vocabulaire d'internet.
\item Les mots de fréquence (\zh{每天}, \zh{常常}, \zh{有时候}, \zh{很少}, \zh{从来不}) se placent \textbf{avant le verbe}, comme en français dans « je cherche \emph{souvent} ».
\item Pour dire « sur internet », le chinois dit \zh{在网上} (zài wǎngshang), littéralement « être sur le réseau », placé avant le verbe d'action.
\end{enumerate}

\courssub{Le vocabulaire d'internet : les mots clés}

Voici le lexique de base de la leçon. Apprends-le par petits groupes logiques : d'abord les \textbf{outils}, puis les \textbf{actions}, puis les \textbf{notions}.

\begin{center}
\begin{tabularx}{\linewidth}{|X|X|X|X|}
\hline
\rowcolor{popblueL} Caractères & Pinyin & Nature & Traduction \\
\hline
互联网 & hùliánwǎng & nom & internet (le réseau) \\
\hline
上网 & shàngwǎng & verbe (séparable) & se connecter à internet \\
\hline
网站 & wǎngzhàn & nom & site web \\
\hline
手机 & shǒujī & nom & téléphone portable \\
\hline
电脑 & diànnǎo & nom & ordinateur \\
\hline
微信 & Wēixìn & nom propre & WeChat \\
\hline
信息 & xìnxī & nom & information, message \\
\hline
聊天 & liáotiān & verbe (séparable) & bavarder, chatter \\
\hline
\end{tabularx}
\end{center}

\begin{center}
\begin{tabularx}{\linewidth}{|X|X|X|X|}
\hline
\rowcolor{popblueL} Caractères & Pinyin & Nature & Traduction \\
\hline
查 & chá & verbe & chercher, consulter, vérifier \\
\hline
发 & fā & verbe & envoyer (message, photo) \\
\hline
下载 & xiàzài & verbe & télécharger \\
\hline
玩 & wán & verbe & jouer, s'amuser \\
\hline
学习 & xuéxí & verbe & apprendre, étudier \\
\hline
用 & yòng & verbe & utiliser \\
\hline
资料 & zīliào & nom & documents, données \\
\hline
有用 & yǒuyòng & adjectif & utile \\
\hline
\end{tabularx}
\end{center}

Quelques remarques de formation des mots, pour t'aider à mémoriser :

\begin{itemize}
\item \zh{手机} (shǒujī) : \zh{手} « main » + \zh{机} « machine » = la « machine de la main ». Logique, non ?
\item \zh{电脑} (diànnǎo) : \zh{电} « électricité » + \zh{脑} « cerveau » = le « cerveau électrique ».
\item \zh{互联网} (hùliánwǎng) : \zh{互} « mutuel » + \zh{联} « relier » + \zh{网} « réseau » = « réseau d'interconnexion ».
\item \zh{网站} (wǎngzhàn) : \zh{网} « réseau » + \zh{站} « station » = « station du réseau », c'est-à-dire un site web.
\end{itemize}

\begin{popfigure}
\begin{tikzpicture}[mindmap, grow cyclic, every node/.style={concept}, concept color=popblue!40,
root concept/.append style={concept color=popblue!70, font=\large\bfseries},
level 1/.append style={level distance=3.2cm, sibling angle=120},
level 2/.append style={level distance=2.2cm, sibling angle=50}]
\node[root concept, align=center]{互联网\\ hùliánwǎng}
  child[concept color=popgreen!50]{ node[align=center]{工具\\ gōngjù\\ (outils)}
    child { node[align=center]{手机\\ shǒujī} }
    child { node[align=center]{电脑\\ diànnǎo} }
  }
  child[concept color=poporange!50]{ node[align=center]{活动\\ huódòng\\ (activités)}
    child { node[align=center]{聊天\\ liáotiān} }
    child { node[align=center]{学习\\ xuéxí} }
    child { node[align=center]{玩游戏\\ wán yóuxì} }
  }
  child[concept color=poppink!50]{ node[align=center]{评价\\ píngjià\\ (jugement)}
    child { node[align=center]{好处\\ hǎochu} }
    child { node[align=center]{坏处\\ huàichu} }
  };
\end{tikzpicture}
\legende{Carte mentale du vocabulaire d'internet : les outils, les activités et le jugement.}
\end{popfigure}

\courssub{Grammaire 1 : le verbe séparable 上网}

\begin{definition}[上网 shàngwǎng : un verbe séparable]
\zh{上网} signifie « se connecter à internet ». C'est un \cle{verbe séparable} : \zh{上} est le verbe (« monter ») et \zh{网} est son objet figé (« le réseau »). Les deux caractères fonctionnent ensemble, mais on peut les \textbf{séparer} pour insérer une durée, \zh{了} ou \zh{过} : \zh{上两个小时网} (shàng liǎng ge xiǎoshí wǎng) — « se connecter pendant deux heures ».
\end{definition}

\begin{propriete}[Structure : verbe séparable + durée]
Schéma : \textbf{Verbe + (了 / 过 / durée) + Objet figé}\\[2pt]
\zh{我今天上了三个小时网。} — Wǒ jīntiān shàng le sān ge xiǎoshí wǎng. — « Aujourd'hui, je me suis connecté pendant trois heures. »\\[2pt]
\zh{你上过这个网站吗？} — Nǐ shàngguo zhège wǎngzhàn ma ? — « As-tu déjà visité ce site ? »\\[2pt]
Même phénomène avec \zh{聊天} : \zh{我们聊了半个小时天。} — Wǒmen liáo le bàn ge xiǎoshí tiān. — « Nous avons bavardé pendant une demi-heure. »
\end{propriete}

\begin{attention}[Piège : ne colle jamais la durée derrière 上网]
Les francophones traduisent mot à mot « je me connecte deux heures » par *\zh{我上网两个小时}. Cette phrase peut se comprendre à l'oral, mais elle est maladroite et sanctionnée à l'écrit scolaire ; la forme correcte et élégante est \zh{我上两个小时网} ou \zh{我上网上了两个小时}. Retiens : la durée \textbf{entre} dans le verbe séparable, elle ne le suit pas.
\end{attention}

\courssub{Grammaire 2 : dire « sur internet » — la structure 在网上 + verbe}

Observons deux phrases du texte : \zh{我常常在网上查信息} et \zh{我可以在网上学习汉语}. Dans les deux cas, le lieu de l'action (« sur internet ») est exprimé par \zh{在网上} placé \textbf{avant le verbe}.

\begin{propriete}[La structure 在网上 + verbe]
Schéma : \textbf{Sujet + 在网上 + verbe (+ complément)}\\[2pt]
En français, le complément de lieu « sur internet » se met après le verbe (« je cherche des documents \emph{sur internet} »). En chinois, il se met \textbf{obligatoirement avant le verbe}. C'est la règle générale : en chinois, les compléments de lieu introduits par \zh{在} précèdent toujours le verbe d'action.
\end{propriete}

\begin{exemplebox}[Exemples traduits]
\zh{我在网上查资料。} — Wǒ zài wǎngshang chá zīliào. — « Je cherche des documents sur internet. »\\[3pt]
\zh{他很少玩游戏。} — Tā hěn shǎo wán yóuxì. — « Il joue rarement aux jeux. »\\[3pt]
\zh{她在网上学习法语。} — Tā zài wǎngshang xuéxí Fǎyǔ. — « Elle apprend le français sur internet. »\\[3pt]
\zh{我们在网上发信息。} — Wǒmen zài wǎngshang fā xìnxī. — « Nous envoyons des messages sur internet. »
\end{exemplebox}

\begin{attention}[Erreur fréquente des francophones]
Ne traduis pas « sur internet » par un complément placé en fin de phrase : *\zh{我学习在网上} est une faute grave (ordre des mots français calqué sur le chinois). La structure naturelle est \textbf{\zh{在网上} + verbe}, placée avant le verbe : \zh{我在网上学习。} De plus, à l'oral courant, on préfère la forme courte \zh{网上} à la forme complète \zh{互联网上}, qui appartient plutôt au style écrit ou journalistique.
\end{attention}

\courssub{Grammaire 3 : l'échelle des fréquences}

Pour parler de tes habitudes, tu as besoin des adverbes de fréquence. Tous se placent \textbf{avant le verbe} (et \zh{每天} peut aussi ouvrir la phrase).

\begin{center}
\begin{tabularx}{\linewidth}{|X|X|X|X|}
\hline
\rowcolor{popblueL} Caractères & Pinyin & Traduction & Fréquence approximative \\
\hline
每天 & měitiān & chaque jour & 100 \% \\
\hline
常常 & chángcháng & souvent & environ 80 \% \\
\hline
有时候 & yǒushíhou & parfois & environ 50 \% \\
\hline
很少 & hěn shǎo & rarement & environ 10 \% \\
\hline
从来不 & cónglái bù & ne... jamais & 0 \% \\
\hline
\end{tabularx}
\end{center}

\begin{propriete}[Emploi et position des adverbes de fréquence]
Schéma : \textbf{Sujet + adverbe de fréquence + verbe}\\[2pt]
\zh{我每天上网。} — Wǒ měitiān shàngwǎng. — « Je me connecte chaque jour. »\\[2pt]
\zh{他从来不玩游戏。} — Tā cónglái bù wán yóuxì. — « Il ne joue jamais aux jeux. »\\[2pt]
Remarque : \zh{从来不} porte déjà sa négation \zh{不} ; n'ajoute jamais un deuxième \zh{不} (*\zh{他从来不不去} est impossible). \zh{每天} et \zh{有时候} peuvent aussi se placer en tête de phrase : \zh{有时候我跟朋友聊天。} — « Parfois, je discute avec des amis. »
\end{propriete}

\begin{exoresolu}[Construire une phrase de fréquence]
Énoncé : traduis en chinois « Mon frère se connecte rarement avec l'ordinateur ; il discute souvent avec ses amis sur le téléphone. »
\tcblower
Solution détaillée :
\begin{enumerate}
\item « Mon frère » = \zh{我哥哥} (wǒ gēge). « Se connecte » = \zh{上网}.
\item « Rarement » = \zh{很少}, placé \textbf{avant} le verbe : \zh{我哥哥很少上网}.
\item « Avec l'ordinateur » = \zh{用电脑} (yòng diànnǎo), complément de moyen placé avant le verbe : \zh{我哥哥很少用电脑上网。}
\item « Il discute souvent avec ses amis » : \zh{常常} avant le verbe, \zh{跟朋友} avant le verbe aussi : \zh{他常常跟朋友聊天。}
\end{enumerate}
Phrase complète : \zh{我哥哥很少用电脑上网，他常常跟朋友聊天。}\\
Wǒ gēge hěn shǎo yòng diànnǎo shàngwǎng, tā chángcháng gēn péngyou liáotiān.
\end{exoresolu}

\courssub{Mise en route orale : l'échauffement en chaîne}

Passons maintenant à la bouche ! L'objectif de cette fin de section est de faire parler chaque élève au moins deux fois, avec des phrases très courtes.

\begin{experience}[Échauffement en chaîne : 你常常上网吗？]
\begin{enumerate}
\item Le professeur pose la question \zh{你常常上网吗？} (Nǐ chángcháng shàngwǎng ma ? — Est-ce que tu te connectes souvent ?) au premier élève.
\item L'élève répond avec un adverbe de fréquence : \zh{我常常上网。} / \zh{我很少上网。} / \zh{我从来不上网。}
\item Puis il se tourne vers son voisin et pose la même question. La chaîne continue jusqu'à ce que toute la classe ait parlé.
\item Variante pour les plus rapides : ajouter un outil — \zh{我每天用手机上网。} (Je me connecte chaque jour avec mon téléphone.)
\end{enumerate}
\end{experience}

\begin{experience}[Mini-sondage de classe : 你用互联网做什么？]
\begin{enumerate}
\item Chaque élève interroge trois camarades : \zh{你用互联网做什么？} (Nǐ yòng hùliánwǎng zuò shénme ? — Que fais-tu sur internet ?)
\item Les réponses possibles : \zh{我查资料。} (Je cherche des documents.) / \zh{我跟朋友聊天。} (Je discute avec des amis.) / \zh{我学习汉语。} (J'apprends le chinois.) / \zh{我玩游戏。} (Je joue.) / \zh{我下载音乐。} (Je télécharge de la musique.)
\item Chacun rapporte ensuite un résultat à la classe : \zh{玛丽常常在网上学习，她从来不玩游戏。} (Mǎlì chángcháng zài wǎngshang xuéxí, tā cónglái bù wán yóuxì. — Marie étudie souvent en ligne, elle ne joue jamais.)
\end{enumerate}
\end{experience}

\begin{methode}[Bien prononcer pendant l'échauffement]
\begin{enumerate}
\item Articule les tons avant la vitesse : \zh{shàngwǎng} = quatrième ton descendant + troisième ton plongeant ; \zh{chángcháng} = deux deuxièmes tons montants.
\item Dans \zh{有时候} (yǒushíhou), le dernier caractère est au ton neutre : dis-le court et léger.
\item Dans \zh{从来不上网}, enchaîne sans pause : cónglái bù shàngwǎng — le rythme est régulier, syllabe par syllabe.
\end{enumerate}
\end{methode}

\begin{savaistu}[WeChat, le couteau suisse chinois]
En Chine, personne n'utilise WhatsApp : on utilise \zh{微信} (Wēixìn, WeChat). Cette application sert à discuter, mais aussi à payer au marché, commander un repas, appeler un taxi ou prendre un rendez-vous chez le médecin. Au Cameroun, le mobile money (Orange Money, MTN MoMo) et WhatsApp sont deux applications séparées ; en Chine, WeChat réunit tout en une seule. Quand un Chinois te demande \zh{你有微信吗？} (Nǐ yǒu Wēixìn ma ? — Tu as WeChat ?), c'est l'équivalent de « je te donne mon numéro ? ».
\end{savaistu}

\begin{aretenir}[L'essentiel de la section]
\begin{itemize}
\item Le vocabulaire d'internet s'organise autour de \zh{网} : \zh{互联网} (le réseau), \zh{上网} (se connecter), \zh{网站} (site web), \zh{在网上} (sur internet).
\item \zh{上网} et \zh{聊天} sont des verbes séparables : la durée s'insère au milieu (\zh{上两个小时网}).
\item « Sur internet » se dit \zh{在网上}, placé \textbf{avant} le verbe : \zh{我在网上查资料。}
\item L'échelle des fréquences : \zh{每天} $>$ \zh{常常} $>$ \zh{有时候} $>$ \zh{很少} $>$ \zh{从来不}, tous placés avant le verbe.
\item Deux questions à maîtriser à l'oral : \zh{你常常上网吗？} et \zh{你用互联网做什么？}
\end{itemize}
\end{aretenir}

\coursec{Dialogue modèle 1 : parler de ses habitudes}

Après avoir découvert le vocabulaire essentiel d'internet dans la section précédente, nous allons maintenant l'ancrer dans une situation de communication authentique. Observer un dialogue naturel entre deux lycéens permet de voir comment s'enchaînent les questions de fréquence, les verbes d'action et le choix de l'outil numérique. C'est la base de toute expression orale fluide sur ce thème.

\courssub{Observation et écoute du dialogue}

Écoutons d'abord le dialogue dans son intégralité. L'élève A interroge l'élève B sur sa fréquence de connexion, la nature de ses activités et le support utilisé. Remarquez l'ordre naturel des échanges : fréquence $\rightarrow$ activités $\rightarrow$ outil.

\begin{textebox}[Dialogue : 你常常上网吗?]
A: 你常常上网吗? \\
B: 对，我每天都上网。 \\
A: 你用互联网做什么? \\
B: 我常常查信息、学习汉语，有时候看视频。 \\
A: 你用电脑还是手机? \\
B: 我一般用手机，很方便。
\end{textebox}

\begin{center}
\begin{tabularx}{\linewidth}{|X|X|X|X|}
\hline
\rowcolor{popblueL} \textbf{Caractères} & \textbf{Pinyin} & \textbf{Nature} & \textbf{Traduction} \\
\hline
常常 & chángcháng & adv. & souvent, fréquemment \\
\hline
上网 & shàngwǎng & VO & se connecter (à internet), aller en ligne \\
\hline
每天 & měitiān & adv. & chaque jour, tous les jours \\
\hline
都 & dōu & adv. & tous (marque la totalité) \\
\hline
用 & yòng & v. & utiliser, employer \\
\hline
互联网 & hùliánwǎng & n. & internet (terme formel) \\
\hline
做什么 & zuò shénme & loc. & faire quoi, à quoi sert... \\
\hline
查信息 & chá xìnxī & VO & chercher des informations \\
\hline
学习 & xuéxí & v. & étudier, apprendre \\
\hline
汉语 & Hànyǔ & n. & chinois (langue) \\
\hline
有时候 & yǒushíhou & adv. & parfois, de temps en temps \\
\hline
看视频 & kàn shìpín & VO & regarder des vidéos \\
\hline
电脑 & diànnǎo & n. & ordinateur \\
\hline
还是 & háishì & conj. & ou (dans une question alternative) \\
\hline
手机 & shǒujī & n. & téléphone portable, smartphone \\
\hline
一般 & yìbān & adv. & généralement, d'ordinaire \\
\hline
方便 & fāngbiàn & adj. & pratique, commode \\
\hline
\end{tabularx}
\end{center}

\courssub{Analyse phrase par phrase et travail des tons}

Travaillez la prononciation phrase par phrase. Le chinois est une langue à tons : une erreur de ton change le sens du mot. Portez une attention particulière aux paires tonales ci-dessous.

\begin{exemplebox}[Phrase 1 : La question de fréquence]
\textbf{A:} \zh{你常常上网吗?} \\
\textbf{Pinyin:} Nǐ chángcháng shàngwǎng ma? \\
\textbf{Traduction:} Est-ce que tu te connectes souvent (à internet) ?
\begin{itemize}
    \item \cle{常常 chángcháng} : 2e ton + 2e ton. La voix monte deux fois.
    \item \cle{上网 shàngwǎng} : \textbf{4e ton (shàng) + 3e ton (wǎng)}. \zh{上} part haut et tombe sèchement ; \zh{网} descend puis remonte (ton creusé).
    \item \cle{吗 ma} : ton neutre, léger, sans accent.
\end{itemize}
\end{exemplebox}

\begin{exemplebox}[Phrase 2 : La confirmation + fréquence]
\textbf{B:} \zh{对，我每天都上网。} \\
\textbf{Pinyin:} Duì, wǒ měitiān dōu shàngwǎng. \\
\textbf{Traduction:} Oui, je me connecte tous les jours.
\begin{itemize}
    \item \cle{对 duì} : 4e ton. Court, affirmatif.
    \item \cle{每天 měitiān} : 3e ton + 1e ton. \zh{每} (3e ton) + \zh{天} (1e ton haut et plat). \textbf{Attention :} \zh{每} est un 3e ton isolé devant un 1e ton $\rightarrow$ il se prononce en demi-3e ton (bas, 21) sans remontée.
    \item \cle{都 dōu} : 1e ton. Place \textbf{avant} le verbe \zh{上网}.
    \item Structure clé : \textbf{Sujet + 每天 + 都 + Verbe}.
\end{itemize}
\end{exemplebox}

\begin{exemplebox}[Phrase 3 : La question sur l'activité]
\textbf{A:} \zh{你用互联网做什么?} \\
\textbf{Pinyin:} Nǐ yòng hùliánwǎng zuò shénme? \\
\textbf{Traduction:} Qu'est-ce que tu fais avec internet ? / À quoi utilises-tu internet ?
\begin{itemize}
    \item \cle{用 yòng} : 4e ton. Verbe « utiliser ».
    \item \cle{互联网 hùliánwǎng} : 4e + 2e + 3e ton. \zh{互} (hù) tombe, \zh{联} (lián) monte, \zh{网} (wǎng) creuse.
    \item \cle{做什么 zuò shénme} : \zh{做} (4e ton) + \zh{什} (2e ton, shén ici car suivi de me ton neutre) + \zh{么} (ton neutre). Pas de \zh{吗} en fin de phrase car \zh{什么} porte l'interrogation.
\end{itemize}
\end{exemplebox}

\begin{exemplebox}[Phrase 4 : L'énumération d'activités]
\textbf{B:} \zh{我常常查信息、学习汉语，有时候看视频。} \\
\textbf{Pinyin:} Wǒ chángcháng chá xìnxī, xuéxí Hànyǔ, yǒushíhou kàn shìpín. \\
\textbf{Traduction:} Je cherche souvent des infos, j'étudie le chinois, parfois je regarde des vidéos.
\begin{itemize}
    \item \cle{查信息 chá xìnxī} : 2e + 4e + 1e ton. \zh{查} (chercher) monte ; \zh{信} (xìn) tombe ; \zh{息} (xī) haut et plat.
    \item \cle{学习汉语 xuéxí Hànyǔ} : 2e + 2e + 4e + 3e ton. \zh{学} (xué) monte, \zh{习} (xí) monte, \zh{汉} (Hàn) tombe, \zh{语} (yǔ) creuse.
    \item \cle{有时候 yǒushíhou} : 3e + 2e + ton neutre. \zh{有} (yǒu) 3e ton $\rightarrow$ demi-3e ton devant \zh{时} (2e ton) ; \zh{时} (shí) 2e ton ; \zh{hou} léger.
    \item \cle{看视频 kàn shìpín} : 4e + 4e + 2e ton. \zh{看} (kàn) tombe ; \zh{视} (shì) tombe ; \zh{频} (pín) monte.
    \item Ponctuation : le \textbf{、} (dùnhào) sépare les items d'une énumération.
\end{itemize}
\end{exemplebox}

\begin{exemplebox}[Phrase 5 : La question alternative (choix de l'outil)]
\textbf{A:} \zh{你用电脑还是手机?} \\
\textbf{Pinyin:} Nǐ yòng diànnǎo háishì shǒujī? \\
\textbf{Traduction:} Tu utilises l'ordinateur ou le téléphone ?
\begin{itemize}
    \item \cle{电脑 diànnǎo} : \textbf{4e ton (diàn) + 3e ton (nǎo)}. \zh{电} tombe sèchement ; \zh{脑} creuse.
    \item \cle{还是 háishì} : 2e + 4e ton. \zh{还} (hái) 2e ton ; \zh{是} (shì) 4e ton. Forme de référence : \textbf{háishì}. À l'oral rapide, on entend parfois \textbf{háishǐ} (2-3), mais la norme est 2-4.
    \item \cle{手机 shǒujī} : \textbf{3e ton (shǒu) + 1e ton (jī)}. Voir règle de sandhi ci-dessous.
    \item \textbf{Pas de 吗 à la fin !} La présence de \zh{还是} marque la question alternative.
\end{itemize}
\end{exemplebox}

\begin{exemplebox}[Phrase 6 : La réponse + avis]
\textbf{B:} \zh{我一般用手机，很方便。} \\
\textbf{Pinyin:} Wǒ yìbān yòng shǒujī, hěn fāngbiàn. \\
\textbf{Traduction:} J'utilise généralement le téléphone, c'est très pratique.
\begin{itemize}
    \item \cle{一般 yìbān} : 4e + 1e ton. \zh{一} (yī) $\rightarrow$ \textbf{yì} (4e ton) devant un 4e ton (\zh{般} bān). Règle de sandhi de \zh{一}.
    \item \cle{很方便 hěn fāngbiàn} : \zh{很} (hěn 3e ton) + \zh{方} (fāng 1e ton) + \zh{便} (biàn 4e ton). \textbf{Sandhi du 3e ton :} \zh{很} $\rightarrow$ \textbf{hén} (2e ton) devant \zh{方} (1e ton). On prononce \textbf{hén fāngbiàn}.
    \item \cle{方便 fāngbiàn} : 1e + 4e ton. Adjectif predicatif (pas de \zh{是} devant).
\end{itemize}
\end{exemplebox}

\courssub{Point de phonétique : Le sandhi du 3e ton (三声变调)}

C'est la règle phonétique \textbf{la plus importante} pour ne pas avoir un accent « haché ». Un 3e ton (ton creusé 214) ne se prononce intégralement que s'il est \textbf{isolé} ou en \textbf{fin de groupe rythmique}. Devant un autre ton, il change.

\begin{propriete}[Règle du sandhi du 3e ton]
\begin{enumerate}
    \item \textbf{3e ton + 3e ton} $\rightarrow$ Le \textbf{premier} devient \textbf{2e ton} (montant 35). Ex: \zh{你好} nǐhǎo $\rightarrow$ \textbf{níhǎo}.
    \item \textbf{3e ton + 1e / 2e / 4e ton / ton neutre} $\rightarrow$ Le 3e ton devient \textbf{demi-3e ton} (bas, 21), sans la remontée finale. Il reste bas.
\end{enumerate}
\end{propriete}

\begin{exemplebox}[Application au dialogue]
\begin{itemize}
    \item \zh{手机 shǒujī} (3e + 1e) : \zh{手} (shǒu) $\rightarrow$ \textbf{demi-3e ton (bas)}. Pas de remontée. \zh{机} (jī) 1e ton haut. \rightarrow \textbf{shǒu(↘)jī(↗)}.
    \item \zh{很好 hěn hǎo} (3e + 3e) : \zh{很} (hěn) $\rightarrow$ \textbf{2e ton (hén)}. \zh{好} (hǎo) reste 3e ton. \rightarrow \textbf{hén hǎo}.
    \item \zh{每天 měitiān} (3e + 1e) : \zh{每} (měi) $\rightarrow$ \textbf{demi-3e ton (bas)}. \zh{天} (tiān) 1e ton. \rightarrow \textbf{měi(↘)tiān(↗)}.
    \item \zh{有空 yǒukòng} (3e + 4e) : \zh{有} $\rightarrow$ \textbf{demi-3e ton}. \zh{空} 4e ton.
\end{itemize}
\end{exemplebox}

\begin{attention}[Piège fréquent : \zh{还是} vs \zh{或者}]
Les francophones confondent souvent « ou » interrogatif et « ou » affirmatif.
\begin{itemize}
    \item \cle{还是 háishì} : \textbf{Uniquement dans les questions alternatives}. \zh{你喝茶还是咖啡?} (Thé ou café ?)
    \item \cle{或者 huòzhě} : \textbf{Dans les phrases affirmatives (ou négatives)}. \zh{我喝茶或者咖啡都行。} (Je bois du thé ou du café, les deux vont.)
    \item \textbf{Erreur typique :} \zh{*我用电脑还是手机。} (Incorrect en affirmation). \textbf{Correct :} \zh{我用电脑或者手机。}
\end{itemize}
\end{attention}

\courssub{Structure grammaticale : La question alternative avec \zh{还是}}

\begin{propriete}[Structure de la question alternative]
\textbf{Sujet + Verbe + Option A + 还是 + Option B ?}
\begin{itemize}
    \item Pas de particule \zh{吗} à la fin.
    \item L'intonation monte légèrement sur la dernière option (B) puis tombe.
    \item La réponse attendue répète l'option choisie (pas juste « oui/non »).
\end{itemize}
\end{propriete}

\begin{center}
\begin{tabularx}{\linewidth}{|X|X|X|}
\hline
\rowcolor{popblueL} \textbf{Structure} & \textbf{Exemple (Caractères + Pinyin)} & \textbf{Traduction} \\
\hline
Suj + V + Nom + 还是 + Nom & 你用 \zh{电脑还是手机}? (Nǐ yòng diànnǎo háishì shǒujī?) & Tu utilises l'ordi ou le tel ? \\
\hline
Suj + V + Verbe + 还是 + Verbe & 我们 \zh{去还是不去}? (Wǒmen qù háishì bù qù?) & On y va ou pas ? \\
\hline
Suj + Adj + 还是 + Adj & 这个 \zh{贵还是便宜}? (Zhège guì háishì piányi?) & Celui-ci est cher ou bon marché ? \\
\hline
\end{tabularx}
\end{center}

\courssub{Exercice de substitution guidée (Entraînement oral)}

Remplacez les éléments soulignés dans le dialogue par vos propres habitudes. Travaillez en binôme : A pose les questions, B répond personnellement.

\begin{exoresolu}[Substitution : Personnaliser le dialogue]
\textbf{Modèle de dialogue à personnaliser :}
A: 你常常上网吗?
B: 对，我每天都上网。 / 不，我不常常上网。 / 我一周上网两次。
A: 你用互联网做什么?
B: 我常常 \uline{查信息、学习汉语}，有时候 \uline{看视频}。
A: 你用电脑还是手机?
B: 我一般用 \uline{手机}，很方便。

\tcblower

\textbf{Propositions de substitution (élève B) :}
\begin{itemize}
    \item \textbf{Activités :} \zh{听音乐} (tīng yīnyuè - écouter de la musique), \zh{玩游戏} (wán yóuxì - jouer à des jeux), \zh{看新闻} (kàn xīnwén - lire les infos), \zh{聊天} (liáotiān - chatter/discuter), \zh{发邮件} (fā yóujiàn - envoyer des mails), \zh{买东西} (mǎi dōngxi - faire des achats).
    \item \textbf{Fréquence :} \zh{一周三次} (yī zhōu sān cì - trois fois par semaine), \zh{每天晚上} (měitiān wǎnshang - tous les soirs), \zh{有空的时候} (yǒu kòng de shíhou - quand j'ai du temps).
    \item \textbf{Outil :} \zh{电脑} (diànnǎo), \zh{平板电脑} (píngbǎn diànnǎo - tablette), \zh{学校的电脑} (xuéxiào de diànnǎo - l'ordi de l'école).
    \item \textbf{Avis :} \zh{很快} (hěn kuài - très rapide), \zh{很慢} (hěn màn - très lent), \zh{不太方便} (bútài fāngbiàn - pas très pratique).
\end{itemize}

\textbf{Exemple de production élève :}
A: 你常常上网吗?
B: 对，我\zh{每天晚上}都上网。
A: 你用互联网做什么?
B: 我常常\zh{听音乐、聊天}，有时候\zh{看新闻}。
A: 你用电脑还是手机?
B: 我一般用\zh{手机}，\zh{很方便}。 / 我一般用\zh{电脑}，\zh{很快}。
\end{exoresolu}

\courssub{Méthode pour mémoriser et jouer le dialogue}

\begin{methode}[4 étapes pour s'approprier un dialogue oral]
\begin{enumerate}
    \item \textbf{Écoute globale :} Écoutez l'audio (ou le prof) 2 fois sans regarder le texte. Identifiez le sujet (internet), les interlocuteurs (2 élèves), le ton (décontracté).
    \item \textbf{Répétition en chœur (影子跟读) :} Répétez \textbf{simultanément} que le modèle (prof ou audio), phrase par phrase. Imitez la mélodie, les tons, la liaison \zh{很方便} (hén fāngbiàn). Ne lisez pas, \textbf{écoutez et parlez}.
    \item \textbf{Récitation à deux (对练) :} Fermez le livre. Jouez le dialogue avec un camarade. A = questions, B = réponses. Inversez les rôles. Si vous bloquez, rouvrez 3 secondes, refermez, continuez.
    \item \textbf{Jeu sans texte + substitution (表演+替换) :} Jouez la scène (gestes : mime le téléphone, l'ordi). Ensuite, B répond \textbf{vraiment} pour lui/elle (voir exercice de substitution ci-dessus). A réagit naturellement : \zh{啊，我也喜欢看视频!} (Ah, moi aussi j'aime regarder des vidéos !).
\end{enumerate}
\end{methode}

\courssub{Visualisation : Schéma de la structure du dialogue}

Le diagramme ci-dessous montre l'enchaînement logique des actes de parole (Question $\rightarrow$ Réponse $\rightarrow$ Relance). Cela aide à structurer sa prise de parole pour l'exposé ou le débat ultérieur.

\begin{popfigure}
\begin{tikzpicture}[
    node distance=1.2cm and 1.5cm,
    box/.style={draw=popblue, thick, rounded corners, fill=popblueL, text width=3.2cm, align=center, font=\small, minimum height=1.8cm},
    arrow/.style={-Stealth, thick, popdark},
    labelbox/.style={font=\tiny, text=popmuted, anchor=north}
]
% Nœuds
\node[box, align=center] (q1){\textbf{A: Question fréquence}\\ \zh{你常常上网吗?}};
\node[box, right=of q1, align=center] (r1){\textbf{B: Réponse + Fréquence}\\ \zh{对，我每天都上网。}};
\node[box, right=of r1, align=center] (q2){\textbf{A: Relance activité}\\ \zh{你用互联网做什么?}};
\node[box, right=of q2, align=center] (r2){\textbf{B: Énumération activités}\\ \zh{查信息、学习汉语、看视频}};
\node[box, right=of r2, align=center] (q3){\textbf{A: Question outil (choix)}\\ \zh{你用电脑还是手机?}};
\node[box, right=of q3, align=center] (r3){\textbf{B: Choix + Justification}\\ \zh{我一般用手机，很方便。}};

% Flèches
\draw[arrow] (q1.east) -- (r1.west) node[midway, above, labelbox] {Confirmation};
\draw[arrow] (r1.east) -- (q2.west) node[midway, above, labelbox] {Approfondir};
\draw[arrow] (q2.east) -- (r2.west) node[midway, above, labelbox] {Lister (、)};
\draw[arrow] (r2.east) -- (q3.west) node[midway, above, labelbox] {Préciser};
\draw[arrow] (q3.east) -- (r3.west) node[midway, above, labelbox] {Choisir + Juger};

% Étiquettes d'actes de parole sous les boîtes
\node[labelbox, below=0.2cm of q1.south] {Interrogation};
\node[labelbox, below=0.2cm of r1.south] {Information};
\node[labelbox, below=0.2cm of q2.south] {Interrogation ouverte};
\node[labelbox, below=0.2cm of r2.south] {Description};
\node[labelbox, below=0.2cm of q3.south] {Alternative (还是)};
\node[labelbox, below=0.2cm of r3.south] {Opinion (Adj)}; 
\end{tikzpicture}
\legende{Structure dialogique : enchaînement fréquence $\rightarrow$ activités $\rightarrow$ outil/avis.}
\end{popfigure}

\courssub{Entraînement phonétique ciblé : Les paires tonales du dialogue}

Isolez ces mots « pièges » pour la dictée tonale ou l'enregistrement individuel. Le professeur dicte le pinyin sans tons, l'élève écrit les tons ; ou l'élève lit la liste en exagérant la courbe de pitch.

\begin{center}
\begin{tabularx}{\linewidth}{|X|X|X|X|}
\hline
\rowcolor{poporangeL} \textbf{Mot} & \textbf{Pinyin brut} & \textbf{Tons réels (Sandhi appliqué)} & \textbf{Note} \\
\hline
上网 & shang wang & \textbf{shàng wǎng} (4-3) & 3e ton final = ton creusé complet \\
\hline
手机 & shou ji & \textbf{shǒu(↘) jī} (3-1) & 3e ton $\rightarrow$ demi-3e (bas) devant 1e ton \\
\hline
电脑 & dian nao & \textbf{diàn nǎo} (4-3) & 3e ton final = creusé \\
\hline
聊天 & liao tian & \textbf{liáo tiān} (2-1) & Deux tons montants/hauts : mélodie joyeuse \\
\hline
方便 & fang bian & \textbf{fāng biàn} (1-4) & 1e ton haut stable + 4e ton chute \\
\hline
很好 & hen hao & \textbf{hén hǎo} (2-3) & \textbf{Sandhi 3+3} : 1er 3e $\rightarrow$ 2e ton \\
\hline
一般 & yi ban & \textbf{yì bān} (4-1) & \textbf{Sandhi de 一} : yī $\rightarrow$ yì devant 4e ton \\
\hline
有时候 & you shi hou & \textbf{yǒu shí hou} (3-2-neutre) & 3e ton $\rightarrow$ demi-3e devant 2e ton \\
\hline
\end{tabularx}
\end{center}

\courssub{Exemples supplémentaires pour l'expression libre}

Utilisez ces phrases modèles pour varier vos réponses lors du jeu de rôle final.

\begin{exemplebox}[Variantes de réponse (Niveau HSK 2-3)]
\begin{itemize}
    \item \textbf{Fréquence :}
    \zh{我一周上网两三次。} (Wǒ yī zhōu shàngwǎng liǎng sān cì.) — Je me connecte 2 ou 3 fois par semaine.
    \zh{我只有周末才上网。} (Wǒ zhǐyǒu zhōumò cái shàngwǎng.) — Je ne me connecte que le week-end. (\zh{才} = seulement, pas avant ce moment).
    
    \item \textbf{Activités (structure : 常常 V1、V2，有时候 V3) :}
    \zh{我常常看中文视频、听中文歌，有时候写作业。} (Wǒ chángcháng kàn Zhōngwén shìpín, tīng Zhōngwén gē, yǒushíhou xiě zuòyè.) — Je regarde souvent des vidéos chinoises, j'écoute des chansons chinoises, parfois je fais mes devoirs.
    \zh{我常常跟朋友聊天，有时候玩游戏。} (Wǒ chángcháng gēn péngyou liáotiān, yǒushíhou wán yóuxì.) — Je chatte souvent avec des amis, parfois je joue.
    
    \item \textbf{Outil + Avis (结构 : 一般用..., 很/不太 + Adj) :}
    \zh{我一般用电脑，屏幕大，看视频很舒服。} (Wǒ yìbān yòng diànnǎo, píngmù dà, kàn shìpín hěn shūfu.) — J'utilise généralement l'ordi, l'écran est grand, regarder des vidéos est confortable.
    \zh{我有时用平板，比较轻，方便带。} (Wǒ yǒushí yòng píngbǎn, bǐjiào qīng, fāngbiàn dài.) — J'utilise parfois la tablette, elle est assez légère, facile à emporter.
\end{itemize}
\end{exemplebox}

\begin{exercice}[Entraînement oral individuel / Enregistreur vocal]
Enregistrez-vous (\textbf{1 minute max}) en répondant aux 3 questions du dialogue \textbf{pour vous-même} :
1. 你常常上网吗? (Fréquence précise)
2. 你用互联网做什么? (3 activités min., dont 1 nouvelle du vocabulaire)
3. 你用电脑还是手机? Pourquoi? (Outil + adjectif justificatif : 方便、快、大、轻、慢)
\emph{Écoutez votre enregistrement : vérifiez les tons \zh{shàngwǎng}, \zh{shǒujī}, \zh{hén fāngbiàn}, \zh{yìbān}.}
\end{exercice}

\begin{corrige}[Exercice : Entraînement oral]
\textbf{Attendus (exemples de réponses correctes) :}
1. \zh{对，我每天都上网。} / \zh{不常常，我一周两次。} / \zh{我每天晚上上网两个小时。}
2. \zh{我常常查资料、学习英语，有时候看电影、听音乐。} (Utilisation de \zh{查资料} chá zīliào, \zh{学习英语} xuéxí Yīngyǔ, \zh{看电影} kàn diànyǐng, \zh{听音乐} tīng yīnyuè).
3. \zh{我一般用手机，很方便。} / \zh{我一般用电脑，屏幕大，很快。} / \zh{我两个都用，手机方便，电脑大。} (Utilisation de \zh{两个都用} liǎng gè dōu yòng).
\textbf{Critères de réussite :} Phrases complètes, tons corrects sur les mots-clés (sandhi \zh{很} $\rightarrow$ \zh{hén}, \zh{一} $\rightarrow$ \zh{yì}), débit fluide sans longues hésitations.
\end{corrige}

\begin{savaistu}[Le clavier chinois : Pinyin vs Wubi]
Au Cameroun comme en France, on tape le chinois via \textbf{Pinyin} (entrées phonétiques) sur clavier AZERTY/QWERTY. En Chine continentale, beaucoup de professionnels (secrétaires, journalistes, joueurs) utilisent le \textbf{Wubi (五笔)} : une méthode par structure de traits (clés radicales), bien plus rapide une fois maîtrisée (pas de choix d'homophones), mais difficile à apprendre. Sur smartphone, le \textbf{Pinyin 9 touches (九键)} ou \textbf{26 touches (全拼)} domine, aidé par la prédiction intelligente (IA). Essayez d'installer un clavier chinois (Sogou, Microsoft Pinyin, Google Pinyin) sur votre téléphone pour vous exercer à taper les phrases du dialogue !
\end{savaistu}

\coursec{Formules fonctionnelles : présenter, opiner, relancer}

Dans la section précédente, nous avons appris à parler de nos habitudes sur internet à travers un dialogue modèle : dire ce que l'on fait en ligne, à quelle fréquence, avec quel appareil. Mais un vrai débat ne s'arrête pas aux habitudes : il faut savoir \cle{prendre la parole}, \cle{donner son avis}, \cle{réagir} aux idées des autres et \cle{relancer} la conversation. C'est exactement l'objet de cette section : construire votre « boîte à outils du débat » en chinois.

\courssub{Observer d'abord : une prise de parole complète}

Lisons d'abord une courte prise de parole d'un élève camerounais lors d'un débat en classe de chinois. Observez les formules soulignées par leur structure.

\begin{textebox}[Une prise de parole en classe]
大家好，今天我要谈谈互联网。我想说三点。\\
Dàjiā hǎo, jīntiān wǒ yào tán tan hùliánwǎng. Wǒ xiǎng shuō sān diǎn.\\
« Bonjour à tous, aujourd'hui je vais parler d'internet. Je veux dire trois choses. »\\[4pt]
第一，我觉得互联网很有用，因为我们可以很快地找到信息。\\
Dì-yī, wǒ juéde hùliánwǎng hěn yǒuyòng, yīnwèi wǒmen kěyǐ hěn kuài de zhǎodào xìnxī.\\
« Premièrement, je trouve qu'internet est très utile, parce qu'on peut trouver des informations très vite. »\\[4pt]
第二，我认为互联网让我们的生活更方便，因为我们可以网上购物、网上学习。\\
Dì-èr, wǒ rènwéi hùliánwǎng ràng wǒmen de shēnghuó gèng fāngbiàn, yīnwèi wǒmen kěyǐ wǎngshàng gòuwù, wǎngshàng xuéxí.\\
« Deuxièmement, je pense qu'internet rend notre vie plus pratique, parce qu'on peut faire des achats et étudier en ligne. »\\[4pt]
第三，对我来说，互联网也有问题：很多年轻人花太多时间玩手机。\\
Dì-sān, duì wǒ láishuō, hùliánwǎng yě yǒu wèntí: hěn duō niánqīngrén huā tài duō shíjiān wán shǒujī.\\
« Troisièmement, pour moi, internet a aussi des problèmes : beaucoup de jeunes passent trop de temps sur leur téléphone. »\\[4pt]
总之，互联网是好东西，但是我们要好好使用它。我的看法就是这样。你呢？\\
Zǒngzhī, hùliánwǎng shì hǎo dōngxi, dànshì wǒmen yào hǎohǎo shǐyòng tā. Wǒ de kànfǎ jiùshì zhèyàng. Nǐ ne?\\
« En conclusion, internet est une bonne chose, mais nous devons bien l'utiliser. Voilà mon point de vue. Et toi ? »
\end{textebox}

\textbf{Questions d'observation :}
\begin{itemize}
\item Quelles formules l'élève utilise-t-il pour commencer et pour annoncer son plan ?
\item Quelles expressions introduisent son avis personnel ? En trouvez-vous trois ?
\item Comment conclut-il, et comment passe-t-il la parole à un camarade ?
\end{itemize}

\courssub{Présenter un exposé : ouvrir et annoncer}

\begin{propriete}[Ouvrir un exposé en chinois]
Pour commencer une prise de parole, on suit un schéma fixe : \cle{salutation} + \cle{annonce du sujet} + \cle{annonce du plan}.
\begin{itemize}
\item Salutation : \zh{大家好！} (Dàjiā hǎo!) « Bonjour à tous ! » ; plus formel : \zh{老师好，同学们好！} (Lǎoshī hǎo, tóngxuémen hǎo!) « Bonjour professeur, bonjour camarades ! »
\item Annonce du sujet : \zh{今天我要谈谈……} (Jīntiān wǒ yào tán tan...) « Aujourd'hui je vais parler de... ». Le verbe \zh{谈} (tán, « parler de, discuter ») est souvent redoublé : \zh{谈谈} (tán tan, deuxième syllabe ton neutre), ce qui adoucit le ton.
\item Annonce du plan : \zh{我想说三点。} (Wǒ xiǎng shuō sān diǎn.) « Je veux dire trois choses. » \zh{点} (diǎn) signifie ici « point, aspect ».
\end{itemize}
\end{propriete}

\begin{exemplebox}[Annoncer le sujet : variations]
\zh{今天我要谈谈手机和互联网。}\\
\emph{Jīntiān wǒ yào tán tan shǒujī hé hùliánwǎng.} « Aujourd'hui je vais parler du téléphone et d'internet. »\\[3pt]
\zh{我想说两点：第一，好处；第二，问题。}\\
\emph{Wǒ xiǎng shuō liǎng diǎn: dì-yī, hǎochù; dì-èr, wèntí.} « Je veux dire deux choses : premièrement, les avantages ; deuxièmement, les problèmes. »
\end{exemplebox}

\begin{attention}[谈谈 et non 说说]
Pour « parler d'un sujet », on dit \zh{谈谈} + sujet. Ne dites pas \zh{我说互联网} pour « je parle d'internet » : \zh{说} (shuō) signifie « dire, parler » mais s'emploie avec un contenu de parole (\zh{说话} « parler », \zh{说汉语} « parler chinois »), pas avec un thème d'exposé. La bonne formule est \zh{谈谈互联网} ou \zh{谈一谈互联网}.
\end{attention}

\courssub{Donner son avis : les quatre formules d'opinion}

\begin{propriete}[Les verbes et tournures d'opinion]
Quatre formules permettent d'exprimer son avis, de la plus courante à la plus formelle :
\begin{enumerate}
\item \zh{我觉得} + phrase (wǒ juéde) : « je trouve que... » — la plus fréquente à l'oral.
\item \zh{我认为} + phrase (wǒ rènwéi) : « je pense que... » — un peu plus soutenu.
\item \zh{对我来说}，+ phrase (duì wǒ láishuō) : « pour moi, en ce qui me concerne... » — se place en tête de phrase, suivi d'une virgule.
\item \zh{我的看法是} + phrase (wǒ de kànfǎ shì) : « mon point de vue est... » — formel, idéal pour conclure.
\end{enumerate}
Contrairement au français « je pense \emph{que} », le chinois n'utilise aucune conjonction : on enchaîne directement \zh{我觉得} + proposition complète.
\end{propriete}

\begin{exemplebox}[Donner son avis]
\zh{我觉得互联网很有用。} \emph{Wǒ juéde hùliánwǎng hěn yǒuyòng.} « Je trouve qu'internet est très utile. »\\[3pt]
\zh{我认为学习汉语很有意思。} \emph{Wǒ rènwéi xuéxí Hànyǔ hěn yǒu yìsi.} « Je pense qu'apprendre le chinois est très intéressant. »\\[3pt]
\zh{对我来说，手机比电脑方便。} \emph{Duì wǒ láishuō, shǒujī bǐ diànnǎo fāngbiàn.} « Pour moi, le téléphone est plus pratique que l'ordinateur. » (On retrouve la structure de comparaison avec \zh{比} déjà étudiée : A + 比 + B + adjectif.)\\[3pt]
\zh{我的看法是：我们应该少玩手机。} \emph{Wǒ de kànfǎ shì: wǒmen yīnggāi shǎo wán shǒujī.} « Mon point de vue est : nous devrions moins jouer avec nos téléphones. »
\end{exemplebox}

\courssub{Structurer son propos : 第一、第二、第三 et 总之}

\begin{propriete}[Numéroter ses arguments]
Pour ordonner ses arguments, on utilise \zh{第一} (dì-yī, « premièrement »), \zh{第二} (dì-èr, « deuxièmement »), \zh{第三} (dì-sān, « troisièmement »), chacun suivi d'une virgule chinoise \zh{，}. Le préfixe \zh{第} (dì) forme les nombres ordinaux. Pour conclure : \zh{总之} (zǒngzhī, « en conclusion, en résumé »), placé en tête de la dernière phrase.
\end{propriete}

\begin{exemplebox}[Une conclusion avec 总之]
\zh{总之，我觉得互联网很有用，但是不能太依赖它。}\\
\emph{Zǒngzhī, wǒ juéde hùliánwǎng hěn yǒuyòng, dànshì bù néng tài yīlài tā.}\\
« En conclusion, je trouve qu'internet est très utile, mais on ne peut pas trop en dépendre. »
\end{exemplebox}

\courssub{Argumenter : la structure 因为……所以……}

\begin{propriete}[因为……所以…… : cause et conséquence]
\zh{因为} (yīnwèi, « parce que ») introduit la cause ; \zh{所以} (suǒyǐ, « donc, c'est pourquoi ») introduit la conséquence. Schéma : \textbf{因为 + cause，所以 + conséquence}.\\
Grande différence avec le français : en français, on ne dit jamais « \emph{parce que}..., \emph{donc}... » dans la même phrase ; en chinois, les deux mots peuvent — et souvent doivent — coexister. On peut aussi n'en utiliser qu'un seul, mais jamais traduire mot à mot le français.
\end{propriete}

\begin{exemplebox}[因为 et 所以 ensemble]
\zh{因为互联网很方便，所以大家都喜欢。}\\
\emph{Yīnwèi hùliánwǎng hěn fāngbiàn, suǒyǐ dàjiā dōu xǐhuan.}\\
« Parce qu'internet est très pratique, tout le monde l'aime. »\\[4pt]
\zh{因为网上有很多信息，所以学习更容易了。}\\
\emph{Yīnwèi wǎngshàng yǒu hěn duō xìnxī, suǒyǐ xuéxí gèng róngyì le.}\\
« Parce qu'il y a beaucoup d'informations en ligne, apprendre est devenu plus facile. »\\[4pt]
\zh{我觉得互联网很有用，因为我们可以很快地找到信息。}\\
\emph{Wǒ juéde hùliánwǎng hěn yǒuyòng, yīnwèi wǒmen kěyǐ hěn kuài de zhǎodào xìnxī.}\\
« Je trouve qu'internet est très utile, parce qu'on peut trouver des informations très vite. » (Ici, \zh{因为} seul, placé après l'opinion, comme « parce que » en français.)
\end{exemplebox}

\begin{attention}[Le piège de la traduction mot à mot]
Les francophones traduisent souvent « parce que... donc... » en gardant un seul mot, ou pire, ils traduisent « parce que » par \zh{所以}. Retenez : \zh{因为} = cause (« parce que »), \zh{所以} = résultat (« donc »). Phrase correcte : \zh{因为我很忙，所以不能上网。} (Yīnwèi wǒ hěn máng, suǒyǐ bù néng shàngwǎng.) « Parce que je suis occupé, je ne peux pas aller sur internet. »
\end{attention}

\courssub{Réagir : exprimer l'accord et le désaccord poli}

\begin{propriete}[L'accord]
\begin{itemize}
\item \zh{我同意你的看法。} (Wǒ tóngyì nǐ de kànfǎ.) « Je suis d'accord avec ton point de vue. »
\item \zh{你说得对。} (Nǐ shuō de duì.) « Tu as raison. » (littéralement « tu dis juste », avec le complément de manière \zh{得}).
\item \zh{我也这么认为。} (Wǒ yě zhème rènwéi.) « Je le pense aussi. »
\end{itemize}
\end{propriete}

\begin{propriete}[Le désaccord poli]
En chinois comme en français, on adoucit le désaccord :
\begin{itemize}
\item \zh{对不起，我不同意。} (Duìbuqǐ, wǒ bù tóngyì.) « Pardon, je ne suis pas d'accord. »
\item \zh{我不这么认为。} (Wǒ bù zhème rènwéi.) « Je ne le pense pas. » (littéralement « je ne pense pas ainsi » ; \zh{这么} = « ainsi »).
\item \zh{你说得有道理，但是……} (Nǐ shuō de yǒu dàolǐ, dànshì...) « Ce que tu dis est raisonnable, mais... » — formule de concession très appréciée dans un débat.
\end{itemize}
\end{propriete}

\begin{attention}[Ne pas dire *我同意你 seul]
Le verbe \zh{同意} (tóngyì, « être d'accord avec ») ne se construit pas directement avec une personne seule. On ne dit pas \zh{我同意你}. Les formes correctes sont : \zh{我同意你的看法} (« je suis d'accord avec ton avis »), \zh{我同意你说的话} (wǒ tóngyì nǐ shuō de huà, « je suis d'accord avec ce que tu dis »), ou simplement \zh{我同意} (« je suis d'accord »).
\end{attention}

\courssub{Relancer la parole}

\begin{propriete}[Passer la parole à un camarade]
\begin{itemize}
\item \zh{你呢？} (Nǐ ne?) « Et toi ? » — la particule \zh{呢} reprend le thème de la question précédente ; c'est la relance la plus simple et la plus naturelle.
\item \zh{你觉得怎么样？} (Nǐ juéde zěnmeyàng?) « Qu'en penses-tu ? » (littéralement « tu trouves comment ? »).
\item \zh{你有什么看法？} (Nǐ yǒu shénme kànfǎ?) « Quel est ton avis ? »
\item \zh{你同意吗？} (Nǐ tóngyì ma?) « Es-tu d'accord ? »
\end{itemize}
\end{propriete}

\courssub{Tableau récapitulatif de la boîte à outils}

\begin{center}
\begin{tabularx}{\linewidth}{|X|X|X|}
\hline
\rowcolor{popblueL} Fonction & Formule chinoise + pinyin & Traduction \\
\hline
Saluer & 大家好！Dàjiā hǎo! & Bonjour à tous ! \\
\hline
Annoncer le sujet & 今天我要谈谈…… Jīntiān wǒ yào tán tan... & Aujourd'hui je vais parler de... \\
\hline
Annoncer le plan & 我想说三点。Wǒ xiǎng shuō sān diǎn. & Je veux dire trois choses. \\
\hline
Donner son avis & 我觉得 / 我认为 / 对我来说 / 我的看法是 & Je trouve que / je pense que / pour moi / mon avis est \\
\hline
Structurer & 第一……第二……第三……总之 & Premièrement... deuxièmement... troisièmement... en conclusion \\
\hline
Argumenter & 因为……所以…… yīnwèi... suǒyǐ... & Parce que... donc... \\
\hline
Accord & 我同意你的看法。/ 你说得对。 & Je suis d'accord avec toi. / Tu as raison. \\
\hline
Désaccord poli & 对不起，我不同意。/ 我不这么认为。 & Pardon, je ne suis pas d'accord. / Je ne le pense pas. \\
\hline
Relancer & 你呢？/ 你觉得怎么样？/ 你有什么看法？ & Et toi ? / Qu'en penses-tu ? / Quel est ton avis ? \\
\hline
\end{tabularx}
\end{center}

\begin{popfigure}
\begin{tikzpicture}[
  node distance=0.45cm,
  box/.style={draw=popblue, fill=popblueL, rounded corners=3pt, align=center, text width=4.6cm, inner sep=5pt, font=\small},
  arr/.style={-{Stealth[length=2.5mm]}, thick, draw=poporange}
]
\node[box, align=center] (a){\textbf{Présenter}\\大家好，今天我要谈谈……};
\node[box, below=of a, align=center] (b){\textbf{Opiner}\\我觉得…… / 我认为……};
\node[box, below=of b, align=center] (c){\textbf{Argumenter}\\因为……所以……};
\node[box, below=of c, align=center] (d){\textbf{Réagir}\\我同意你的看法。/ 我不这么认为。};
\node[box, below=of d, align=center] (e){\textbf{Relancer}\\你呢？/ 你觉得怎么样？};
\draw[arr] (a) -- (b);
\draw[arr] (b) -- (c);
\draw[arr] (c) -- (d);
\draw[arr] (d) -- (e);
\draw[arr, dashed, draw=poppurple] (e.east) .. controls +(1.6,0) and +(1.6,0) .. (b.east);
\end{tikzpicture}
\legende{Organigramme : la boîte à outils du débat — après la relance, l'échange reprend (flèche en pointillés).}
\end{popfigure}

\courssub{Méthode et entraînement}

\begin{methode}[Réussir un mini-exposé en cinq étapes]
\begin{enumerate}
\item \textbf{Saluer} : \zh{大家好！}
\item \textbf{Annoncer le sujet} : \zh{今天我要谈谈互联网。}
\item \textbf{Annoncer le plan} : \zh{我想说三点。}
\item \textbf{Développer 2 ou 3 arguments} avec \zh{第一……第二……第三……}, chacun appuyé par une cause en \zh{因为}.
\item \textbf{Conclure} avec \zh{总之} + votre avis, puis \textbf{relancer} : \zh{你呢？}
\end{enumerate}
\end{methode}

\begin{exoresolu}[Compléter une prise de parole]
Énoncé : complétez avec la formule fonctionnelle qui convient.\\
1. \zh{……，今天我要谈谈互联网。} (saluer)\\
2. \zh{……互联网很有用，因为我们可以网上购物。} (donner son avis)\\
3. \zh{……手机很方便，……大家都喜欢用手机上网。} (cause + conséquence)\\
4. A : \zh{我觉得互联网对学习很有帮助。} B : \zh{……你的看法。} (exprimer l'accord)\\
5. A a fini son exposé et passe la parole : \zh{……？} (relancer)
\tcblower
Solution détaillée :\\
1. \zh{大家好} (Dàjiā hǎo) : c'est la salutation d'ouverture obligatoire.\\
2. \zh{我觉得} (wǒ juéde) : formule d'opinion la plus naturelle à l'oral ; \zh{我认为} était aussi possible.\\
3. \zh{因为……所以……} : \zh{因为手机很方便，所以大家都喜欢用手机上网。} (Yīnwèi shǒujī hěn fāngbiàn, suǒyǐ dàjiā dōu xǐhuan yòng shǒujī shàngwǎng.) En chinois, les deux mots coexistent, contrairement au français.\\
4. \zh{我同意} : \zh{我同意你的看法。} (Wǒ tóngyì nǐ de kànfǎ.) Attention : jamais \zh{我同意你} seul.\\
5. \zh{你呢？} (Nǐ ne?) ou \zh{你有什么看法？} (Nǐ yǒu shénme kànfǎ?) : les deux relances sont correctes.
\end{exoresolu}

\begin{exercice}[À vous de parler]
Préparez oralement un mini-exposé de cinq phrases sur le thème « 手机和互联网 » (le téléphone et internet) en suivant la méthode : salutation, annonce du sujet, deux arguments avec \zh{因为}, conclusion avec \zh{总之}, relance avec \zh{你呢？}. Puis, par groupes de deux, réagissez à l'exposé de votre camarade avec une formule d'accord ou de désaccord poli.
\end{exercice}

Ces formules acquises, nous pourrons, dans la section suivante, les entendre en action dans un deuxième dialogue modèle consacré au débat sur les avantages et les inconvénients d'internet.

\coursec{Dialogue modèle 2 : débattre des avantages et inconvénients}

Dans la section précédente, tu as appris à \cle{présenter} un sujet (\zh{我谈谈...} \emph{wǒ tántan...} « je vais parler de... »), à \cle{donner ton avis} (\zh{我觉得...} \emph{wǒ juéde...} « je trouve que... ») et à \cle{relancer} ton camarade (\zh{你呢？你觉得怎么样？} \emph{Nǐ ne? Nǐ juéde zěnmeyàng?} « Et toi ? Qu'en penses-tu ? »). Ces formules sont les outils du débat. Nous allons maintenant les mettre en action dans un dialogue complet, où deux élèves discutent des \cle{avantages} et des \cle{inconvénients} d'internet. Tu vas d'abord écouter et observer, puis nous analyserons chaque structure avant de te faire débattre à ton tour.

\courssub{Le dialogue-débat : observation}

\begin{textebox}[Dialogue 2 : 互联网好不好？]
A：我觉得互联网很好，我们可以学习很多东西。\\
\emph{Wǒ juéde hùliánwǎng hěn hǎo, wǒmen kěyǐ xuéxí hěn duō dōngxi.}\\
« Je trouve qu'internet est très bien, on peut apprendre beaucoup de choses. »

B：你说得对，可是我觉得互联网也有坏处。\\
\emph{Nǐ shuō de duì, kěshì wǒ juéde hùliánwǎng yě yǒu huàichù.}\\
« Tu as raison, mais je trouve qu'internet a aussi des inconvénients. »

A：什么坏处？\\
\emph{Shénme huàichù?}\\
« Quels inconvénients ? »

B：有些学生玩游戏太多，不学习。\\
\emph{Yǒuxiē xuésheng wán yóuxì tài duō, bù xuéxí.}\\
« Certains élèves jouent trop aux jeux et n'étudient pas. »

A：我同意，所以我们应该少玩游戏，多学习。\\
\emph{Wǒ tóngyì, suǒyǐ wǒmen yīnggāi shǎo wán yóuxì, duō xuéxí.}\\
« Je suis d'accord, donc nous devons moins jouer et plus étudier. »
\end{textebox}

\begin{methode}[Lecture dramatisée à deux voix]
\begin{enumerate}
\item \textbf{Première écoute} : le professeur lit le dialogue, les élèves écoutent sans lire, livre fermé. Question globale : les deux élèves sont-ils d'accord ou pas d'accord ?
\item \textbf{Lecture à voix haute} : deux élèves lisent les rôles A et B en suivant le pinyin, en exagérant les tons.
\item \textbf{Lecture dramatisée} : A doit avoir l'air enthousiaste, B sceptique mais poli. Le ton monte sur la question \zh{什么坏处？} et la phrase finale est prononcée avec conviction.
\item \textbf{Échange des rôles}, puis lecture sans pinyin, uniquement avec les caractères.
\end{enumerate}
\end{methode}

\courssub{Le vocabulaire du débat}

\begin{center}
\begin{tabularx}{\linewidth}{|X|X|X|X|}
\hline
\rowcolor{popblueL} Caractères & Pinyin & Nature & Traduction \\
\hline
好处 & hǎochù & nom & avantage, bon côté \\
\hline
坏处 & huàichù & nom & inconvénient, mauvais côté \\
\hline
有用 & yǒuyòng & adjectif & utile \\
\hline
危险 & wēixiǎn & adjectif & dangereux \\
\hline
应该 & yīnggāi & verbe modal & devoir (il faut) \\
\hline
同意 & tóngyì & verbe & être d'accord \\
\hline
有些 & yǒuxiē & déterminant & certains, quelques \\
\hline
浪费 & làngfèi & verbe & gaspiller, perdre \\
\hline
信息 & xìnxī & nom & information \\
\hline
联系 & liánxì & verbe & contacter, être en contact \\
\hline
\end{tabularx}
\end{center}

Observe la formation des deux mots clés : \zh{好处} et \zh{坏处} sont construits sur le même modèle : \cle{adjectif} (\zh{好} « bon » / \zh{坏} « mauvais ») + \zh{处} « côté, aspect ». Le chinois oppose donc naturellement « le bon côté » et « le mauvais côté » d'une chose, exactement comme le français « le pour et le contre ». Retiens aussi la paire d'adverbes \zh{多} « plus, davantage » et \zh{少} « moins », placés \textbf{avant le verbe} : \zh{多学习} \emph{duō xuéxí} « étudier plus », \zh{少玩游戏} \emph{shǎo wán yóuxì} « jouer moins ».

\begin{definition}[\zh{有些} : « certains », le sujet indéfini du débat]
\zh{有些} (\emph{yǒuxiē}) signifie « certains, quelques ». Il se place devant un nom pour désigner une partie d'un groupe, sans le nommer précisément. C'est l'outil idéal pour argumenter sans généraliser :\\
\zh{有些学生玩游戏太多。} \emph{Yǒuxiē xuésheng wán yóuxì tài duō.} « Certains élèves jouent trop. »\\
\zh{有些人认为互联网很危险。} \emph{Yǒuxiē rén rènwéi hùliánwǎng hěn wēixiǎn.} « Certaines personnes pensent qu'internet est dangereux. »\\
Structure : \textbf{有些 + nom + verbe/adjectif}. Ne confonds pas avec \zh{一些} (\emph{yìxiē}) « quelques, un peu de », qui sert à quantifier un complément (\zh{买一些水果} « acheter quelques fruits ») : pour ouvrir un débat avec un sujet indéfini, c'est \zh{有些} qu'il faut.
\end{definition}

\begin{propriete}[\zh{太...了} : exprimer l'excès]
La structure \textbf{太 + adjectif/verbe + 了} exprime l'excès (« trop ») ou l'intensité forte.\\
\zh{太多了} \emph{tài duō le} « c'est trop » ; \zh{太贵了} \emph{tài guì le} « c'est trop cher » ; \zh{太好了} \emph{tài hǎo le} « c'est formidable ! »\\
Avec un verbe d'action : \zh{玩游戏太多了} \emph{wán yóuxì tài duō le} « jouer trop aux jeux » ; \zh{玩手机的时间太长了} \emph{wán shǒujī de shíjiān tài cháng le} « passer trop de temps sur son téléphone ».\\
Attention au sens : selon l'adjectif, \zh{太...了} est négatif (trop cher, trop dangereux) ou admiratif (\zh{太好了} « génial »). Le contexte décide. Dans le dialogue, \zh{太多} apparaît sans \zh{了} final (\zh{玩游戏太多}) : à l'intérieur d'une phrase qui continue, le \zh{了} peut tomber ; en fin de phrase, on le garde.
\end{propriete}

\courssub{Contre-argumenter : \zh{可是} et \zh{但是}}

\begin{propriete}[Le connecteur d'opposition]
Pour opposer un argument, le chinois utilise \zh{可是} (\emph{kěshì}) ou \zh{但是} (\emph{dànshì}), tous deux « mais ». \zh{可是} est plus oral, \zh{但是} un peu plus soutenu. Ils se placent \textbf{au début de la seconde proposition}, comme en français :\\
\zh{互联网很有用，可是也有坏处。} \emph{Hùliánwǎng hěn yǒuyòng, kěshì yě yǒu huàichù.} « Internet est très utile, mais il a aussi des inconvénients. »\\
\zh{你说得对，但是我不同意。} \emph{Nǐ shuō de duì, dànshì wǒ bù tóngyì.} « Tu as raison, mais je ne suis pas d'accord. »\\
Formule de politesse du débat : on \textbf{concède d'abord} (\zh{你说得对} « tu as raison »), puis on oppose avec \zh{可是}. C'est plus élégant qu'un refus direct. Remarque au passage la forme \zh{你说得对} : le mot-outil \zh{得} (\emph{de}) introduit ici le complément de manière « bien, correctement » après le verbe \zh{说} « dire » ; littéralement : « tu dis de façon juste ».
\end{propriete}

\begin{propriete}[\zh{所以} : conclure le débat]
\zh{所以} (\emph{suǒyǐ}) « donc, c'est pourquoi » introduit la conclusion. Il se place \textbf{devant le sujet} de la proposition de conséquence :\\
\zh{互联网有好处，也有坏处，所以我们应该少玩游戏，多学习。}\\
\emph{Hùliánwǎng yǒu hǎochù, yě yǒu huàichù, suǒyǐ wǒmen yīnggāi shǎo wán yóuxì, duō xuéxí.}\\
« Internet a des avantages et aussi des inconvénients, donc nous devons moins jouer et plus étudier. »\\
Structure complète du débat équilibré : \textbf{avantage + 可是 + inconvénient + 所以 + conclusion}.
\end{propriete}

\begin{center}
\begin{tabularx}{\linewidth}{|X|X|X|}
\hline
\rowcolor{popblueL} Fonction & Formule chinoise + pinyin & Traduction \\
\hline
Concéder & 你说得对。Nǐ shuō de duì. & Tu as raison. \\
\hline
Opposer & 可是... / 但是... kěshì / dànshì & mais... \\
\hline
Nuancer & 有些... yǒuxiē... & certains... \\
\hline
Conclure & 所以我们应该... suǒyǐ wǒmen yīnggāi... & donc nous devons... \\
\hline
\end{tabularx}
\end{center}

\courssub{Le pour et le contre d'internet}

\begin{popfigure}
\begin{tikzpicture}
\node[draw=popgreen, fill=popgreenL, rounded corners, thick, minimum width=5.2cm, align=center, font=\bfseries] (hao) at (-3.2,0) {好处 hǎochù\\ avantages};
\node[draw=poppink, fill=poppinkL, rounded corners, thick, minimum width=5.2cm, align=center, font=\bfseries] (huai) at (3.2,0) {坏处 huàichù\\ inconvénients};
\node[draw=popgreen, fill=white, rounded corners, align=center] at (-3.2,-1.3) {学习很多东西\\ apprendre beaucoup};
\node[draw=popgreen, fill=white, rounded corners, align=center] at (-3.2,-2.5) {查信息\\ chercher des informations};
\node[draw=popgreen, fill=white, rounded corners, align=center] at (-3.2,-3.7) {跟朋友联系\\ contacter ses amis};
\node[draw=poppink, fill=white, rounded corners, align=center] at (3.2,-1.3) {玩游戏太多\\ jouer trop};
\node[draw=poppink, fill=white, rounded corners, align=center] at (3.2,-2.5) {浪费时间\\ perdre du temps};
\node[draw=poppink, fill=white, rounded corners, align=center] at (3.2,-3.7) {坏信息\\ mauvaises informations};
\draw[-{Stealth}, popgreen, thick] (hao) -- (-3.2,-0.8);
\draw[-{Stealth}, poppink, thick] (huai) -- (3.2,-0.8);
\node[draw=popblue, fill=popblueL, rounded corners, align=center] at (0,-4.7) {所以我们应该少玩游戏，多学习。\\ Donc nous devons moins jouer et plus étudier.};
\end{tikzpicture}
\legende{Carte mentale du débat : le pour et le contre d'internet, avec la conclusion équilibrée.}
\end{popfigure}

\courssub{Travail phonétique : les tons du débat}

Quatre mots du dialogue demandent une attention particulière :

\begin{itemize}
\item \zh{好处} \emph{hǎochù} : 3\up{e} ton + 4\up{e} ton. Le 3\up{e} ton ne se prononce entièrement (descente puis remontée) que lorsqu'il est isolé ou en fin de phrase. Devant un autre ton, il devient un \textbf{demi-3\up{e} ton} : la voix descend bas et ne remonte pas. Prononce donc un demi-3\up{e} ton suivi d'un 4\up{e} ton sec : \emph{hǎo-chù}.
\item \zh{应该} \emph{yīnggāi} : 1\up{er} ton + 1\up{er} ton. Deux tons hauts et plats, sans descendre : garde la voix haute et régulière sur les deux syllabes.
\item \zh{危险} \emph{wēixiǎn} : 1\up{er} ton + 3\up{e} ton. Haut et plat, puis descente-remontée complète, car le 3\up{e} ton est ici en fin de mot accentué.
\item \zh{可是} \emph{kěshì} : 3\up{e} ton + 4\up{e} ton. Même profil que \zh{好处} : demi-3\up{e} ton puis chute nette. C'est le mot qui porte ton opposition : appuie sur le 4\up{e} ton de \zh{是}.
\end{itemize}

\begin{attention}[\zh{应该} : un piège de prononciation et de sens]
Le caractère \zh{应} a deux lectures : \emph{yīng} (1\up{er} ton) dans \zh{应该} « devoir », et \emph{yìng} (4\up{e} ton) dans \zh{答应} \emph{dāying} « répondre, promettre ». Dans le débat, dis toujours \emph{yīnggāi} avec le 1\up{er} ton.\\
Autre piège : ne confonds pas \zh{应该} « devoir » (conseil, obligation morale) avec \zh{可以} « pouvoir » (permission, possibilité) :\\
\zh{我们应该多学习。} \emph{Wǒmen yīnggāi duō xuéxí.} « Nous devons étudier davantage. » (c'est un devoir)\\
\zh{我们可以用互联网学习。} \emph{Wǒmen kěyǐ yòng hùliánwǎng xuéxí.} « Nous pouvons utiliser internet pour étudier. » (c'est une possibilité)
\end{attention}

\begin{exoresolu}[Construis une contre-argumentation]
Énoncé : ton camarade dit \zh{互联网很有用。} « Internet est très utile. » Réponds en trois étapes : concède, oppose avec \zh{可是} et un argument en \zh{有些}, conclus avec \zh{应该}.
\tcblower
Solution détaillée :\\
1. Concéder : \zh{你说得对，互联网很有用。} \emph{Nǐ shuō de duì, hùliánwǎng hěn yǒuyòng.} « Tu as raison, internet est très utile. »\\
2. Opposer : \zh{可是有些学生玩游戏太多，浪费时间。} \emph{Kěshì yǒuxiē xuésheng wán yóuxì tài duō, làngfèi shíjiān.} « Mais certains élèves jouent trop et perdent du temps. »\\
3. Conclure : \zh{所以我们应该少玩游戏，多学习。} \emph{Suǒyǐ wǒmen yīnggāi shǎo wán yóuxì, duō xuéxí.} « Donc nous devons moins jouer et plus étudier. »\\
Réponse complète : \zh{你说得对，可是有些学生玩游戏太多，浪费时间，所以我们应该少玩游戏，多学习。}
\end{exoresolu}

\begin{savaistu}[Internet des jeunes : la Chine et le Cameroun]
En Chine, le temps de jeu en ligne des mineurs est limité par la loi : les plateformes de jeux doivent vérifier l'âge des joueurs et restreindre les horaires de connexion des enfants et adolescents. L'objectif officiel est de protéger les études et la santé des jeunes. Au Cameroun, le débat porte plutôt sur l'usage du téléphone portable en classe : beaucoup d'établissements de Yaoundé ou de Douala l'interdisent pendant les cours, car il distrait les élèves. Dans les deux pays, la conclusion ressemble à celle de notre dialogue : \zh{少玩游戏，多学习！} \emph{Shǎo wán yóuxì, duō xuéxí!} « Moins de jeux, plus d'études ! »
\end{savaistu}

\begin{exercice}[À ton tour de débattre]
Par groupes de deux, construisez un mini-débat de six répliques sur le modèle étudié : A défend internet (deux arguments de la colonne \zh{好处}), B oppose deux arguments de la colonne \zh{坏处} avec \zh{可是} et \zh{有些}, et A conclut avec \zh{所以我们应该...}. Soignez les tons de \zh{好处}, \zh{可是}, \zh{应该} et \zh{危险}. Les meilleurs dialogues seront joués devant la classe.
\end{exercice}

\coursec{Travail phonétique : les tons dans le débat}

Dans la section précédente, nous avons appris à débattre des avantages et des inconvénients d'internet avec des phrases comme \zh{我觉得上网很方便。} (\emph{Wǒ juéde shàngwǎng hěn fāngbiàn.} « Je trouve qu'aller sur internet est très pratique. »). Mais une phrase juste sur le papier peut devenir incompréhensible à l'oral si les tons sont faux. En chinois, le ton fait partie du mot, comme une lettre en français : on ne peut pas le négliger. Cette section est donc un entraînement phonétique complet pour que votre débat sonne vraiment chinois.

\courssub{Révision rapide des quatre tons et du ton neutre}

\begin{definition}[Les tons du chinois mandarin]
Le chinois mandarin possède \cle{quatre tons} et un \cle{ton neutre}. Le ton est la mélodie de la syllabe : à consonnes et voyelles identiques, deux tons différents donnent deux mots différents.
\end{definition}

Observons d'abord les mots de notre leçon sur internet :

\begin{center}
\begin{tabularx}{\linewidth}{|X|X|X|X|}
\hline
\rowcolor{popblueL} Ton & Mot de la leçon & Pinyin & Traduction \\
\hline
1er ton (haut et plat) & 网 & wǎng & internet, réseau \\
\hline
2e ton (montant) & 聊 & liáo & bavarder \\
\hline
3e ton (bas, descendant-montant) & 好 & hǎo & bon \\
\hline
4e ton (descendant brusque) & 用 & yòng & utiliser \\
\hline
Ton neutre (léger, court) & 吗 & ma & particule de question \\
\hline
\end{tabularx}
\end{center}

\begin{popfigure}
\begin{tikzpicture}[scale=1.05]
% cadre commun : 5 niveaux de hauteur
\foreach \x/\titre/\mot/\py in {0/{1er ton}/网/wǎng, 3.6/{2e ton}/聊/liáo, 7.2/{3e ton}/好/hǎo, 10.8/{4e ton}/用/yòng}{
  \draw[popline, dashed] (\x,0) -- (\x+2.6,0);
  \draw[popline, dashed] (\x,0.7) -- (\x+2.6,0.7);
  \draw[popline, dashed] (\x,1.4) -- (\x+2.6,1.4);
  \draw[popline, dashed] (\x,2.1) -- (\x+2.6,2.1);
  \draw[popline, dashed] (\x,2.8) -- (\x+2.6,2.8);
  \node[align=center] at (\x+1.3,3.4) {\textbf{\titre}};
  \node[align=center, popdark] at (\x+1.3,-0.6) {\zh{\mot} \ \emph{\py}};
}
% 1er ton : ligne haute plate
\draw[line width=2.2pt, popblue] (0,2.8) -- (2.6,2.8);
% 2e ton : montante
\draw[line width=2.2pt, poporange] (3.6,1.4) -- (6.2,2.8);
% 3e ton : descendante-montante
\draw[line width=2.2pt, popgreen] (7.2,2.1) .. controls (7.9,0.4) and (8.6,0.4) .. (9.8,1.8);
% 4e ton : descendante brusque
\draw[line width=2.2pt, poppink] (10.8,2.8) -- (13.4,0.2);
\end{tikzpicture}
\legende{Les contours des quatre tons : hauteur de la voix (grave en bas, aiguë en haut) en fonction du temps.}
\end{popfigure}

\begin{aretenir}[Comment mémoriser chaque ton]
\begin{itemize}
\item \cle{1er ton} : voix haute et plate, comme quand le médecin vous demande de dire « aaaa ».
\item \cle{2e ton} : voix qui monte, comme quand on demande « hein ? » en français.
\item \cle{3e ton} : voix qui descend puis remonte ; dans la parole rapide, on ne prononce souvent que la partie basse (on parle alors de \cle{demi-3e ton}).
\item \cle{4e ton} : voix qui tombe brusquement, bref et sec.
\item \cle{Ton neutre} : syllabe courte et légère, sans mélodie propre (\zh{吗} ma, \zh{的} de).
\end{itemize}
\end{aretenir}

\courssub{Paires minimales : le ton et le contexte décident du sens}

Deux mots peuvent avoir le même son et le même ton (homophones), ou ne différer que par le ton. C'est le contexte — et votre prononciation — qui évite le malentendu.

\begin{center}
\begin{tabularx}{\linewidth}{|X|X|X|X|}
\hline
\rowcolor{popblueL} Paire & Pinyin & Différence & Traduction \\
\hline
查 / 茶 & chá / chá & même ton, même son : le contexte décide & consulter, chercher / thé \\
\hline
发 / 罚 & fā / fá & 1er ton contre 2e ton & envoyer / punir \\
\hline
问 / 吻 & wèn / wěn & 4e ton contre 3e ton & demander / embrasser \\
\hline
买 / 卖 & mǎi / mài & 3e ton contre 4e ton & acheter / vendre \\
\hline
\end{tabularx}
\end{center}

\begin{exemplebox}[Les paires minimales en contexte]
\zh{我想在网上查资料。} \emph{Wǒ xiǎng zài wǎngshang chá zīliào.} « Je veux chercher des informations sur internet. »\\
\zh{爷爷喜欢喝茶。} \emph{Yéye xǐhuan hē chá.} « Grand-père aime boire du thé. »\\
\zh{我发了一个邮件。} \emph{Wǒ fā le yí ge yóujiàn.} « J'ai envoyé un courriel. »\\
\zh{老师罚他写作业。} \emph{Lǎoshī fá tā xiě zuòyè.} « Le professeur le punit en lui faisant écrire ses devoirs. »
\end{exemplebox}

\begin{savaistu}[Un mauvais ton peut tout changer]
Un élève qui veut dire \zh{我想问你} (\emph{Wǒ xiǎng wèn nǐ}, « je veux te demander ») mais prononce un 3e ton au lieu du 4e dit \zh{我想吻你} (\emph{Wǒ xiǎng wěn nǐ}, « je veux t'embrasser ») ! D'où l'importance capitale des tons à l'oral, surtout dans un débat où l'on pose beaucoup de questions.
\end{savaistu}

\courssub{Le sandhi du troisième ton}

\begin{propriete}[Sandhi du 3e ton]
Quand deux syllabes au \cle{3e ton} se suivent, la première se prononce comme un \cle{2e ton}. C'est une règle obligatoire de prononciation, mais l'écriture pinyin garde le ton d'origine.
\begin{itemize}
\item \zh{很好} s'écrit hěn hǎo mais se prononce \emph{hén hǎo} « très bien ».
\item \zh{你好} s'écrit nǐ hǎo mais se prononce \emph{ní hǎo} « bonjour ».
\item \zh{我可以} s'écrit wǒ kěyǐ mais se prononce \emph{wó kěyǐ} « je peux ».
\end{itemize}
\cle{Exception} : devant un autre ton (1er, 2e, 4e ou neutre), il n'y a pas de sandhi ; le 3e ton garde sa forme, souvent réduite à sa partie basse (demi-3e ton) : \zh{好处} hǎochù « avantage » (3 + 4), \zh{网上} wǎngshang « sur internet » (3 + neutre).
\end{propriete}

\begin{exemplebox}[Le 3e ton dans la parole rapide]
\zh{我觉得上网有好处。} \emph{Wǒ juéde shàngwǎng yǒu hǎochù.} « Je trouve qu'internet a des avantages. »\\
Dans la parole rapide, le 3e ton de \zh{我} wǒ devant \zh{觉得} juéde (2e ton) ne se prononce pas en entier : on n'entend que sa partie basse (demi-3e ton), sans la remontée finale. C'est normal : ne cherchez pas à prononcer la descente-montée complète à l'intérieur d'une phrase rapide. Le sandhi complet (3e → 2e ton) n'a lieu que devant un autre 3e ton.
\end{exemplebox}

\begin{propriete}[Changement de ton de 不]
\zh{不} « ne... pas » est au 4e ton (bù), mais il devient \cle{2e ton} (bú) devant une autre syllabe au 4e ton.
\begin{itemize}
\item \zh{不是} bù shì → se prononce \emph{bú shì} « ce n'est pas ».
\item \zh{不用} bù yòng → se prononce \emph{bú yòng} « pas besoin ».
\item \zh{不看} bù kàn → se prononce \emph{bú kàn} « ne pas regarder ».
\item Devant les autres tons, il reste au 4e ton : \zh{不说} bù shuō « ne pas dire », \zh{不来} bù lái « ne pas venir », \zh{不好} bù hǎo « pas bon », \zh{不同意} bù tóngyì « ne pas être d'accord » (\zh{同} tóng est au 2e ton, donc pas de changement).
\end{itemize}
\end{propriete}

\courssub{L'intonation de la question avec 吗}

\begin{propriete}[Question avec 吗 : l'intonation monte, les tons ne changent pas]
Pour transformer une phrase affirmative en question, on ajoute \zh{吗} ma à la fin. La voix monte légèrement sur la fin de la phrase, mais \cle{les tons des mots ne changent pas} : seule la mélodie globale de la phrase s'élève, comme en français dans « Tu aimes internet ? ».
\end{propriete}

\begin{exemplebox}[Affirmation et question]
\zh{你喜欢上网。} \emph{Nǐ xǐhuan shàngwǎng.} « Tu aimes aller sur internet. » (voix qui descend à la fin)\\
\zh{你喜欢上网吗？} \emph{Nǐ xǐhuan shàngwǎng ma?} « Aimes-tu aller sur internet ? » (voix qui monte à la fin, mais 喜 reste au 3e ton et 上 au 4e ton)
\end{exemplebox}

\begin{attention}[Erreur fréquente des francophones : le 4e ton agressif]
En français, une voix qui descend brusquement marque souvent la colère ou l'insistance (« Non ! »). Les francophones transforment donc le 4e ton en accent d'insistance, ce qui les fait paraître agressifs. Le 4e ton chinois est \cle{bref et sec}, pas violent : dans \zh{上网} shàngwǎng « aller sur internet », laissez simplement la voix tomber vite, sans appuyer. De même, ne montez pas la voix sur \zh{吗} comme sur un mot interrogatif français : la montée est légère et globale.
\end{attention}

\begin{methode}[Travailler les tons : la méthode du rythme frappé]
\begin{enumerate}
\item \cle{Frappez le rythme de la main} : la main reste haute et plate pour le 1er ton, monte pour le 2e, descend-remonte pour le 3e, tombe d'un coup sec pour le 4e.
\item Prononcez chaque mot lentement en accompagnant le geste : \zh{网} wǎng, \zh{聊} liáo, \zh{好} hǎo, \zh{用} yòng.
\item Enchaînez les mots deux par deux en gardant le geste : \zh{上网} shàngwǎng, \zh{聊天} liáotiān.
\item Dites la phrase entière \cle{sans ralentir}, en gardant le rythme intérieur : \zh{我觉得上网很好。} Wǒ juéde shàngwǎng hěn hǎo.
\item Répétez sans le geste, à vitesse normale de conversation.
\end{enumerate}
\end{methode}

\begin{exoresolu}[Ping-pong tonal]
\textbf{Énoncé.} L'enseignant dit un mot ; l'élève répond avec la phrase complète du débat en respectant les tons. Appliquez le sandhi du 3e ton et le changement de ton de 不.\\
1. Professeur : \zh{方便} fāngbiàn « pratique ».\\
2. Professeur : \zh{同意} tóngyì « être d'accord ».\\
3. Professeur : \zh{好处} hǎochù « avantage ».
\tcblower
\textbf{Solution détaillée.}\\
1. \zh{我觉得上网很方便。} \emph{Wǒ juéde shàngwǎng hěn fāngbiàn.} « Je trouve qu'aller sur internet est très pratique. » Attention : \zh{很} hěn devant \zh{方} fāng (1er ton) garde son 3e ton, sans sandhi ; dans la parole rapide, on ne prononce que sa partie basse.\\
2. \zh{我不同意你的看法。} \emph{Wǒ bù tóngyì nǐ de kànfǎ.} « Je ne suis pas d'accord avec ton point de vue. » Attention : \zh{不} devant \zh{同} tóng (2e ton) reste au 4e ton : on prononce \emph{bù tóngyì}, sans changement. Le changement bù → bú n'a lieu que devant un 4e ton (par exemple \zh{不是} bú shì).\\
3. \zh{上网有很多好处。} \emph{Shàngwǎng yǒu hěn duō hǎochù.} « Internet a beaucoup d'avantages. » Attention : \zh{很多} hěn duō ne change pas, car \zh{多} duō est au 1er ton (pas de sandhi) ; \zh{好处} hǎochù (3 + 4) ne change pas non plus. En revanche, dans \zh{很好} hěn hǎo (3 + 3), on prononcerait \emph{hén hǎo}.
\end{exoresolu}

\begin{experience}[Enregistrement et auto-évaluation]
Par groupes de deux :
\begin{enumerate}
\item Chaque élève s'enregistre (téléphone ou magnétophone) en lisant le dialogue du débat : \zh{A：你喜欢上网吗？B：喜欢，我觉得很方便。A：我不同意，上网也有很多问题。} (\emph{A : Nǐ xǐhuan shàngwǎng ma? B : Xǐhuan, wǒ juéde hěn fāngbiàn. A : Wǒ bù tóngyì, shàngwǎng yě yǒu hěn duō wèntí.} « A : Tu aimes aller sur internet ? B : Oui, je trouve ça très pratique. A : Je ne suis pas d'accord, internet pose aussi beaucoup de problèmes. »)
\item Réécoutez-vous et comparez avec le modèle lu par l'enseignant.
\item Notez sur une fiche : tons justes / tons faux, sandhi appliqué ou non, 4e ton sec ou agressif.
\item Échangez les fiches avec votre camarade et refaites l'enregistrement en corrigeant une seule erreur à la fois.
\end{enumerate}
\end{experience}

\begin{exercice}[À vous de jouer]
1. Indiquez le ton de chaque syllabe et les changements de prononciation : \zh{我} , \zh{不} , \zh{很好} , \zh{不是} , \zh{好处} , \zh{可以}.\\
2. Lisez à voix haute en frappant le rythme de la main : \zh{我觉得上网很有意思，但是也有问题。}\\
3. Transformez en question avec \zh{吗} et lisez-la avec l'intonation montante : \zh{你常常聊天。}
\end{exercice}

\begin{corrige}[Exercice]
1. \zh{我} wǒ : 3e ton ; \zh{不} bù : 4e ton, devient bú devant un 4e ton ; \zh{很好} hěn hǎo → hén hǎo (sandhi, 3 + 3) ; \zh{不是} bù shì → bú shì (不 devant un 4e ton) ; \zh{好处} hǎochù : 3 + 4, pas de changement ; \zh{可以} kěyǐ → kéyǐ (sandhi, 3 + 3).\\
2. \emph{Wǒ juéde shàngwǎng hěn yǒu yìsi, dànshì yě yǒu wèntí.} « Je trouve internet très intéressant, mais il y a aussi des problèmes. » Vérifier : \zh{我} wǒ devant \zh{觉得} juéde (2e ton) garde son 3e ton réduit à sa partie basse ; \zh{很有} hěn yǒu (3 + 3) se prononce \emph{hén yǒu} (sandhi) ; \zh{意思} yìsi : 4e ton + ton neutre.\\
3. \zh{你常常聊天吗？} \emph{Nǐ chángcháng liáotiān ma?} « Tu bavardes souvent (en ligne) ? » La voix monte légèrement à la fin ; \zh{聊} garde son 2e ton et \zh{天} son 1er ton.
\end{corrige}

\coursec{Jeux de rôle et débat guidé de classe}

Dans la section précédente, nous avons travaillé la prononciation des tons dans les phrases du débat : vous savez maintenant dire \zh{我认为} (\emph{Wǒ rènwéi}, « je pense que ») avec le troisième ton suivi du quatrième, sans les confondre. Il est temps de passer à la pratique réelle : parler, réagir, argumenter. Cette dernière étape de la leçon vous transforme en acteurs : journaliste, parent, enfant, animateur de débat. L'objectif n'est plus de réciter, mais de \cle{communiquer spontanément} en chinois, comme vous le ferez à l'évaluation orale.

\courssub{Jeu de rôle 1 : l'interview sur les habitudes internet}

\begin{experience}[Journaliste et interviewé]
Par groupes de deux. L'élève A joue le \cle{journaliste} (\zh{记者}, \emph{jìzhě}) d'un journal lycéen de Yaoundé ; l'élève B joue l'\cle{interviewé} (\zh{被采访的人}, \emph{bèi cǎifǎng de rén}, « la personne interviewée »). Le journaliste pose au moins quatre questions sur les habitudes internet de son camarade ; l'interviewé répond par des phrases complètes, puis on inverse les rôles.
\end{experience}

Observez d'abord ce modèle d'échange avant de jouer :

\begin{exemplebox}[Modèle d'interview]
\zh{记者：你好！我可以问你几个问题吗？}\\
\emph{Jìzhě : Nǐ hǎo ! Wǒ kěyǐ wèn nǐ jǐ ge wèntí ma ?}\\
« Journaliste : Bonjour ! Puis-je te poser quelques questions ? »\\[4pt]
\zh{学生：当然可以。}\\
\emph{Xuésheng : Dāngrán kěyǐ.}\\
« Élève : Bien sûr. »\\[4pt]
\zh{记者：你每天上网多长时间？}\\
\emph{Jìzhě : Nǐ měitiān shàngwǎng duō cháng shíjiān ?}\\
« Journaliste : Combien de temps passes-tu sur internet chaque jour ? »\\[4pt]
\zh{学生：我每天上网两个小时，我常常用手机查资料。}\\
\emph{Xuésheng : Wǒ měitiān shàngwǎng liǎng ge xiǎoshí, wǒ chángcháng yòng shǒujī chá zīliào.}\\
« Élève : Je passe deux heures par jour sur internet, j'utilise souvent mon téléphone pour chercher des informations. »\\[4pt]
\zh{记者：你有什么看法？}\\
\emph{Jìzhě : Nǐ yǒu shénme kànfǎ ?}\\
« Journaliste : Quel est ton avis ? »\\[4pt]
\zh{学生：我认为我们应该少玩游戏，多学习。}\\
\emph{Xuésheng : Wǒ rènwéi wǒmen yīnggāi shǎo wán yóuxì, duō xuéxí.}\\
« Élève : Je pense que nous devrions moins jouer aux jeux et plus étudier. »
\end{exemplebox}

\begin{attention}[La durée : \zh{多长时间} et non \zh{多少时间}]
Pour demander une durée, on dit \zh{多长时间} (\emph{duō cháng shíjiān}, « combien de temps ») : le mot \zh{长} (\emph{cháng}, « long ») est indispensable. Ne traduisez pas mot à mot « combien de temps » par \zh{多少时间}, qui est incorrect ou très lourd. De même, pour répondre, la durée se place après le verbe : \zh{我每天上网两个小时} (littéralement « je chaque jour monte-sur-le-réseau deux heures »), jamais \zh{我两个小时上网} dans ce contexte.
\end{attention}

\begin{methode}[Réussir son jeu de rôle en trois étapes]
\begin{enumerate}
\item \textbf{Préparer} : écrivez trois phrases en utilisant les formules vues dans la leçon (\zh{我认为…}, \zh{我觉得…}, \zh{你每天上网多长时间？}).
\item \textbf{Écouter} : pendant que votre partenaire parle, regardez-le et repérez un mot-clé de sa réponse pour construire la question suivante.
\item \textbf{Réagir} : avant de répondre ou d'enchaîner, marquez votre position avec \zh{我同意} (\emph{Wǒ tóngyì}, « je suis d'accord ») ou \zh{我不同意} (\emph{Wǒ bù tóngyì}, « je ne suis pas d'accord »), puis développez.
\end{enumerate}
\end{methode}

\courssub{Jeu de rôle 2 : discussion parent-enfant}

\begin{experience}[Le parent inquiet et l'élève qui se défend]
Par groupes de deux. L'élève A joue le \cle{parent} (\zh{爸爸} \emph{bàba} ou \zh{妈妈} \emph{māma}), convaincu qu'internet est dangereux : trop de jeux, mauvaises informations, moins de temps pour les devoirs. L'élève B joue l'\cle{enfant} (\zh{孩子}, \emph{háizi}), qui défend l'usage d'internet pour les études : chercher des cours, réviser le chinois, communiquer avec les camarades. Chacun doit utiliser au moins deux arguments et une formule d'avis.
\end{experience}

\begin{exemplebox}[Répliques utiles pour chaque camp]
Côté parent :\\
\zh{我觉得上网太危险，你应该少用手机。}\\
\emph{Wǒ juéde shàngwǎng tài wēixiǎn, nǐ yīnggāi shǎo yòng shǒujī.}\\
« Je trouve qu'internet est trop dangereux, tu devrais moins utiliser ton téléphone. »\\[4pt]
Côté enfant :\\
\zh{我不同意。我用网络查资料，这对我的学习很有帮助。}\\
\emph{Wǒ bù tóngyì. Wǒ yòng wǎngluò chá zīliào, zhè duì wǒ de xuéxí hěn yǒu bāngzhù.}\\
« Je ne suis pas d'accord. J'utilise internet pour chercher des informations, cela aide beaucoup mes études. »
\end{exemplebox}

\begin{propriete}[Exprimer l'utilité avec \zh{对…有帮助}]
Structure : \textbf{A + 对 + B + 有（很）帮助}, « A est (très) utile à B ». Exemple : \zh{网络对我的学习很有帮助。} (\emph{Wǎngluò duì wǒ de xuéxí hěn yǒu bāngzhù.}) « Internet est très utile à mes études. » Attention à l'ordre : en chinois, on dit « internet \emph{envers} mes études a de l'aide », et non « mes études ont de l'aide d'internet ». La négation est \zh{对…没有帮助} : \zh{玩游戏太多对学习没有帮助。} « Jouer trop n'est pas utile aux études. »
\end{propriete}

\begin{attention}[Ne lisez pas vos notes mot à mot]
Pendant le jeu de rôle et le débat, beaucoup d'élèves francophones écrivent tout leur texte et le lisent les yeux baissés : l'examinateur ne peut alors ni entendre les tons ni évaluer la communication réelle. Préparez seulement des \cle{mots-clés} (trois ou quatre caractères par idée), regardez votre interlocuteur dans les yeux et parlez naturellement, même lentement. Une phrase simple bien prononcée vaut mieux qu'un long texte récité sans regard.
\end{attention}

\courssub{Le débat guidé : organisation et rôles}

La classe se divise en deux camps. Le premier groupe défend la thèse \zh{互联网好处多} (\emph{Hùliánwǎng hǎochù duō}, « internet a plus d'avantages ») ; le second défend \zh{互联网坏处多} (\emph{Hùliánwǎng huàichù duō}, « internet a plus d'inconvénients »). Trois rôles sont distribués : l'animateur, les orateurs des deux groupes et le chronométreur.

\begin{center}
\begin{tabularx}{\linewidth}{|X|X|X|X|}
\hline
\rowcolor{popblueL} Caractères & Pinyin & Nature & Traduction \\
\hline
主持人 & zhǔchírén & nom & animateur, présentateur \\
\hline
第一组 & dì-yī zǔ & nom & le premier groupe \\
\hline
第二组 & dì-èr zǔ & nom & le deuxième groupe \\
\hline
计时员 & jìshíyuán & nom & chronométreur \\
\hline
发言 & fāyán & verbe / nom & prendre la parole, intervention \\
\hline
提问 & tíwèn & verbe & poser une question \\
\hline
总结 & zǒngjié & verbe / nom & résumer, conclusion \\
\hline
看法 & kànfǎ & nom & avis, point de vue \\
\hline
题目 & tímù & nom & sujet, thème \\
\hline
补充 & bǔchōng & verbe & compléter, ajouter \\
\hline
\end{tabularx}
\end{center}

Le \cle{chronométreur} (\zh{计时员}) donne à chaque orateur une minute de parole. Chaque orateur doit obligatoirement utiliser \textbf{au moins une formule d'avis} (\zh{我认为…} / \zh{我觉得…}) et \textbf{au moins une relance} (\zh{你有什么看法？} « quel est ton avis ? » ou \zh{你同意吗？} « es-tu d'accord ? »).

\begin{propriete}[La question alternative avec \zh{还是}]
Pour poser une question qui propose un choix entre deux options, on utilise \zh{还是} (\emph{háishi}, « ou » dans les questions) : \textbf{option A + 还是 + option B + ？}. Exemple : \zh{互联网好处多还是坏处多？} (\emph{Hùliánwǎng hǎochù duō háishi huàichù duō ?}) « Internet a-t-il plus d'avantages ou d'inconvénients ? » Attention : dans une phrase affirmative, « ou » se dit \zh{或者} (\emph{huòzhě}) ; \zh{还是} est réservé aux questions. Ne dites donc pas \zh{我喝茶还是咖啡} pour « je bois du thé ou du café », mais \zh{我喝茶或者咖啡}.
\end{propriete}

\begin{textebox}[L'animateur ouvre le débat]
主持人：大家好！今天的题目是“互联网好处多还是坏处多”。请第一组先说。\\
\emph{Zhǔchírén : Dàjiā hǎo ! Jīntiān de tímù shì “hùliánwǎng hǎochù duō háishi huàichù duō”. Qǐng dì-yī zǔ xiān shuō.}\\
« L'animateur : Bonjour à tous ! Le sujet d'aujourd'hui est : internet a-t-il plus d'avantages ou d'inconvénients ? Le premier groupe commence, s'il vous plaît. »\\[4pt]
第一组学生：我认为互联网好处多，因为我们可以查资料、学习汉语。\\
\emph{Dì-yī zǔ xuésheng : Wǒ rènwéi hùliánwǎng hǎochù duō, yīnwèi wǒmen kěyǐ chá zīliào, xuéxí Hànyǔ.}\\
« Élève du premier groupe : Je pense qu'internet a plus d'avantages, parce que nous pouvons chercher des informations et apprendre le chinois. »\\[4pt]
第二组学生：我不同意。我觉得互联网坏处多，很多年轻人玩游戏太多。你有什么看法？\\
\emph{Dì-èr zǔ xuésheng : Wǒ bù tóngyì. Wǒ juéde hùliánwǎng huàichù duō, hěn duō niánqīngrén wán yóuxì tài duō. Nǐ yǒu shénme kànfǎ ?}\\
« Élève du deuxième groupe : Je ne suis pas d'accord. Je trouve qu'internet a plus d'inconvénients, beaucoup de jeunes jouent trop aux jeux. Quel est ton avis ? »\\[4pt]
主持人：很好！请第二组继续说，然后大家互相提问。\\
\emph{Zhǔchírén : Hěn hǎo ! Qǐng dì-èr zǔ jìxù shuō, ránhòu dàjiā hùxiāng tíwèn.}\\
« L'animateur : Très bien ! Le deuxième groupe continue, puis tout le monde se pose des questions mutuellement. »
\end{textebox}

\begin{popfigure}
\begin{tikzpicture}[node distance=0.9cm and 1.6cm,
  etape/.style={rectangle, rounded corners=4pt, draw=popblue, fill=popblueL, thick, align=center, minimum height=1cm, text width=3.4cm, font=\small},
  fleche/.style={-{Stealth[length=3mm]}, thick, popdark}]
\node[etape, align=center] (a){主持人介绍主题\\\emph{Zhǔchírén jièshào zhǔtí}\\L'animateur présente le sujet};
\node[etape, right=of a, fill=popgreenL, draw=popgreen, align=center] (b){第一组发言\\\emph{Dì-yī zǔ fāyán}\\Le groupe 1 prend la parole};
\node[etape, right=of b, fill=poporangeL, draw=poporange, align=center] (c){第二组发言\\\emph{Dì-èr zǔ fāyán}\\Le groupe 2 prend la parole};
\node[etape, below=of b, fill=poppurpleL, draw=poppurple, align=center] (d){互相提问\\\emph{Hùxiāng tíwèn}\\Questions croisées};
\node[etape, left=of d, fill=popgoldL, draw=popgold, align=center] (e){主持人总结\\\emph{Zhǔchírén zǒngjié}\\L'animateur conclut};
\draw[fleche] (a) -- (b);
\draw[fleche] (b) -- (c);
\draw[fleche] (c) |- (d);
\draw[fleche] (d) -- (e);
\end{tikzpicture}
\legende{Déroulement du débat guidé en cinq étapes : de la présentation du sujet à la conclusion de l'animateur.}
\end{popfigure}

\begin{propriete}[Règle du tour de parole]
En chinois comme en français, un débat respecte un ordre strict : on ne coupe jamais la parole. Pour demander la parole poliment, on lève la main et on dit \zh{我可以说吗？} (\emph{Wǒ kěyǐ shuō ma ?}, « puis-je parler ? ») ou \zh{我想补充一点} (\emph{Wǒ xiǎng bǔchōng yì diǎn}, « je voudrais ajouter un point »). L'animateur distribue la parole avec \zh{请你说} (\emph{qǐng nǐ shuō}, « à toi, s'il te plaît »).
\end{propriete}

\courssub{Grille d'observation et débriefing collectif}

Pendant le débat, les élèves qui ne parlent pas ne sont pas spectateurs : ils observent et remplissent la grille suivante pour chaque orateur (note de 1 à 4 par critère).

\begin{center}
\begin{tabularx}{\linewidth}{|X|X|X|}
\hline
\rowcolor{popblueL} Critère & Ce qu'on observe & Note (1 à 4) \\
\hline
Prononciation des tons & Les tons sont-ils corrects, surtout les 3\textsuperscript{e} et 4\textsuperscript{e} tons ? & \\
\hline
Formules fonctionnelles & Au moins une formule d'avis et une relance utilisées ? & \\
\hline
Fluidité & Le débit est-il naturel, sans longues hésitations ? & \\
\hline
Tour de parole & L'orateur écoute-t-il les autres sans les interrompre ? & \\
\hline
Regard et posture & Parle-t-il en regardant son interlocuteur, sans lire ? & \\
\hline
\end{tabularx}
\end{center}

Après le débat, la classe fait un \cle{débriefing collectif} : chaque orateur entend d'abord ses \textbf{points forts} (un ton bien réussi, une belle relance), puis un seul \textbf{axe de progrès} concret (par exemple : « fais attention au troisième ton de \zh{我} »). Le professeur termine en félicitant le groupe qui a le mieux respecté les règles du débat, pas nécessairement celui qui a « gagné ».

\begin{exoresolu}[Préparer une intervention de débat]
Énoncé : vous êtes dans le groupe \zh{互联网好处多}. Rédigez une intervention de trois phrases avec une formule d'avis, un argument avec \zh{因为} et une relance.
\tcblower
Solution détaillée :
\begin{enumerate}
\item Formule d'avis : \zh{我认为互联网好处多。} (\emph{Wǒ rènwéi hùliánwǎng hǎochù duō.}) « Je pense qu'internet a plus d'avantages. »
\item Argument avec \zh{因为} : \zh{因为我们可以在网上查资料、学习外语。} (\emph{Yīnwèi wǒmen kěyǐ zài wǎngshàng chá zīliào, xuéxí wàiyǔ.}) « Parce que nous pouvons chercher des informations et apprendre des langues étrangères en ligne. »
\item Relance : \zh{你同意吗？你有什么看法？} (\emph{Nǐ tóngyì ma ? Nǐ yǒu shénme kànfǎ ?}) « Es-tu d'accord ? Quel est ton avis ? »
\end{enumerate}
Vérifiez : avis + argument + relance, tons correctement prononcés, regard vers l'interlocuteur.
\end{exoresolu}

\begin{exercice}[À vous de jouer]
\begin{enumerate}
\item En binôme, rejouez le jeu de rôle parent-enfant en inversant les rôles : chacun doit défendre le camp opposé au sien.
\item Préparez en cinq minutes une intervention pour le camp \zh{互联网坏处多} avec une formule d'avis, un argument et une relance.
\item Pendant le débat de la classe, remplissez la grille d'observation pour deux orateurs et formulez pour chacun un point fort et un axe de progrès.
\end{enumerate}
\end{exercice}

\begin{corrige}[Exercice 2 — éléments de réponse]
Intervention possible : \zh{我觉得互联网坏处多，因为很多年轻人玩游戏太多，不学习。你同意吗？} (\emph{Wǒ juéde hùliánwǎng huàichù duō, yīnwèi hěn duō niánqīngrén wán yóuxì tài duō, bù xuéxí. Nǐ tóngyì ma ?}) « Je trouve qu'internet a plus d'inconvénients, parce que beaucoup de jeunes jouent trop et n'étudient pas. Es-tu d'accord ? » — Toute réponse contenant une formule d'avis, un argument introduit par \zh{因为} et une relance est acceptée.
\end{corrige}

\begin{savaistu}[À l'évaluation orale, la spontanéité compte]
Au Bac blanc et aux évaluations orales de chinois, l'examinateur valorise la capacité à \cle{réagir spontanément}, pas seulement à réciter. Un élève qui écoute la question, dit \zh{我同意} ou \zh{我不同意} et répond avec une phrase simple mais correcte obtient une meilleure note qu'un élève qui récite un long texte appris par cœur sans rapport avec la question. Entraînez-vous donc à improviser : c'est la compétence la mieux payée le jour de l'examen.
\end{savaistu}

\newpage

\coursec{Activité d'intégration}

\courssub{Objectifs de l'activité}

Cette activité d'intégration clôt la leçon 37 par une tâche finale de communication orale : un \cle{grand débat} de classe sur le thème \zh{互联网，好处多还是坏处多？} (\emph{Hùliánwǎng, hǎochù duō háishi huàichù duō?} « Internet : plus d'avantages ou plus d'inconvénients ? »). À l'issue de l'activité, l'élève doit être capable de :

\begin{itemize}
\item présenter oralement un point de vue simple en chinois, en 4 à 6 phrases ;
\item argumenter avec la structure \zh{因为……所以……} (\emph{yīnwèi... suǒyǐ...} « parce que... donc... ») ;
\item réagir à l'avis d'un camarade avec \zh{我同意} / \zh{我不同意} et relancer la parole avec \zh{你呢？} ou \zh{你觉得怎么样？} ;
\item respecter la prononciation des tons dans un échange spontané ;
\item conclure un débat avec \zh{总之} (\emph{zǒngzhī} « en résumé »).
\end{itemize}

\courssub{Situation}

Le club de chinois du lycée de Biyem-Assi, à Yaoundé, prépare une journée portes ouvertes. Les élèves de Première A doivent montrer aux visiteurs qu'ils savent discuter en chinois d'un sujet de société. Le professeur a choisi un thème qui concerne tout le monde : internet. Presque tous les élèves utilisent WhatsApp, regardent des vidéos ou font leurs recherches en ligne, mais les parents et les professeurs critiquent souvent le temps passé sur les écrans. La classe se divise donc en deux équipes : l'équipe A défend les avantages d'internet (\zh{好处多}), l'équipe B défend les inconvénients (\zh{坏处多}). Un élève animateur dirige le débat, et un jury de quatre élèves note les prestations.

Pour lancer la réflexion, le professeur affiche au tableau un court texte d'appui :

\begin{textebox}[Texte d'appui : internet chez les jeunes]
现在，很多年轻人每天都上网。他们用手机聊天、看视频、学习汉语。互联网很有用：我们可以在网上学习，也可以和很远的朋友联系。但是，上网太多也有问题：有的学生晚上不睡觉，上课的时候很累；还有的人在网上看不好的东西。所以，我们应该好好用互联网。\\
\emph{Xiànzài, hěn duō niánqīngrén měi tiān dōu shàngwǎng. Tāmen yòng shǒujī liáotiān, kàn shìpín, xuéxí Hànyǔ. Hùliánwǎng hěn yǒuyòng: wǒmen kěyǐ zài wǎngshang xuéxí, yě kěyǐ hé hěn yuǎn de péngyou liánxì. Dànshì, shàngwǎng tài duō yě yǒu wèntí: yǒu de xuésheng wǎnshang bù shuìjiào, shàngkè de shíhou hěn lèi; hái yǒu de rén zài wǎngshang kàn bù hǎo de dōngxi. Suǒyǐ, wǒmen yīnggāi hǎohāo yòng hùliánwǎng.}\\
« Aujourd'hui, beaucoup de jeunes se connectent chaque jour. Ils utilisent leur téléphone pour discuter, regarder des vidéos, apprendre le chinois. Internet est très utile : on peut étudier en ligne et rester en contact avec des amis lointains. Mais trop de connexion pose aussi des problèmes : certains élèves ne dorment pas la nuit et sont fatigués en classe ; d'autres regardent de mauvaises choses en ligne. Nous devons donc bien utiliser internet. »
\end{textebox}

\courssub{Consignes}

\begin{enumerate}
\item \textbf{Compréhension du texte d'appui (5 min).} Lis le texte affiché au tableau. Souligne trois mots du lexique d'internet et relève une phrase avec \zh{因为} ou \zh{所以}. Restitue à voix haute, en une phrase, l'idée principale du texte.
\item \textbf{Préparation en équipe (10 min).} Dans ton équipe (A : avantages ; B : inconvénients), rédige par écrit \cle{trois arguments}. Chaque argument doit contenir : une formule d'opinion (\zh{我觉得……}, \zh{我认为……}), un connecteur logique (\zh{因为……所以……}, \zh{但是}, \zh{而且}) et un mot du vocabulaire d'internet (\zh{上网}, \zh{手机}, \zh{网站}, \zh{聊天}, \zh{信息}).
\item \textbf{Présentation des arguments (2 x 3 min).} L'animateur ouvre le débat avec la formule \zh{大家好！今天我们讨论……}. Chaque orateur parle \cle{une minute maximum}, sans lire intégralement sa fiche : il présente l'argument de son équipe en regardant la classe.
\item \textbf{Questions-réponses croisées (5 min).} Chaque équipe pose une question à l'équipe adverse. Pour réagir, tu dois obligatoirement commencer par \zh{我同意你的看法} ou \zh{我不同意，因为……}, puis \cle{relancer} un camarade avec \zh{你呢？} ou \zh{你觉得怎么样？}.
\item \textbf{Évaluation par le jury.} Le jury note chaque orateur avec la grille d'observation : prononciation des tons (2 points), formules fonctionnelles utilisées (2 points), fluidité et prise de parole sans lecture (1 point). Le jury désigne l'équipe la plus convaincante.
\item \textbf{Synthèse finale (2 min).} L'animateur conclut le débat à voix haute en commençant par \zh{总之……} et en donnant sa conclusion personnelle équilibrée.
\end{enumerate}

\courssub{Ressources à mobiliser}

\begin{aretenir}[Boîte à outils du débatteur]
\textbf{Présenter :} \zh{大家好！今天我们讨论互联网的问题。} (\emph{Dàjiā hǎo! Jīntiān wǒmen tǎolùn hùliánwǎng de wèntí.} « Bonjour à tous ! Aujourd'hui nous discutons de la question d'internet. »)\\
\textbf{Donner son avis :} \zh{我觉得……} / \zh{我认为……} / \zh{对我来说，……}\\
\textbf{Argumenter :} \zh{因为……，所以……} ; \zh{第一……，第二……} ; \zh{不但……，而且……}\\
\textbf{Réagir :} \zh{我同意你的看法。} / \zh{我不同意，因为……}\\
\textbf{Relancer :} \zh{你呢？} / \zh{你觉得怎么样？} / \zh{你有什么看法？}\\
\textbf{Conclure :} \zh{总之，……}
\end{aretenir}

\begin{center}
\begin{tabularx}{\linewidth}{|X|X|X|X|}
\hline
\rowcolor{popblueL} Caractères & Pinyin & Nature & Traduction \\ \hline
上网 & shàngwǎng & verbe & se connecter à internet \\ \hline
好处 & hǎochù & nom & avantage \\ \hline
坏处 & huàichù & nom & inconvénient \\ \hline
聊天 & liáotiān & verbe & bavarder (en ligne) \\ \hline
信息 & xìnxī & nom & information \\ \hline
有用 & yǒuyòng & adjectif & utile \\ \hline
浪费时间 & làngfèi shíjiān & loc. verb. & perdre du temps \\ \hline
看法 & kànfǎ & nom & point de vue, opinion \\ \hline
\end{tabularx}
\end{center}

\begin{attention}[Pièges à éviter à l'oral]
\begin{itemize}
\item Ne dis pas \zh{我觉得是好处多} : après \zh{觉得}, pas de \zh{是} devant un adjectif. Dis \zh{我觉得好处很多。}
\item Attention aux tons de \zh{好处} (\emph{hǎo} 3\up{e} ton) et \zh{坏处} (\emph{huài} 4\up{e} ton) : les confondre inverse ton argument !
\item \zh{因为……所以……} : en chinois, on peut utiliser les deux ensemble, contrairement au français (« parce que... donc » est correct en chinois).
\end{itemize}
\end{attention}

\courssub{Proposition de production et corrigé commenté}

\begin{exoresolu}[Grand débat : production modèle]
Énoncé : tu es l'orateur de l'équipe A (avantages). Présente un argument en 4 à 6 phrases, réagis à l'équipe B, puis l'animateur conclut.
\tcblower
\textbf{Orateur A :} \zh{大家好！我觉得互联网的好处很多。因为我们可以从网上找到很多信息，所以学习更方便了。而且，我们还可以和外国朋友聊天，练习汉语。}\\
\emph{Dàjiā hǎo! Wǒ juéde hùliánwǎng de hǎochù hěn duō. Yīnwèi wǒmen kěyǐ cóng wǎngshang zhǎodào hěn duō xìnxī, suǒyǐ xuéxí gèng fāngbiàn le. Érqiě, wǒmen hái kěyǐ hé wàiguó péngyou liáotiān, liànxí Hànyǔ.}\\
« Bonjour à tous ! Je pense qu'internet a beaucoup d'avantages. Parce qu'on peut trouver beaucoup d'informations en ligne, étudier est devenu plus pratique. De plus, on peut discuter avec des amis étrangers et pratiquer le chinois. »

\textbf{Réaction de l'équipe B :} \zh{我不同意。因为很多学生上网时间太长，所以上课的时候很累。我觉得坏处也很多。你呢？}\\
\emph{Wǒ bù tóngyì. Yīnwèi hěn duō xuésheng shàngwǎng shíjiān tài cháng, suǒyǐ shàngkè de shíhou hěn lèi. Wǒ juéde huàichù yě hěn duō. Nǐ ne?}\\
« Je ne suis pas d'accord. Parce que beaucoup d'élèves passent trop de temps en ligne, ils sont fatigués en classe. Je pense que les inconvénients sont nombreux aussi. Et toi ? »

\textbf{Synthèse de l'animateur :} \zh{总之，互联网有好处，也有坏处。我们应该好好用它，不要浪费时间。谢谢大家！}\\
\emph{Zǒngzhī, hùliánwǎng yǒu hǎochù, yě yǒu huàichù. Wǒmen yīnggāo hǎohāo yòng tā, búyào làngfèi shíjiān. Xièxie dàjiā!}\\
« En résumé, internet a des avantages et des inconvénients. Nous devons bien l'utiliser et ne pas perdre de temps. Merci à tous ! »

\textbf{Commentaire des choix de langue :} l'orateur ouvre avec la formule protocolaire \zh{大家好}, annonce son avis avec \zh{我觉得}, puis enchaîne deux arguments reliés par \zh{因为……所以……} et \zh{而且}. La réaction respecte la consigne : \zh{我不同意} + justification + relance \zh{你呢？}. La synthèse utilise \zh{总之} et une structure équilibrée \zh{有……，也有……}, ce qui montre une prise de parole nuancée, attendue au niveau Première.
\end{exoresolu}

\courssub{Bilan de l'activité}

À l'issue du débat, chaque élève a mobilisé l'ensemble des savoir-faire de la leçon : comprendre un texte sur internet, réutiliser le lexique (\zh{上网}, \zh{好处}, \zh{坏处}, \zh{信息}), construire un argument avec \zh{因为……所以……}, réagir avec \zh{我同意} / \zh{我不同意}, relancer avec \zh{你呢？} et conclure avec \zh{总之}. La grille du jury a permis d'objectiver la prononciation des tons et la fluidité. Pour prolonger l'activité à la maison, chaque élève peut enregistrer sur son téléphone une prise de position d'une minute sur le même sujet, en veillant aux tons de \zh{hǎochù} et \zh{huàichù}, et la faire écouter à un camarade au début de la séance suivante.

\newpage

\coursec{Méthodes et exercices résolus}

\courssub{Méthodes et exercices résolus}

\begin{methode}[Structurer un exposé oral court (2--3 minutes)]
\textbf{Objectif :} Présenter un sujet lié à Internet de façon claire, fluide et structurée.
\begin{enumerate}
    \item \textbf{Analyser le sujet} : identifier les mots-clés (ex. : \zh{网购} wǎnggòu, \zh{社交媒体} shèjiāo méitǐ, \zh{网课} wǎngkè). Définir l'angle d'attaque (mon expérience / avantages-inconvénients / conseil).
    \item \textbf{Construire le squelette} (3 parties) :
    \begin{itemize}
        \item \textbf{Introduction} : accroche + annonce du plan (\zh{今天我想谈谈……} Jīntiān wǒ xiǎng tántan… / \zh{我的题目是……} Wǒ de tímù shì…).
        \item \textbf{Développement} : 2--3 arguments reliés par des connecteurs logiques (\zh{首先……其次……最后……} shǒuxiān… qícì… zuìhòu… / \zh{一方面……另一方面……} yì fāngmiàn… lìng yì fāngmiàn…). Illustrer par un exemple personnel (Cameroun/Chine).
        \item \textbf{Conclusion} : résumé + opinion personnelle + ouverture (\zh{总之……我认为……希望……} zǒngzhī… wǒ rènwéi… xīwàng…).
    \end{itemize}
    \item \textbf{Rédiger des notes-mots} (pas de phrases complètes) : mots-clés + connecteurs + exemples.
    \item \textbf{Répéter à voix haute} : chronométrer, vérifier les tons des mots-clés, fluidifier les transitions. S'enregistrer si possible.
\end{enumerate}
\cle{Structure type : 开头 (début) -- 主体 (corps) -- 结尾 (fin)}
\end{methode}

\begin{exoresolu}[Exposé : « Mon utilisation d'Internet pour les études »]
\textbf{Énoncé :} Préparez un exposé de 2 minutes sur le thème : \zh{我如何利用互联网学习} (Wǒ rúhé lìyòng hùliánwǎng xuéxí — Comment j'utilise Internet pour apprendre). Utilisez au moins 5 mots du vocabulaire de la leçon et 3 connecteurs logiques. Rédigez le texte intégral de l'exposé (caractères, pinyin, traduction).
\tcblower
\textbf{Solution détaillée :}
\begin{enumerate}
    \item \textbf{Plan choisi :} introduction (rôle général) $\rightarrow$ 2 usages précis (cours vidéo, dictionnaires/applis) $\rightarrow$ problème (distraction) $\rightarrow$ conclusion (équilibre).
    \item \textbf{Vocabulaire ciblé :} \zh{网课} (wǎngkè), \zh{搜索} (sōusuǒ), \zh{翻译软件} (fānyì ruǎnjiàn), \zh{分心} (fēnxīn), \zh{自律} (zìlǜ).
    \item \textbf{Connecteurs :} \zh{首先} (shǒuxiān), \zh{而且} (érqiě), \zh{但是} (dànshì), \zh{所以} (suǒyǐ), \zh{总之} (zǒngzhī).
\end{enumerate}

\textbf{Exposé modèle (caractères + pinyin + traduction) :}

各位同学，大家好。\\
Gèwèi tóngxué, dàjiā hǎo.\\
« Chers camarades, bonjour à tous. »

今天我想谈谈我如何利用互联网学习。\\
Jīntiān wǒ xiǎng tántan wǒ rúhé lìyòng hùliánwǎng xuéxí.\\
« Aujourd'hui, je voudrais parler de la façon dont j'utilise Internet pour étudier. »

首先，我经常上网课。比如中文课，我可以看视频复习语法。\\
Shǒuxiān, wǒ jīngcháng shàng wǎngkè. Bǐrú Zhōngwén kè, wǒ kěyǐ kàn shìpín fùxí yǔfǎ.\\
« D'abord, je suis souvent des cours en ligne. Par exemple, en cours de chinois, je peux regarder des vidéos pour réviser la grammaire. »

而且，搜索资料很方便。我常用翻译软件查生词。\\
Érqiě, sōusuǒ zīliào hěn fāngbiàn. Wǒ cháng yòng fānyì ruǎnjiàn chá shēngcí.\\
« De plus, chercher des documents est très pratique. J'utilise souvent un logiciel de traduction pour chercher les mots nouveaux. »

但是，网上也有很多视频和游戏，容易让我分心。\\
Dànshì, wǎngshàng yě yǒu hěn duō shìpín hé yóuxì, róngyì ràng wǒ fēnxīn.\\
« Mais sur Internet, il y a aussi beaucoup de vidéos et de jeux, qui me distraient facilement. »

所以，自律很重要。我每天定计划，只学习不玩手机。\\
Suǒyǐ, zìlǜ hěn zhòngyào. Wǒ měi tiān dìng jìhuà, zhǐ xuéxí bù wán shǒujī.\\
« Donc l'autodiscipline est très importante. Chaque jour, je fixe un programme : je ne fais qu'étudier, je ne joue pas avec mon téléphone. »

总之，互联网是把双刃剑。只要用得好，它就是最好的老师。\\
Zǒngzhī, hùliánwǎng shì bǎ shuāngrènjiàn. Zhǐyào yòng de hǎo, tā jiù shì zuì hǎo de lǎoshī.\\
« En résumé, Internet est une épée à double tranchant. Tant qu'on l'utilise bien, c'est le meilleur des professeurs. »

谢谢大家。\\
Xièxie dàjiā.\\
« Merci à tous. »

\begin{center}
\begin{tabularx}{\linewidth}{|X|X|X|X|}
\hline
\rowcolor{popblueL} Caractères & Pinyin & Nature & Traduction \\
\hline
各位同学 & gèwèi tóngxué & NP & Chers camarades \\
利用 & lìyòng & V & Utiliser \\
复习 & fùxí & V & Réviser \\
方便 & fāngbiàn & Adj & Pratique / commode \\
分心 & fēnxīn & V & Se distraire / perdre sa concentration \\
自律 & zìlǜ & V/Adj & Faire preuve d'autodiscipline \\
双刃剑 & shuāngrènjiàn & N & Épée à double tranchant (image) \\
只要……就…… & zhǐyào… jiù… & Conj & Tant que… alors… \\
\hline
\end{tabularx}
\end{center}

\textbf{Justifications grammaticales :}
\begin{itemize}
    \item \zh{利用互联网学习} : structure V + O + V (le second verbe exprime le but).
    \item \zh{容易让我分心} : \zh{容易} (facile de) + \zh{让} (faire, causatif) + objet + verbe.
    \item \zh{只要用得好，它就是……} : structure conditionnelle \zh{只要……就……} ; \zh{用得好} emploie le complément de degré avec \zh{得}.
    \item Tons : attention à \zh{让} (ràng, 4\up{e} ton) ; \zh{好} (hǎo, 3\up{e} ton) est ici adjectif « bien ».
\end{itemize}
\textbf{Conclusion :} l'exposé respecte le temps, la structure annoncée, utilise le vocabulaire ciblé et des connecteurs variés. La prononciation des termes techniques (\zh{翻译软件}, \zh{自律}) doit être travaillée.
\end{exoresolu}

\begin{methode}[Participer à un débat : exprimer accord, désaccord et nuance]
\textbf{Objectif :} Réagir spontanément aux arguments d'autrui en utilisant les formules fonctionnelles adéquates.
\begin{enumerate}
    \item \textbf{Écouter activement} : repérer le mot-clé de l'argument adverse (ex. : \zh{方便} fāngbiàn).
    \item \textbf{Choisir sa position} : \textbf{accord} $\rightarrow$ \zh{我同意……因为……} (Wǒ tóngyì… yīnwei…) ; \textbf{désaccord poli} $\rightarrow$ \zh{我不同意……我觉得……} (Wǒ bù tóngyì… wǒ juéde…) ; \textbf{nuance} $\rightarrow$ \zh{你说得对，但是……} (Nǐ shuō de duì, dànshì…) / \zh{一方面……另一方面……}.
    \item \textbf{Utiliser les marqueurs de prise de parole} : \zh{我想补充一句} (Wǒ xiǎng bǔchōng yí jù — je voudrais ajouter un mot), \zh{请问} (qǐngwèn — puis-je demander).
    \item \textbf{Appuyer sur un exemple concret} (Cameroun : \zh{流量很贵} liúliàng hěn guì — les données coûtent cher ; Chine : \zh{5G很快} 5G hěn kuài) pour ancrer le débat.
    \item \textbf{Clôturer son tour} : \zh{我就说这么多} (Wǒ jiù shuō zhème duō — j'en reste là), \zh{谢谢} (xièxie).
\end{enumerate}
\cle{Formules clés : 我同意 / 我不同意 / 我想补充 / 举个例子}
\end{methode}

\begin{exoresolu}[Débat : « Internet rapproche-t-il ou éloigne-t-il les gens ? »]
\textbf{Énoncé :} Trois élèves débattent. Complétez les répliques de B et C en choisissant la structure grammaticale et le vocabulaire appropriés (caractères + pinyin). Traduisez.
\begin{itemize}
    \item A : \zh{我认为网络拉近了人与人的距离，微信视频随时能见面。} (Je pense qu'Internet a rapproché les gens : avec la vidéo WeChat, on peut se voir à tout moment.)
    \item B (accord + exemple personnel) : « J'approuve totalement. Ma famille est à Douala, moi à Yaoundé, on se voit chaque semaine par vidéo. »
    \item C (nuance : accord partiel + contre-argument « téléphone à table ») : « Tu as raison, mais parfois on est réunis physiquement et chacun regarde son téléphone, c'est assez ironique. »
    \item A (relance) : \zh{那怎么解决？} (Comment résoudre ça ?)
    \item B (proposition : une règle) : « Je propose qu'on fixe une règle : pas de téléphone pendant les repas. »
\end{itemize}
\tcblower
\textbf{Solution détaillée :}

\textbf{A} : 我认为网络拉近了人与人的距离，微信视频随时能见面。\\
Wǒ rènwéi wǎngluò lājìn le rén yǔ rén de jùlí, Wēixìn shìpín suíshí néng jiànmiàn.\\
« Je pense qu'Internet a réduit la distance entre les personnes : avec la vidéo WeChat, on peut se voir à tout moment. »

\textbf{B} : 我完全同意。比如我家在杜阿拉，我在雅温得，我们每周都能视频见面。\\
Wǒ wánquán tóngyì. Bǐrú wǒ jiā zài Dù'ālā, wǒ zài Yǎwēndé, wǒmen měi zhōu dōu néng shìpín jiànmiàn.\\
« Je suis tout à fait d'accord. Par exemple, ma famille est à Douala et moi à Yaoundé : nous pouvons nous voir en vidéo chaque semaine. »\\
\emph{Analyse :} \zh{完全同意} (approuver totalement). \zh{比如} introduit l'exemple. \zh{都能} : « pouvoir à chaque fois » (\zh{都} renforce la régularité).

\textbf{C} : 你说得对，但是有时候大家聚在一起，却都在看手机，挺讽刺的。\\
Nǐ shuō de duì, dànshì yǒushíhou dàjiā jù zài yìqǐ, què dōu zài kàn shǒujī, tǐng fěngcì de.\\
« Tu as raison, mais parfois tout le monde est réuni, et pourtant chacun regarde son téléphone : c'est assez ironique. »\\
\emph{Analyse :} \zh{你说得对} (validation, avec \zh{得} complément de degré). \zh{但是} (pivot). \zh{聚在一起} (se réunir ; \zh{一起} se prononce yìqǐ). \zh{却} (pourtant, opposition forte). \zh{在看} (action en cours). \zh{挺……的} (assez…).

\textbf{A} : 那怎么解决？\\
Nà zěnme jiějué?\\
« Alors comment résoudre ça ? »

\textbf{B} : 我建议定个规矩：吃饭的时候不许看手机。\\
Wǒ jiànyì dìng ge guīju: chīfàn de shíhou bùxǔ kàn shǒujī.\\
« Je propose de fixer une règle : pendant les repas, interdiction de regarder son téléphone. »\\
\emph{Analyse :} \zh{建议} (suggérer) + V + O. \zh{定规矩} (collocation : établir une règle). \zh{的时候} (au moment où). \zh{不许} (interdiction ferme émanant du groupe).

\textbf{Points de grammaire vérifiés :}
\begin{itemize}
    \item \zh{拉近……距离} (réduire la distance) : verbe + complément de résultat \zh{近} + objet.
    \item \zh{随时} (à tout moment) vs \zh{有时候} (parfois).
    \item \zh{不许} vs \zh{不能} : \zh{不许} implique l'autorité de celui qui parle (règle de groupe), \zh{不能} exprime l'incapacité ou une interdiction générale.
\end{itemize}
\end{exoresolu}

\begin{methode}[Relancer et approfondir l'échange (interaction orale)]
\textbf{Objectif :} Éviter les silences et faire avancer la conversation par des questions de relance ciblées.
\begin{enumerate}
    \item \textbf{Identifier le point d'accroche} dans la phrase du partenaire (un nom, un verbe, un adjectif).
    \item \textbf{Choisir le type de question} :
    \begin{itemize}
        \item \textbf{Précision factuelle} : \zh{具体怎么……？} (jùtǐ zěnme…?) / \zh{能举个例子吗？} (néng jǔ ge lìzi ma?)
        \item \textbf{Opinion / évaluation} : \zh{你觉得……怎么样？} (Nǐ juéde… zěnmeyàng?) / \zh{为什么？} (Wèishénme?)
        \item \textbf{Conséquence / futur} : \zh{那以后呢？} (Nà yǐhòu ne?) / \zh{有什么影响？} (Yǒu shénme yǐngxiǎng?)
        \item \textbf{Comparaison (Cameroun/Chine)} : \zh{在中国也是这样吗？} (Zài Zhōngguó yě shì zhèyàng ma?)
    \end{itemize}
    \item \textbf{Utiliser les particules modales} pour adoucir : \zh{呢} (ne), \zh{吧} (ba, suggestion), \zh{吗} (ma, question fermée).
    \item \textbf{Enchaîner} : reformuler la réponse reçue (\zh{原来是这样，那……} Yuánlái shì zhèyàng, nà… — ah c'est comme ça, alors…) avant de poser la question suivante.
\end{enumerate}
\cle{Technique : privilégier les questions ouvertes (怎么， 为什么， 怎么样) aux questions fermées (吗).}
\end{methode}

\begin{exoresolu}[Entraînement à la relance : « Le commerce en ligne »]
\textbf{Énoncé :} L'élève A dit : \zh{我现在都在网上买衣服，又便宜又快。} (Wǒ xiànzài dōu zài wǎngshàng mǎi yīfu, yòu piányi yòu kuài. — Maintenant j'achète mes vêtements en ligne, c'est à la fois bon marché et rapide.)
Rédigez 3 relances différentes de l'élève B selon les 3 types ci-dessous. Pour chaque relance : caractères, pinyin, traduction, et justification du choix.
\begin{enumerate}
    \item Relance « précision / expérience » (problème de taille/qualité).
    \item Relance « opinion / comparaison » (magasin physique vs en ligne).
    \item Relance « contexte Cameroun » (livraison, paiement Mobile Money).
\end{enumerate}
\tcblower
\textbf{Solution détaillée :}
\begin{enumerate}
    \item \textbf{Précision :}\\
    \zh{网上买衣服，尺码不合适怎么办？退货方便吗？}\\
    Wǎngshàng mǎi yīfu, chǐmǎ bù héshì zěnme bàn? Tuìhuò fāngbiàn ma?\\
    « Quand on achète des vêtements en ligne, si la taille ne convient pas, comment faire ? Le retour est-il facile ? »\\
    \emph{Justification :} cible le point faible implicite (\zh{又便宜又快} cache le risque taille/qualité). Vocabulaire : \zh{尺码} (taille), \zh{退货} (retourner un article). Question ouverte \zh{怎么办} + question fermée \zh{吗}.

    \item \textbf{Opinion / comparaison :}\\
    \zh{那你觉得实体店还有必要去吗？试穿不是更保险吗？}\\
    Nà nǐ juéde shítǐdiàn hái yǒu bìyào qù ma? Shìchuān bú shì gèng bǎoxiǎn ma?\\
    « Alors tu penses qu'il est encore nécessaire d'aller dans les magasins physiques ? Essayer sur place, n'est-ce pas plus sûr ? »\\
    \emph{Justification :} \zh{实体店} (magasin physique) opposé à \zh{网店} (boutique en ligne). \zh{试穿} (essayer un vêtement). Structure rhétorique \zh{不是……吗？} (bú shì… ma) qui attend l'accord de l'interlocuteur.

    \item \textbf{Contexte Cameroun :}\\
    \zh{在喀麦隆网购，物流怎么样？可以用移动支付付款吗？}\\
    Zài Kāmàilóng wǎnggòu, wùliú zěnmeyàng? Kěyǐ yòng yídòng zhīfù fùkuǎn ma?\\
    « Au Cameroun, pour les achats en ligne, comment est la logistique ? Peut-on payer par paiement mobile ? »\\
    \emph{Justification :} ancrage local indispensable. \zh{物流} (logistique/livraison), \zh{移动支付} (paiement mobile, type Mobile Money), \zh{付款} (payer). Question d'opinion \zh{怎么样} + question fermée \zh{吗}.
\end{enumerate}

\textbf{Bilan :} l'élève B démontre une écoute active (il reprend \zh{网上买衣服}), varie les types de questions (fait, opinion, contexte) et utilise le vocabulaire technique (\zh{尺码}, \zh{物流}, \zh{退货}).
\end{exoresolu}

\begin{methode}[Maîtriser la phonétique du débat : tons, neutralisation et rythme]
\textbf{Objectif :} Produire un discours intelligible et naturel en gérant les changements de tons (sandhi) et l'accentuation de phrase.
\begin{enumerate}
    \item \textbf{Repérer les « 3\up{es} tons consécutifs »} : appliquer la règle du sandhi (le 1\up{er} 3\up{e} ton devient un 2\up{e} ton). Ex. : \zh{很好} (hěn hǎo $\rightarrow$ hén hǎo), \zh{可以} (kěyǐ $\rightarrow$ kéyǐ), \zh{老板} (lǎobǎn $\rightarrow$ láobǎn).
    \item \textbf{Gérer \zh{不} (bù) et \zh{一} (yī)} :
    \begin{itemize}
        \item \zh{不} se prononce \textbf{bú} (2\up{e} ton) devant un 4\up{e} ton (\zh{不是} bù shì $\rightarrow$ bú shì) ; il reste \textbf{bù} (4\up{e} ton) ailleurs.
        \item \zh{一} se prononce \textbf{yí} (2\up{e} ton) devant un 4\up{e} ton (\zh{一样} yīyàng $\rightarrow$ yíyàng) ; \textbf{yì} (4\up{e} ton) devant un 1\up{er}, 2\up{e} ou 3\up{e} ton (\zh{一起} yìqǐ) ; \textbf{yī} (1\up{er} ton) quand il est un chiffre ou en fin d'énumération.
    \end{itemize}
    \item \textbf{Neutraliser les particules et suffixes} : \zh{了} (le), \zh{着} (zhe), \zh{过} (guo), \zh{们} (men), \zh{子} (zi), \zh{的/得/地} (de) $\rightarrow$ ton neutre (léger, court).
    \item \textbf{Rythme binaire (groupes de 2 ou 4 syllabes)} : \zh{互联网 | 很方便} ; marquer une pause (//) aux frontières syntaxiques (sujet // prédicat).
    \item \textbf{Accentuer l'information nouvelle (focus)} : hausse d'intensité + léger allongement sur le mot-clé (ex. : dans \zh{我反对这个观点} Wǒ fǎnduì zhège guāndiǎn, on accentue \zh{反对}).
\end{enumerate}
\cle{Règle d'or : « les tons des mots isolés ne sont pas toujours les tons dans la phrase ». Entraînez-vous sur des phrases entières.}
\end{methode}

\begin{exoresolu}[Phonétique : transcription tonale et lecture d'une phrase de débat]
\textbf{Énoncé :} Pour la phrase ci-dessous :
\begin{enumerate}
    \item Écrivez le pinyin \textbf{avec les tons réellement prononcés} (après sandhi et neutralisation).
    \item Marquez les pauses (//) et soulignez le mot accentué (focus).
    \item Identifiez 2 cas de ton neutre et 1 cas de sandhi.
\end{enumerate}
\textbf{Phrase :} \zh{我不太同意你的看法，因为网络虽然方便，但也让人很孤独。}
\tcblower
\textbf{Solution détaillée :}

\textbf{1. Pinyin corrigé (tons réels) :}\\
Wǒ \textbf{bú} tài tóngyì nǐ \textbf{de} kànfa, // yīnwei wǎngluò suīrán fāngbiàn, // dàn yě ràng rén hěn gūdú.\\
\emph{Note :} \zh{不} $\rightarrow$ \textbf{bú} (2\up{e} ton) car suivi de \zh{太} (tài, 4\up{e} ton). \zh{的} est au ton neutre (de). \zh{看法} se prononce kànfa : la seconde syllabe \zh{法} s'allège en ton neutre dans ce mot composé.

\textbf{2. Pauses // et focus :}\\
Wǒ bú tài tóngyì nǐ de kànfa, // yīnwei wǎngluò suīrán fāngbiàn, // dàn yě ràng rén hěn \textbf{gūdú}.\\
\emph{Justification :} pause après la première proposition (\zh{看法}) ; pause après la concession (\zh{虽然方便}) ; focus final sur \zh{孤独} (gūdú, argument clé et information nouvelle).

\textbf{3. Identification demandée :}
\begin{itemize}
    \item \textbf{Ton neutre 1 :} \zh{的} (de) dans \zh{你的看法} (particule structurelle).
    \item \textbf{Ton neutre 2 :} \zh{法} dans \zh{看法} (kànfa) : seconde syllabe de mot composé, prononcée légère.
    \item \textbf{Sandhi :} \zh{不} (bù) $\rightarrow$ \textbf{bú} devant \zh{太} (tài, 4\up{e} ton). C'est le sandhi obligatoire de \zh{不}. (Il n'y a pas ici de suite 3\up{e} + 3\up{e} ton du type \zh{很好}.)
\end{itemize}

\textbf{Conseil pédagogique :} faire lire en chœur en tapant le rythme (main sur la table) pour ancrer les groupes rythmiques. Insister sur le contraste \zh{bú tài} (montée) / \zh{bù tóng} (descente).
\end{exoresolu}

\begin{methode}[Préparer et réussir un jeu de rôle (saynète) noté]
\textbf{Objectif :} Structurer une interaction à deux (A/B) avec répartition des rôles, gestion du temps et couverture des fonctions langagières exigées.
\begin{enumerate}
    \item \textbf{Lire la consigne} : identifier le rôle (A = client / B = vendeur ; A = parent / B = adolescent), le lieu, le but (convaincre, se plaindre, décider).
    \item \textbf{Lister les « points obligatoires »} (checklist) :
    \begin{itemize}
        \item Salutations / prise de congé formelles (\zh{您好} nín hǎo / \zh{再见} zàijiàn).
        \item 3 arguments minimum par rôle.
        \item 1 structure complexe visée (ex. : phrase en \zh{比}, phrase en \zh{如果……就……}).
        \item 2 formules de politesse / relance (\zh{请问} qǐngwèn / \zh{没关系} méi guānxi).
    \end{itemize}
    \item \textbf{Rédiger un « squelette dialogué »} (5--6 répliques chacun) au brouillon : réplique A $\rightarrow$ réplique B $\rightarrow$ réplique A… Vérifier l'enchaînement logique.
    \item \textbf{Répétition chronométrée} : 2--3 minutes. Vérifier : tons (noms propres, mots-clés), fluidité (pas de « euh » français), volume, contact visuel.
    \item \textbf{Stratégie « sauvetage »} : en cas de trou de mémoire $\rightarrow$ \zh{对不起，我想一想……} (Duìbuqǐ, wǒ xiǎng yi xiǎng… — pardon, je réfléchis…) / \zh{怎么说呢……} (Zěnme shuō ne… — comment dire…) / reformuler plus simplement.
\end{enumerate}
\cle{Évaluation indicative : contenu (50\%) + langue (30\%) + phonétique et interaction (20\%).}
\end{methode}

\begin{exoresolu}[Jeu de rôle : « Achat d'un smartphone : négocier le prix »]
\textbf{Énoncé :} \textbf{Rôle A (client) :} vous voulez acheter un téléphone \zh{华为} (Huáwéi) dernier modèle. Budget maximum : 4000 元. Vous trouvez le modèle exposé un peu cher. Négociez le prix, une coque et un film protecteur.
\textbf{Rôle B (vendeur) :} modèle exposé à 4500 元. Vous pouvez offrir le film protecteur. La coque coûte 50 元. Pas de baisse possible sur le téléphone, mais vous proposez l'ancien modèle à 3800 元.
\textbf{Consigne :} rédigez le squelette dialogué (6 répliques chacun) en caractères + pinyin. Intégrez : la phrase en \zh{比}, \zh{如果……就……}, \zh{还是}, \zh{划算}.
\tcblower
\textbf{Solution détaillée (dialogue modèle) :}

\textbf{A} : 你好，这款华为手机多少钱？\\
Nǐ hǎo, zhè kuǎn Huáwéi shǒujī duōshao qián?\\
« Bonjour, combien coûte ce modèle de Huawei ? »

\textbf{B} : 你好！这是最新款，4500元。屏幕很大，5G特别快。\\
Nǐ hǎo! Zhè shì zuì xīn kuǎn, sìqiān wǔbǎi yuán. Píngmù hěn dà, 5G tèbié kuài.\\
« Bonjour ! C'est le tout dernier modèle, 4500 yuans. L'écran est grand et la 5G est très rapide. »

\textbf{A} : 4500元有点贵。那个旧款呢？如果买旧款，就送耳机吗？\\
Sìqiān wǔbǎi yuán yǒudiǎn guì. Nà ge jiù kuǎn ne? Rúguǒ mǎi jiù kuǎn, jiù sòng ěrjī ma?\\
« 4500 yuans, c'est un peu cher. Et l'ancien modèle ? Si j'achète l'ancien modèle, vous offrez les écouteurs ? »\\
\emph{Grammaire :} \zh{有点} + adjectif (un peu trop). \zh{如果……就……} (hypothèse). \zh{送} (offrir).

\textbf{B} : 旧款3800元，屏幕比新款小，而且没有货了。不过，如果您今天定新款，可以送钢化膜。\\
Jiù kuǎn sānqiān bābǎi yuán, píngmù bǐ xīn kuǎn xiǎo, érqiě méiyǒu huò le. Búguò, rúguǒ nín jīntiān dìng xīn kuǎn, kěyǐ sòng gānghuà mó.\\
« L'ancien modèle coûte 3800 yuans, son écran est plus petit que celui du nouveau, et en plus il n'y en a plus en stock. En revanche, si vous prenez le nouveau aujourd'hui, je peux offrir le film trempé. »\\
\emph{Grammaire :} phrase en \zh{比} : \zh{屏幕比新款小} (A + \zh{比} + B + adjectif). \zh{不过} (transition). \zh{钢化膜} (film trempé).

\textbf{A} : 只送膜还不够。我的预算只有4000元。要不您少收200元，我再要一个保护套？\\
Zhǐ sòng mó hái bùgòu. Wǒ de yùsuàn zhǐyǒu sìqiān yuán. Yàobu nín shǎo shōu èrbǎi yuán, wǒ zài yào yí ge bǎohùtào?\\
« Offrir seulement le film ne suffit pas. Mon budget n'est que de 4000 yuans. Sinon, vous me faites 200 yuans de moins et je prends en plus une coque ? »\\
\emph{Vocabulaire :} \zh{预算} (budget). \zh{要不} (sinon, proposition). \zh{少收} (facturer moins). \zh{保护套} (coque de protection). Classificateur : \zh{一个保护套}.

\textbf{B} : 手机的价格不能动。但是保护套原价50元，给您打八折，40元。这样一共4540元，很划算吧？\\
Shǒujī de jiàgé bù néng dòng. Dànshì bǎohùtào yuánjià wǔshí yuán, gěi nín dǎ bāzhé, sìshí yuán. Zhèyàng yígòng sìqiān wǔbǎi sìshí yuán, hěn huásuàn ba?\\
« Le prix du téléphone ne peut pas bouger. Mais la coque, au lieu de 50 yuans, je vous fais une remise de 20\% : 40 yuans. Cela fait 4540 yuans au total, c'est une bonne affaire, non ? »\\
\emph{Grammaire :} \zh{不能} (interdiction liée à une règle). \zh{打八折} (dǎ bāzhé) : littéralement « faire 80\% du prix », donc 20\% de remise. \zh{划算} (avantageux, qui vaut le coup). \zh{吧} adoucit et invite à l'accord.

\textbf{A} : 行，那就买新款吧。刷卡还是扫码？\\
Xíng, nà jiù mǎi xīn kuǎn ba. Shuākǎ háishì sǎomǎ?\\
« D'accord, alors je prends le nouveau modèle. Carte bancaire ou scan ? »\\
\emph{Clôture :} \zh{行} (OK). \zh{那就} (alors, dans ce cas). \zh{还是} (alternative dans une question).

\textbf{B} : 扫码快。微信、支付宝都可以。谢谢惠顾，慢走！\\
Sǎomǎ kuài. Wēixìn, Zhīfùbǎo dōu kěyǐ. Xièxie huìgù, màn zǒu!\\
« Le scan est rapide. WeChat et Alipay marchent tous les deux. Merci de votre visite, au revoir ! »

\textbf{Analyse des objectifs atteints :}
\begin{itemize}
    \item Structures ciblées : phrase en \zh{比} (réplique B), \zh{如果……就……} (A et B), \zh{还是} (A).
    \item Vocabulaire spécialisé : \zh{钢化膜} (film trempé), \zh{保护套} (coque), \zh{打折} (faire une remise), \zh{扫码} (scanner), \zh{预算} (budget).
    \item Interaction : négociation réelle (concession sur le film $\rightarrow$ demande de la coque $\rightarrow$ remise sur la coque $\rightarrow$ accord).
    \item Tons : \zh{便宜} (piányi), \zh{贵} (guì), \zh{划算} (huásuàn, 2\up{e}--4\up{e} ton), \zh{华为} (Huáwéi, 2\up{e}--2\up{e} ton).
\end{itemize}
\end{exoresolu}

\begin{methode}[Réviser et activer le vocabulaire thématique (carte mentale $\rightarrow$ parole)]
\textbf{Objectif :} Passer de la reconnaissance passive à la production active du lexique « Internet et télécommunications ».
\begin{enumerate}
    \item \textbf{Carte mentale chronométrée (3 min)} : au centre \zh{互联网}. 4 branches : \zh{设备} (appareils : \zh{手机}, \zh{电脑}), \zh{动作} (actions : \zh{上网}, \zh{刷视频}, \zh{网购}, \zh{视频通话}), \zh{概念} (concepts : \zh{流量}, \zh{密码}, \zh{账号}), \zh{评价} (évaluation : \zh{方便}, \zh{上瘾}, \zh{浪费时间}).
    \item \textbf{Classement par collocations (blocs lexicaux)} : ne pas apprendre \zh{流量} seul, mais \zh{流量不够用} (les données ne suffisent pas), \zh{充流量} (recharger des données). \zh{密码} $\rightarrow$ \zh{设置密码} (créer un mot de passe), \zh{忘记密码} (oublier son mot de passe).
    \item \textbf{Jeu des « 3 mots »} : tirer 3 mots au hasard $\rightarrow$ construire une phrase cohérente les reliant. Ex. : \zh{网购} + \zh{方便} + \zh{便宜} $\rightarrow$ \zh{网购又方便又便宜。} (Wǎnggòu yòu fāngbiàn yòu piányi.)
    \item \textbf{Traduction inversée (français $\rightarrow$ chinois)} : 10 phrases types du programme $\rightarrow$ écrire en caractères + pinyin sans regarder le cours.
    \item \textbf{Cartes mémoire orales} : recto = français / verso = caractères + pinyin. Révision espacée (J0, J1, J3, J7).
\end{enumerate}
\cle{Blocs lexicaux > mots isolés. Contexte camerounais : \zh{流量包} (forfait de données), \zh{充值卡} (carte de recharge).}
\end{methode}

\begin{exoresolu}[Activation lexicale : de la carte mentale à l'argumentaire]
\textbf{Énoncé :} À partir de la branche « évaluation / problèmes », rédigez un paragraphe argumenté de 80--100 caractères sur : \zh{短视频对青少年的影响} (Duǎnshìpín duì qīngshàonián de yǐngxiǎng — L'impact des vidéos courtes sur les adolescents). Structures imposées : \zh{虽然……但是……}, \zh{不仅……而且……}, \zh{建议}. Utilisez au moins 6 mots de la liste : \zh{上瘾}, \zh{视力}, \zh{知识}, \zh{放松}, \zh{自控力}, \zh{引导}, \zh{浪费时间}.
\tcblower
\textbf{Solution détaillée (rédaction modèle) :}

短视频虽然能让青少年放松，还能学到一些知识，但是很容易让人上瘾，伤害视力。不仅浪费时间，而且降低自控力。建议家长和学校加强引导，规定每天用手机的时间，让孩子多读书、多运动。\\
Duǎnshìpín suīrán néng ràng qīngshàonián fàngsōng, hái néng xuédào yìxiē zhīshi, dànshì hěn róngyì ràng rén shàngyǐn, shānghài shìlì. Bùjǐn làngfèi shíjiān, érqiě jiàngdī zìkònglì. Jiànyì jiāzhǎng hé xuéxiào jiāqiáng yǐndǎo, guīdìng měi tiān yòng shǒujī de shíjiān, ràng háizi duō dúshū, duō yùndòng.\\
« Les vidéos courtes, bien qu'elles permettent aux adolescents de se détendre et d'apprendre quelques connaissances, rendent très facilement accro et abîment la vue. Non seulement elles font perdre du temps, mais en plus elles diminuent la maîtrise de soi. Je suggère que les parents et l'école renforcent l'encadrement, fixent le temps d'utilisation quotidienne du téléphone, et fassent lire et faire du sport davantage aux enfants. »

\begin{center}
\begin{tabularx}{\linewidth}{|X|X|X|X|}
\hline
\rowcolor{popblueL} Caractères & Pinyin & Nature & Traduction \\
\hline
短视频 & duǎnshìpín & N & Vidéo courte (type TikTok/Douyin) \\
青少年 & qīngshàonián & N & Adolescent / jeune \\
放松 & fàngsōng & V & Se détendre \\
上瘾 & shàngyǐn & V & Devenir accro \\
视力 & shìlì & N & Vue / acuité visuelle \\
浪费时间 & làngfèi shíjiān & Loc. V & Perdre du temps \\
自控力 & zìkònglì & N & Maîtrise de soi / volonté \\
引导 & yǐndǎo & V/N & Guider / encadrement \\
规定 & guīdìng & V & Fixer une règle / stipuler \\
\hline
\end{tabularx}
\end{center}

\textbf{Analyse structurelle :}
\begin{itemize}
    \item \zh{虽然……但是……} : concession standard (HSK 3). L'ordre des deux propositions est fixe.
    \item \zh{不仅……而且……} : progression (HSK 4). Le sujet commun est sous-entendu ; \zh{不仅} précède le premier verbe (\zh{浪费}), \zh{而且} introduit le second (\zh{降低}).
    \item \zh{建议} + sujet + verbe : structure de conseil. Ici \zh{建议家长和学校加强引导}.
    \item \textbf{Tons :} \zh{自控力} (zì kòng lì, 4-4-4) ; \zh{引导} (yǐndǎo, 3\up{e} + 3\up{e} ton $\rightarrow$ prononcé yíndǎo par sandhi) ; \zh{可以} type de mots à surveiller dans toute production orale.
\end{itemize}
\textbf{Note :} le texte atteint environ 95 caractères. Il est dense, cohérent, utilise les mots imposés et les structures cibles.
\end{exoresolu}

\begin{attention}[Les 5 erreurs qui coûtent des points à l'oral et à l'écrit]
\begin{enumerate}
    \item \textbf{Tons : « le 3\up{e} ton qui ne change pas ».} \zh{很好} lu hěn hǎo (3-3) au lieu de hén hǎo (2-3) ; \zh{可以} lu kěyǐ au lieu de kéyǐ. \textbf{Règle :} 3\up{e} ton + 3\up{e} ton $\rightarrow$ 2\up{e} ton + 3\up{e} ton. Systématiser sur les mots très fréquents (\zh{我}, \zh{你}, \zh{好}, \zh{很}, \zh{也}, \zh{都}).
    \item \textbf{Ordre des mots : « complément de temps ou de lieu après le verbe » (calque du français).} *\zh{我去图书馆昨天} est incorrect. \textbf{Correct :} \zh{我昨天去图书馆。} (Wǒ zuótiān qù túshūguǎn. — Sujet + temps + verbe) ou \zh{昨天我去图书馆。} Le complément de temps se place avant le verbe ; le lieu avec \zh{在} aussi : \zh{我在图书馆学习。} (Wǒ zài túshūguǎn xuéxí.)
    \item \textbf{Mots-outils \zh{的/得/地} : confusion systématique.} \zh{安静的房间} (adjectif + \zh{的} + nom), \zh{跑得快} (verbe + \zh{得} + complément de degré), \zh{安静地走} (adverbe + \zh{地} + verbe). \textbf{Astuce :} \zh{的} annonce un nom, \zh{得} suit un verbe, \zh{地} précède un verbe.
    \item \textbf{Classificateurs oubliés ou erronés.} \zh{一个手机} est compréhensible mais \zh{一部手机} est meilleur ; *\zh{两个衣服} est incorrect $\rightarrow$ \zh{两件衣服}. \textbf{Liste à connaître :} \zh{个} (générique), \zh{本} (livres), \zh{张} (objets plats : papier, billet, table), \zh{条} (objets longs : pantalon, rue), \zh{件} (vêtements, affaires), \zh{部/台} (appareils), \zh{辆} (véhicules), \zh{位} (personnes, poli).
    \item \textbf{Structures \zh{比} / \zh{没有} : inversion du sens.} \zh{A 比 B + Adj} signifie « A est plus Adj que B » : \zh{我比他高。} (Wǒ bǐ tā gāo. — Je suis plus grand que lui.) \zh{A 没有 B + Adj} signifie « A n'est pas aussi Adj que B » : \zh{我没有他高。} (Wǒ méiyǒu tā gāo. — Je ne suis pas aussi grand que lui.) \textbf{Piège :} ne jamais ajouter \zh{很} ou \zh{更} dans une phrase en \zh{比} : *\zh{我比他很高} est incorrect ; on dit \zh{我比他高} ou, pour insister, \zh{我比他高得多} (Wǒ bǐ tā gāo de duō. — Je suis bien plus grand que lui.)
\end{enumerate}
\cle{Check-list avant de parler : tons (3\up{e} ton, 不, 一) -- ordre des mots (temps/lieu avant le verbe) -- 的/得/地 -- classificateurs -- 比/没有.}
\end{attention}

\newpage

\coursec{Exercices}

\courssub{Je vérifie mes connaissances}

\begin{exercice}[QCM : le vocabulaire d'internet]
Entoure la bonne réponse.

\begin{enumerate}
\item \zh{上网} (shàngwǎng) signifie :
\begin{enumerate}
\item regarder la télévision
\item se connecter à internet
\item écrire une lettre
\end{enumerate}
\item \zh{网站} (wǎngzhàn) signifie :
\begin{enumerate}
\item un site web
\item un téléphone portable
\item un mot de passe
\end{enumerate}
\item Pour dire « envoyer un courriel », on dit :
\begin{enumerate}
\item \zh{发电子邮件} (fā diànzǐ yóujiàn)
\item \zh{打电话} (dǎ diànhuà)
\item \zh{买东西} (mǎi dōngxi)
\end{enumerate}
\item \zh{我觉得} (wǒ juéde) sert à :
\begin{enumerate}
\item poser une question
\item donner son avis
\item dire au revoir
\end{enumerate}
\end{enumerate}
\end{exercice}

\begin{exercice}[Vrai ou faux ? Justifie ta réponse]
Réponds par V (vrai) ou F (faux) et justifie en une phrase en français.

\begin{enumerate}
\item \zh{网络} (wǎngluò) signifie « le réseau, internet ».
\item La structure \zh{因为……所以……} (yīnwei... suǒyǐ...) exprime la cause et la conséquence.
\item Pour demander à quelqu'un son avis, on peut dire \zh{你呢？} (nǐ ne ?).
\item \zh{好处} (hǎochu) signifie « inconvénient ».
\item La question alternative se forme avec \zh{还是} (háishi).
\end{enumerate}
\end{exercice}

\begin{exercice}[Trouve le mot chinois]
Complète chaque phrase française par le mot chinois correct (caractères + pinyin) choisi dans la liste : \zh{聊天} (liáotiān), \zh{信息} (xìnxī), \zh{游戏} (yóuxì), \zh{新闻} (xīnwén), \zh{学习} (xuéxí).

\begin{enumerate}
\item Sur internet, je cherche des \_\_\_\_\_\_\_\_ (informations) pour mes devoirs.
\item Le soir, mon frère joue à des \_\_\_\_\_\_\_\_ (jeux) en ligne.
\item Ma mère lit les \_\_\_\_\_\_\_\_ (nouvelles) du Cameroun sur son téléphone.
\item J'aime \_\_\_\_\_\_\_\_ (bavarder) avec mes amis sur les réseaux sociaux.
\item Internet nous aide à \_\_\_\_\_\_\_\_ (étudier) le chinois.
\end{enumerate}
\end{exercice}

\begin{exercice}[Associe le pinyin aux caractères]
Relie chaque caractère ou mot à son pinyin. Attention aux tons !

\begin{center}
\begin{tabularx}{\linewidth}{|X|X|}
\hline
\rowcolor{popblueL} Caractères & Pinyin \\
\hline
1. 网 & a. shǒujī \\
\hline
2. 手机 & b. wǎng \\
\hline
3. 电脑 & c. hěn yǒuyòng \\
\hline
4. 很有用 & d. diànnǎo \\
\hline
5. 同意 & e. tóngyì \\
\hline
\end{tabularx}
\end{center}
\end{exercice}

\courssub{Je m'entraîne}

\begin{exercice}[Traduis en chinois]
Traduis les phrases suivantes en chinois (caractères + pinyin).

\begin{enumerate}
\item J'utilise internet tous les jours.
\item Internet est très utile pour nos études.
\item Je pense que le téléphone portable est plus pratique que l'ordinateur.
\item Parce qu'internet est rapide, nous aimons l'utiliser.
\item Et toi, qu'est-ce que tu fais sur internet ?
\end{enumerate}
\end{exercice}

\begin{exercice}[Transforme en question alternative avec 还是]
Transforme chaque paire de phrases en une question alternative avec \zh{还是} (háishi), puis traduis ta question en français.

Exemple résolu : \zh{我用手机。} / \zh{我用电脑。} → \zh{你用手机还是用电脑？} (Nǐ yòng shǒujī háishi yòng diànnǎo ?) « Tu utilises le téléphone ou l'ordinateur ? »

\begin{enumerate}
\item \zh{我在网上学习。} / \zh{我在网上玩游戏。}
\item \zh{互联网有好处。} / \zh{互联网有坏处。}
\item \zh{你早上上网。} / \zh{你晚上上网。}
\item \zh{他用电脑发邮件。} / \zh{他用手机发邮件。}
\end{enumerate}
\end{exercice}

\begin{exercice}[Dialogue à trous : les formules fonctionnelles]
Complète le dialogue suivant avec les formules de la liste : \zh{我觉得} (wǒ juéde), \zh{我同意} (wǒ tóngyì), \zh{你呢} (nǐ ne), \zh{因为} (yīnwei), \zh{所以} (suǒyǐ).

\begin{textebox}[Dialogue : 互联网好不好？]
甲：\_\_\_\_\_\_\_\_ 互联网很有用，我们可以学习很多东西。\\
(Jiǎ : \_\_\_\_\_\_\_\_ hùliánwǎng hěn yǒuyòng, wǒmen kěyǐ xuéxí hěn duō dōngxi.)\\
乙：\_\_\_\_\_\_\_\_ ！\_\_\_\_\_\_\_\_ 网上有很多信息，\_\_\_\_\_\_\_\_ 学习很方便。\\
(Yǐ : \_\_\_\_\_\_\_\_ ! \_\_\_\_\_\_\_\_ wǎngshang yǒu hěn duō xìnxī, \_\_\_\_\_\_\_\_ xuéxí hěn fāngbiàn.)\\
甲：可是上网时间太长也不好。\_\_\_\_\_\_\_\_ ？\\
(Kěshì shàngwǎng shíjiān tài cháng yě bù hǎo. \_\_\_\_\_\_\_\_ ?)\\
乙：对，我们要少玩游戏，多学习。\\
(Duì, wǒmen yào shǎo wán yóuxì, duō xuéxí.)
\end{textebox}

Traduction du dialogue complet à rédiger ensuite en français.
\end{exercice}

\begin{exercice}[Remets les mots dans l'ordre]
Remets les mots dans le bon ordre pour former des phrases correctes, écris-les en caractères avec le pinyin, puis traduis-les.

\begin{enumerate}
\item 我 / 每天 / 上网 / 都
\item 觉得 / 互联网 / 我 / 有用 / 很
\item 因为 / 方便 / 很 / 所以 / 喜欢 / 大家 / 都 / 它
\item 手机 / 用 / 聊天 / 朋友们 / 和 / 我
\end{enumerate}
\end{exercice}

\courssub{Je me prépare au Bac}

\begin{exercice}[Compréhension de l'écrit et questions de langue]
Lis le texte suivant, puis réponds aux questions.

\begin{textebox}[互联网和喀麦隆的学生]
现在，很多喀麦隆的学生每天都上网。他们用手机或者电脑学习、聊天、看新闻。\\
(Xiànzài, hěn duō Kāmàilóng de xuésheng měitiān dōu shàngwǎng. Tāmen yòng shǒujī huòzhě diànnǎo xuéxí, liáotiān, kàn xīnwén.)\\
我觉得互联网很有用，因为网上有很多信息，所以学习很方便。\\
(Wǒ juéde hùliánwǎng hěn yǒuyòng, yīnwei wǎngshang yǒu hěn duō xìnxī, suǒyǐ xuéxí hěn fāngbiàn.)\\
可是，有些学生玩游戏的时间太长，这不好。我们应该少玩游戏，多学习。\\
(Kěshì, yǒuxiē xuésheng wán yóuxì de shíjiān tài cháng, zhè bù hǎo. Wǒmen yīnggāi shǎo wán yóuxì, duō xuéxí.)\\
总之，互联网是好东西，但是我们要好好用它。\\
(Zǒngzhī, hùliánwǎng shì hǎo dōngxi, dànshì wǒmen yào hǎohǎo yòng tā.)
\end{textebox}

\textbf{Questions de compréhension} (réponds en français) :
\begin{enumerate}
\item Que font les élèves camerounais sur internet ? Cite trois activités.
\item Pourquoi l'auteur trouve-t-il internet utile ?
\item Quel problème l'auteur signale-t-il ?
\item Quel conseil l'auteur donne-t-il à la fin du texte ?
\end{enumerate}

\textbf{Questions de langue} :
\begin{enumerate}
\setcounter{enumi}{4}
\item Relève dans le texte une phrase avec la structure \zh{因为……所以……} et traduis-la.
\item Trouve dans le texte le mot qui signifie « en résumé, bref » et donne son pinyin.
\item Transforme en question alternative : \zh{他们用手机。} / \zh{他们用电脑。}
\end{enumerate}
\end{exercice}

\begin{exercice}[Thème et version]
\textbf{Thème.} Traduis en chinois (caractères + pinyin) :
\begin{enumerate}
\item Chaque jour, j'utilise mon téléphone portable pour étudier et bavarder avec mes amis.
\item Je pense qu'internet a des avantages et des inconvénients.
\item Parce qu'internet est très rapide, nous pouvons trouver beaucoup d'informations.
\end{enumerate}

\textbf{Version.} Traduis en français :
\begin{enumerate}
\item \zh{网络很方便，我们可以在家学习。}(Wǎngluò hěn fāngbiàn, wǒmen kěyǐ zài jiā xuéxí.)
\item \zh{我不同意你的看法，我觉得上网时间不能太长。}(Wǒ bù tóngyì nǐ de kànfǎ, wǒ juéde shàngwǎng shíjiān bù néng tài cháng.)
\item \zh{你喜欢用手机还是电脑？}(Nǐ xǐhuan yòng shǒujī háishi diànnǎo ?)
\end{enumerate}
\end{exercice}

\begin{exercice}[Expression écrite et préparation à l'oral]
\textbf{Sujet d'expression écrite.} Rédige en chinois un court texte de \textbf{60 à 80 caractères} sur le sujet : \zh{互联网和我们的生活} (Hùliánwǎng hé wǒmen de shēnghuó) « Internet et notre vie ». Utilise obligatoirement la structure \zh{第一……第二……总之……} (dì-yī... dì-èr... zǒngzhī...) et au moins une phrase avec \zh{因为……所以……}.

\textbf{Préparation à l'oral (exposé).} Prépare un exposé oral de 5 phrases sur le même sujet. Tu devras :
\begin{itemize}
\item commencer par une phrase de présentation : \zh{我给大家介绍一下……} (Wǒ gěi dàjiā jièshào yīxià...) ;
\item donner ton avis avec \zh{我觉得……} ;
\item citer un avantage et un inconvénient ;
\item conclure avec \zh{总之……}.
\end{itemize}

\textbf{Jeu de rôle évalué (en binôme).} Interviewe un camarade sur son usage d'internet : minimum 4 échanges, avec au moins 2 relances (\zh{你呢？}, \zh{为什么？}, \zh{真的吗？}) et une question alternative avec \zh{还是}. La prononciation des tons sera notée.
\end{exercice}

\newpage

\coursec{Corrigés des exercices}

\begin{corrige}[Exercice 1 — QCM : le vocabulaire d'internet]
\begin{enumerate}
\item \textbf{b} — \zh{上网} (shàngwǎng) signifie « se connecter à internet ». \zh{上} (shàng) exprime ici l'idée de « monter sur, accéder à » et \zh{网} (wǎng) signifie « le filet, le réseau ».
\item \textbf{a} — \zh{网站} (wǎngzhàn) signifie « un site web » : \zh{网} (wǎng) « réseau » + \zh{站} (zhàn) « station ».
\item \textbf{a} — \zh{发电子邮件} (fā diànzǐ yóujiàn) signifie « envoyer un courriel ». \zh{打电话} (dǎ diànhuà) = « téléphoner » ; \zh{买东西} (mǎi dōngxi) = « faire des achats ».
\item \textbf{b} — \zh{我觉得} (wǒ juéde) sert à donner son avis : « je pense que, je trouve que ». C'est une formule fonctionnelle essentielle du débat.
\end{enumerate}
\end{corrige}

\begin{corrige}[Exercice 2 — Vrai ou faux ?]
\begin{enumerate}
\item \textbf{V} — \zh{网络} (wǎngluò) signifie bien « le réseau, internet » : \zh{网} « filet » + \zh{络} « relier ».
\item \textbf{V} — \zh{因为……所以……} (yīnwei... suǒyǐ...) exprime la cause puis la conséquence : « parce que... donc... ». Exemple : \zh{因为互联网很方便，所以我很喜欢。} (Yīnwei hùliánwǎng hěn fāngbiàn, suǒyǐ wǒ hěn xǐhuan.) « Parce qu'internet est pratique, je l'aime beaucoup. »
\item \textbf{V} — \zh{你呢？} (nǐ ne ?) est une relance qui signifie « et toi ? » ; elle permet de demander l'avis de l'autre.
\item \textbf{F} — \zh{好处} (hǎochu) signifie « avantage » ; « inconvénient » se dit \zh{坏处} (huàichu). Attention à l'antonymie \zh{好}/\zh{坏}.
\item \textbf{V} — La question alternative se forme avec \zh{还是} (háishi) : \zh{A 还是 B ?} = « A ou B ? ». Exemple : \zh{你用手机还是电脑？} (Nǐ yòng shǒujī háishi diànnǎo ?)
\end{enumerate}
\end{corrige}

\begin{corrige}[Exercice 3 — Trouve le mot chinois]
\begin{enumerate}
\item informations → \zh{信息} (xìnxī) : « Sur internet, je cherche des informations pour mes devoirs. » \zh{我在网上找信息。} (Wǒ zài wǎngshang zhǎo xìnxī.)
\item jeux → \zh{游戏} (yóuxì) : \zh{晚上，我哥哥玩网上游戏。} (Wǎnshang, wǒ gēge wán wǎngshang yóuxì.)
\item nouvelles → \zh{新闻} (xīnwén) : \zh{妈妈用手机看喀麦隆的新闻。} (Māma yòng shǒujī kàn Kāmàilóng de xīnwén.)
\item bavarder → \zh{聊天} (liáotiān) : \zh{我喜欢和朋友在网上聊天。} (Wǒ xǐhuan hé péngyou zài wǎngshang liáotiān.)
\item étudier → \zh{学习} (xuéxí) : \zh{互联网帮助我们学习汉语。} (Hùliánwǎng bāngzhù wǒmen xuéxí Hànyǔ.)
\end{enumerate}
\textbf{Point de vigilance :} ne pas confondre \zh{信息} (xìnxī, 4e + 1er ton) et \zh{新闻} (xīnwén, 1er + 2e ton) : les tons changent le sens.
\end{corrige}

\begin{corrige}[Exercice 4 — Associe le pinyin aux caractères]
\begin{center}
\begin{tabularx}{\linewidth}{|X|X|X|}
\hline
\rowcolor{popblueL} Caractères & Pinyin & Traduction \\
\hline
1. 网 → b & wǎng & le réseau \\
\hline
2. 手机 → a & shǒujī & le téléphone portable \\
\hline
3. 电脑 → d & diànnǎo & l'ordinateur \\
\hline
4. 很有用 → c & hěn yǒuyòng & très utile \\
\hline
5. 同意 → e & tóngyì & être d'accord \\
\hline
\end{tabularx}
\end{center}
\textbf{Point de vigilance :} \zh{手机} (shǒujī, 3e + 1er ton) : « main-machine » ; \zh{电脑} (diànnǎo, 4e + 3e ton) : « électricité-cerveau ». Ces images aident à mémoriser.
\end{corrige}

\begin{corrige}[Exercice 5 — Traduis en chinois]
\begin{enumerate}
\item \zh{我每天用互联网。} (Wǒ měitiān yòng hùliánwǎng.) — Le complément de temps \zh{每天} se place avant le verbe.
\item \zh{互联网对我们的学习很有用。} (Hùliánwǎng duì wǒmen de xuéxí hěn yǒuyòng.) — \zh{对……很有用} = « très utile pour... ».
\item \zh{我觉得手机比电脑方便。} (Wǒ juéde shǒujī bǐ diànnǎo fāngbiàn.) — Structure comparative : A + \zh{比} + B + adjectif.
\item \zh{因为互联网很快，所以我们喜欢用。} (Yīnwei hùliánwǎng hěn kuài, suǒyǐ wǒmen xǐhuan yòng.) — En chinois, on peut garder \zh{因为} et \zh{所以} dans la même phrase (contrairement au français).
\item \zh{你呢？你在网上做什么？} (Nǐ ne ? Nǐ zài wǎngshang zuò shénme ?) — \zh{什么} (shénme) « quoi » reste à la place du complément, pas en tête de phrase.
\end{enumerate}
\end{corrige}

\begin{corrige}[Exercice 6 — Transforme en question alternative avec 还是]
\begin{enumerate}
\item \zh{你在网上学习还是玩游戏？} (Nǐ zài wǎngshang xuéxí háishi wán yóuxì ?) « Sur internet, tu étudies ou tu joues ? »
\item \zh{互联网有好处还是有坏处？} (Hùliánwǎng yǒu hǎochu háishi yǒu huàichu ?) « Internet a-t-il des avantages ou des inconvénients ? »
\item \zh{你早上上网还是晚上上网？} (Nǐ zǎoshang shàngwǎng háishi wǎnshang shàngwǎng ?) « Tu te connectes le matin ou le soir ? »
\item \zh{他用电脑发邮件还是用手机发邮件？} (Tā yòng diànnǎo fā yóujiàn háishi yòng shǒujī fā yóujiàn ?) « Il envoie ses courriels avec l'ordinateur ou avec le téléphone ? »
\end{enumerate}
\textbf{Point de vigilance :} \zh{还是} s'emploie uniquement dans les \emph{questions}. Dans une phrase affirmative, « ou » se dit \zh{或者} (huòzhě).
\end{corrige}

\begin{corrige}[Exercice 7 — Dialogue à trous]
Dialogue complété :

\begin{textebox}[Dialogue corrigé : 互联网好不好？]
甲：\textbf{我觉得}互联网很有用，我们可以学习很多东西。\\
(Jiǎ : Wǒ juéde hùliánwǎng hěn yǒuyòng, wǒmen kěyǐ xuéxí hěn duō dōngxi.)\\
乙：\textbf{我同意}！\textbf{因为}网上有很多信息，\textbf{所以}学习很方便。\\
(Yǐ : Wǒ tóngyì ! Yīnwei wǎngshang yǒu hěn duō xìnxī, suǒyǐ xuéxí hěn fāngbiàn.)\\
甲：可是上网时间太长也不好。\textbf{你呢}？\\
(Jiǎ : Kěshì shàngwǎng shíjiān tài cháng yě bù hǎo. Nǐ ne ?)\\
乙：对，我们要少玩游戏，多学习。\\
(Yǐ : Duì, wǒmen yào shǎo wán yóuxì, duō xuéxí.)
\end{textebox}

\textbf{Traduction :} « A : Je pense qu'internet est très utile, on peut apprendre beaucoup de choses. — B : Je suis d'accord ! Parce qu'il y a beaucoup d'informations sur internet, étudier est très pratique. — A : Mais passer trop de temps en ligne n'est pas bon non plus. Et toi ? — B : Oui, nous devons moins jouer et plus étudier. »

\textbf{Justification :} le premier trou exige une formule d'opinion (\zh{我觉得}) ; le second une formule d'accord (\zh{我同意}) ; \zh{因为……所以……} relie cause et conséquence ; \zh{你呢} relance le partenaire.
\end{corrige}

\begin{corrige}[Exercice 8 — Remets les mots dans l'ordre]
\begin{enumerate}
\item \zh{我每天都上网。} (Wǒ měitiān dōu shàngwǎng.) « Je me connecte tous les jours. » — Ordre : sujet + temps + \zh{都} + verbe.
\item \zh{我觉得互联网很有用。} (Wǒ juéde hùliánwǎng hěn yǒuyòng.) « Je pense qu'internet est très utile. » — \zh{很} précède l'adjectif \zh{有用}.
\item \zh{因为很方便，所以大家都喜欢它。} (Yīnwei hěn fāngbiàn, suǒyǐ dàjiā dōu xǐhuan tā.) « Parce que c'est très pratique, tout le monde l'aime. » — \zh{大家} « tout le monde » est sujet de la seconde proposition ; \zh{都} se place après.
\item \zh{我用手机和朋友们聊天。} (Wǒ yòng shǒujī hé péngyoumen liáotiān.) « Je bavarde avec mes amis avec mon téléphone. » — Ordre : sujet + \zh{用} + outil + \zh{和} + personne + verbe.
\end{enumerate}
\textbf{Point de vigilance :} en chinois, les compléments de temps, de moyen (\zh{用}...) et de compagnie (\zh{和}...) se placent \emph{avant} le verbe, contrairement au français.
\end{corrige}

\begin{corrige}[Exercice 9 — Compréhension de l'écrit et questions de langue]
\textbf{Questions de compréhension :}
\begin{enumerate}
\item Les élèves camerounais utilisent internet pour étudier (\zh{学习}), bavarder (\zh{聊天}) et lire les nouvelles (\zh{看新闻}).
\item L'auteur trouve internet utile parce qu'il y a beaucoup d'informations en ligne, ce qui rend l'étude très pratique (\zh{网上有很多信息，所以学习很方便}).
\item Il signale que certains élèves passent trop de temps à jouer aux jeux en ligne (\zh{玩游戏的时间太长}), ce qui n'est pas bon.
\item Il conseille de moins jouer et de plus étudier, et de bien utiliser internet (\zh{少玩游戏，多学习；好好用它}).
\end{enumerate}
\textbf{Questions de langue :}
\begin{enumerate}
\setcounter{enumi}{4}
\item \zh{因为网上有很多信息，所以学习很方便。} (Yīnwei wǎngshang yǒu hěn duō xìnxī, suǒyǐ xuéxí hěn fāngbiàn.) « Parce qu'il y a beaucoup d'informations sur internet, étudier est très pratique. »
\item \zh{总之} (zǒngzhī) = « en résumé, bref » ; c'est une formule de conclusion utile à l'oral.
\item \zh{他们用手机还是用电脑？} (Tāmen yòng shǒujī háishi yòng diànnǎo ?) « Utilisent-ils le téléphone ou l'ordinateur ? »
\end{enumerate}
\textbf{Barème indicatif (sur 10) :} compréhension 4 × 1,5 pt = 6 pts ; questions de langue 5 et 6 : 1 pt chacune ; question 7 : 2 pts (structure correcte 1 pt, tons et caractères 1 pt).
\end{corrige}

\begin{corrige}[Exercice 10 — Thème et version]
\textbf{Thème :}
\begin{enumerate}
\item \zh{每天我用手机学习，也和朋友聊天。} (Měitiān wǒ yòng shǒujī xuéxí, yě hé péngyou liáotiān.)
\item \zh{我觉得互联网有好处，也有坏处。} (Wǒ juéde hùliánwǎng yǒu hǎochu, yě yǒu huàichu.)
\item \zh{因为互联网很快，所以我们可以找到很多信息。} (Yīnwei hùliánwǎng hěn kuài, suǒyǐ wǒmen kěyǐ zhǎodào hěn duō xìnxī.)
\end{enumerate}
\textbf{Version :}
\begin{enumerate}
\item « Internet est très pratique, nous pouvons étudier à la maison. »
\item « Je ne suis pas d'accord avec ton point de vue, je pense que le temps passé en ligne ne doit pas être trop long. »
\item « Tu préfères utiliser le téléphone portable ou l'ordinateur ? »
\end{enumerate}
\textbf{Points de vigilance :} \zh{找到} (zhǎodào) exprime le résultat « trouver » ; \zh{看法} (kànfǎ) = « point de vue » ; ne pas oublier \zh{也} pour rendre « et... aussi » dans la phrase 2 du thème.
\textbf{Barème indicatif (sur 10) :} thème 3 × 2 pts = 6 pts ; version 3 × 1 pt + 1 pt d'ensemble = 4 pts.
\end{corrige}

\begin{corrige}[Exercice 11 — Expression écrite et préparation à l'oral]
\textbf{Proposition de rédaction modèle (environ 70 caractères) :}

\begin{textebox}[互联网和我们的生活 — texte modèle]
互联网和我们的生活关系很大。第一，互联网很有用，我们可以在网上学习、看新闻。第二，因为网上有很多信息，所以找东西很方便。可是，玩游戏时间太长也不好。总之，互联网是好东西，我们要好好用它。\\
(Hùliánwǎng hé wǒmen de shēnghuó guānxì hěn dà. Dì-yī, hùliánwǎng hěn yǒuyòng, wǒmen kěyǐ zài wǎngshang xuéxí, kàn xīnwén. Dì-èr, yīnwei wǎngshang yǒu hěn duō xìnxī, suǒyǐ zhǎo dōngxi hěn fāngbiàn. Kěshì, wán yóuxì shíjiān tài cháng yě bù hǎo. Zǒngzhī, hùliánwǎng shì hǎo dōngxi, wǒmen yào hǎohǎo yòng tā.)
\end{textebox}

\textbf{Traduction :} « Internet a un grand rapport avec notre vie. Premièrement, internet est très utile : nous pouvons étudier et lire les nouvelles en ligne. Deuxièmement, parce qu'il y a beaucoup d'informations sur internet, chercher des choses est très pratique. Mais jouer trop longtemps n'est pas bon non plus. En résumé, internet est une bonne chose, nous devons bien l'utiliser. »

\textbf{Modèle de plan d'exposé oral (5 phrases) :}
\begin{enumerate}
\item Présentation : \zh{我给大家介绍一下互联网。} (Wǒ gěi dàjiā jièshào yīxià hùliánwǎng.)
\item Avis : \zh{我觉得互联网很有用。}
\item Avantage : \zh{我们可以在网上学习汉语。}
\item Inconvénient : \zh{可是玩游戏太多不好。}
\item Conclusion : \zh{总之，我们要好好用互联网。}
\end{enumerate}

\textbf{Exemple de jeu de rôle (4 échanges minimum) :} A : \zh{你每天上网吗？} — B : \zh{对，我每天都上网。你呢？} — A : \zh{我也上网。你用手机还是电脑？} — B : \zh{我用手机。为什么你问？} — A : \zh{因为我想知道大家的习惯。真的吗？你上网学习吗？} — B : \zh{当然！互联网对学习很有用。}

\textbf{Points de vigilance :}
\begin{itemize}
\item Respecter la structure \zh{第一……第二……总之……} sans oublier la virgule chinoise \zh{，}.
\item Dans \zh{因为……所以……}, garder les deux connecteurs (c'est correct en chinois).
\item À l'oral, soigner les tons : \zh{wǎng} (3e ton), \zh{hěn} (3e ton), \zh{yǒuyòng} (3e + 4e ton).
\item Compter les caractères : 60 à 80, ni plus ni moins.
\end{itemize}
\textbf{Barème indicatif (sur 20) :} écrit — respect du sujet et de la longueur 4 pts, structures imposées 4 pts, correction grammaticale 6 pts, caractères 2 pts ; oral — prononciation des tons 2 pts, fluidité et relances 2 pts.
\end{corrige}

\newpage

\coursec{Fiche bilan}

\begin{popfigure}
\begin{tikzpicture}[mindmap, grow cyclic, every node/.style={concept, text=white, font=\small\bfseries, align=center}, root concept/.append style={concept color=popblue, font=\large\bfseries}, level 1/.append style={level distance=3.4cm, sibling angle=51}, level 2/.append style={level distance=2.2cm, font=\scriptsize}]
\node[root concept, align=center]{Expression orale :\\ 互联网 Internet}
  child[concept color=poporange]{ node[align=center]{Lexique\\ 上网、网站、手机} }
  child[concept color=popgreen]{ node[align=center]{Dialogue 1\\ 谈上网习惯} }
  child[concept color=poppurple]{ node[align=center]{Formules\\ 我认为…} }
  child[concept color=popgold]{ node[align=center]{Dialogue 2\\ 优点和缺点} }
  child[concept color=poppink]{ node[align=center]{Phonétique\\ 四声练习} }
  child[concept color=popblue!70]{ node[align=center]{Jeux de rôle\\ 班级辩论} }
  child[concept color=popgreen!70!black]{ node[align=center]{Méthode\\ 介绍—观点—结论} };
\end{tikzpicture}
\legende{Carte mentale de la leçon : parler d'internet en chinois}
\end{popfigure}

\begin{bilan}
\textbf{1. Lexique clé à connaître par cœur}
\begin{center}
\begin{tabularx}{\linewidth}{|X|X|X|X|}\hline
\rowcolor{popblueL} Caractères & Pinyin & Nature & Traduction \\ \hline
互联网 & hùliánwǎng & n. & internet \\ \hline
上网 & shàngwǎng & v. & aller sur internet \\ \hline
网站 & wǎngzhàn & n. & site web \\ \hline
手机 & shǒujī & n. & téléphone portable \\ \hline
聊天 & liáotiān & v. & bavarder, chatter \\ \hline
发电子邮件 & fā diànzǐ yóujiàn & v. & envoyer un courriel \\ \hline
查资料 & chá zīliào & v. & chercher des informations \\ \hline
优点 & yōudiǎn & n. & avantage \\ \hline
缺点 & quēdiǎn & n. & inconvénient \\ \hline
有用 & yǒuyòng & adj. & utile \\ \hline
浪费时间 & làngfèi shíjiān & v. & perdre du temps \\ \hline
\end{tabularx}
\end{center}

\textbf{2. Formules fonctionnelles du débat}
\begin{center}
\begin{tabularx}{\linewidth}{|X|X|X|}\hline
\rowcolor{popblueL} Fonction & Formule + pinyin & Traduction \\ \hline
Présenter & 我今天要谈一谈互联网。Wǒ jīntiān yào tányitan hùliánwǎng. & Aujourd'hui je vais parler d'internet. \\ \hline
Donner son avis & 我认为互联网很有用。Wǒ rènwéi hùliánwǎng hěn yǒuyòng. & Je pense qu'internet est très utile. \\ \hline
Être d'accord & 我同意你的看法。Wǒ tóngyì nǐ de kànfǎ. & Je suis d'accord avec toi. \\ \hline
Ne pas être d'accord & 我不太同意。Wǒ bú tài tóngyì. & Je ne suis pas tout à fait d'accord. \\ \hline
Relancer & 你觉得呢？Nǐ juéde ne？ & Et toi, qu'en penses-tu ? \\ \hline
Conclure & 所以，我们应该好好用互联网。Suǒyǐ, wǒmen yīnggāi hǎohāo yòng hùliánwǎng. & Donc, nous devons bien utiliser internet. \\ \hline
\end{tabularx}
\end{center}

\textbf{3. Grammaire à revoir}
\begin{itemize}
\item \zh{认为} + proposition : \zh{我认为上网很方便。} (Wǒ rènwéi shàngwǎng hěn fāngbiàn.) « Je pense qu'aller sur internet est pratique. »
\item \zh{虽然……但是……} : \zh{虽然互联网有用，但是也会浪费时间。} « Bien qu'internet soit utile, il fait aussi perdre du temps. »
\item \zh{应该} + verbe pour le conseil : \zh{我们应该少玩手机。} « Nous devrions moins utiliser le portable. »
\item Durée avec \zh{了} : \zh{我上了两个小时的网。} « J'ai surfé deux heures. »
\end{itemize}

\textbf{4. Méthode express pour l'exposé oral}
\begin{enumerate}
\item \cle{Introduction} : annoncer le sujet (我今天要谈一谈……).
\item \cle{Développement} : 2 ou 3 arguments avec exemples (优点 / 缺点).
\item \cle{Conclusion} : donner son avis personnel (我认为……所以……).
\item Parler lentement, soigner les \cle{tons}, regarder l'auditoire.
\end{enumerate}
\end{bilan}

\begin{aretenir}[Checklist avant le Bac]
\begin{itemize}
\item[$\square$] Je connais le lexique d'internet : 互联网、上网、网站、手机、聊天.
\item[$\square$] Je sais présenter un sujet : 我今天要谈一谈……
\item[$\square$] Je sais donner mon avis avec 我认为 et 我觉得.
\item[$\square$] Je sais être d'accord (我同意) ou pas (我不太同意) poliment.
\item[$\square$] Je sais relancer le débat : 你觉得呢？
\item[$\square$] Je maîtrise la structure 虽然……但是…… pour nuancer.
\item[$\square$] Je sais exprimer avantages (优点) et inconvénients (缺点).
\item[$\square$] Je prononce correctement les quatre tons, surtout dans 上网 (shàngwǎng) et 手机 (shǒujī).
\item[$\square$] Je sais conclure un exposé avec 所以 et 应该.
\item[$\square$] Je peux parler une à deux minutes sans lire mes notes.
\end{itemize}
\end{aretenir}

\findecours

\end{document}
